\documentclass[11pt]{amsart}
\pdfinfoomitdate=1
\pdftrailerid{}
\pdfsuppressptexinfo=15
\usepackage[T1]{fontenc}
\usepackage{lmodern,amsmath,amssymb,amsthm,mathtools}
\usepackage[margin=1.08in]{geometry}
\usepackage[expansion=false]{microtype}
\newcommand{\DoNotLoadEpstopdf}{}
\usepackage{enumitem,booktabs,longtable,array,needspace,tikz-cd}
\usepackage[hidelinks,pdfusetitle]{hyperref}
\usepackage[capitalise,noabbrev,nameinlink]{cleveref}
\hypersetup{pdftitle={Orbifold and logarithmic Iitaka subadditivity},pdfauthor={OpenAI},bookmarksdepth=2}
\newcommand{\C}{\mathbf C}
\newcommand{\Q}{\mathbf Q}
\newcommand{\R}{\mathbf R}
\newcommand{\Z}{\mathbf Z}
\newcommand{\N}{\mathbf N}
\newcommand{\cO}{\mathcal O}
\newcommand{\cV}{\mathcal V}
\newcommand{\cA}{\mathcal A}
\newcommand{\ii}{\sqrt{-1}}
\newcommand{\ddc}{dd^c}
\newcommand{\kbar}{\bar\kappa}
\newcommand{\red}{\mathrm{red}}
\newcommand{\hor}{\mathrm{hor}}
\DeclareMathOperator{\Gr}{Gr}
\DeclareMathOperator{\End}{End}
\DeclareMathOperator{\Hom}{Hom}
\DeclareMathOperator{\Sym}{Sym}
\DeclareMathOperator{\im}{im}
\DeclareMathOperator{\rank}{rank}
\DeclareMathOperator{\rk}{rk}
\DeclareMathOperator{\ord}{ord}
\DeclareMathOperator{\Supp}{Supp}
\DeclareMathOperator{\Var}{Var}
\DeclareMathOperator{\vol}{vol}
\DeclareMathOperator{\trdeg}{trdeg}
\DeclareMathOperator{\Frac}{Frac}
\DeclareMathOperator{\Spec}{Spec}
\DeclareMathOperator{\Gal}{Gal}
\DeclareMathOperator{\Bir}{Bir}
\DeclareMathOperator{\Aut}{Aut}
\DeclareMathOperator{\Pic}{Pic}
\DeclareMathOperator{\lct}{lct}
\DeclareMathOperator{\mult}{mult}
\DeclareMathOperator{\id}{id}
\newtheorem{theorem}{Theorem}[section]
\newtheorem{lemma}[theorem]{Lemma}
\newtheorem{proposition}[theorem]{Proposition}
\newtheorem{corollary}[theorem]{Corollary}
\theoremstyle{definition}
\newtheorem{definition}[theorem]{Definition}
\newtheorem{construction}[theorem]{Construction}
\newtheorem{convention}[theorem]{Convention}
\theoremstyle{remark}
\newtheorem{remark}[theorem]{Remark}
\newtheorem{example}[theorem]{Example}
\numberwithin{equation}{section}
\setcounter{tocdepth}{1}
\allowdisplaybreaks[2]
\emergencystretch=1em
\title[Orbifold and logarithmic Iitaka subadditivity]{Orbifold and logarithmic Iitaka subadditivity}
\author{OpenAI}
\date{September 26, 2026}
\begin{document}
\begin{abstract}
We prove Campana's orbifold Iitaka subadditivity conjecture for rational
simple normal crossing boundaries on compact manifolds in Fujiki class
$\mathcal C$, including coefficient one. Ordinary and logarithmic
subadditivity follow.
\end{abstract}
\maketitle
\tableofcontents
\section{Introduction}\label{intro:section}

The Kodaira dimension measures the growth of pluricanonical forms.
For a surjective morphism $f:X\to Y$ of smooth projective complex varieties
with connected fibers, Iitaka's subadditivity problem asks whether
\[
                  \kappa(X)\geq\kappa(F)+\kappa(Y),
\]
where $F$ is a general fiber. The logarithmic version permits poles
along reduced boundary divisors and therefore applies to open
varieties. We prove a common generalization in which rational boundary
coefficients and multiple fibers both contribute to the inequality.

Iitaka's divisor-dimension and logarithmic theories provide the
language for this question \cite{IitakaOriginal,IitakaLog}; an
explicit logarithmic subadditivity conjecture appears in
\cite[\S3]{IitakaClassification}.
The lower bound requires global pluricanonical forms that retain
the growth of both the fibers and the base. Direct-image positivity
established central cases: Kawamata treated curve bases and fibers
admitting good minimal models, Viehweg developed weak positivity
and addition over a base of general type, and Koll\'ar proved
addition for general-type fibers
\cite{KawamataCurves,KawamataKodaira,ViehwegAdditivity,Kollar87}.
Weak positivity controls symmetric powers after small positive
ample twists; a general-type base provides enough canonical
positivity to absorb them.

Other advances exploit dimension or the geometry of the base:
Birkar treated total dimension at most six, Cao--P\u{a}un and
Hacon--Popa--Schnell treated abelian or maximal-Albanese-dimension
bases, and Cao treated projective klt pairs over surfaces
\cite{BirkarIitakaSix,CaoPaunAbelian,HaconPopaSchnellAbelian,CaoSurfaces}.
Thus the established cases draw positivity from different parts of
the fibration, rather than imposing one uniform condition on it.

For logarithmic pairs, Kov\'acs--Patakfalvi proved the projective
log-general-type-fiber case \cite[Theorem~9.6]{KovacsPatakfalvi},
while Hashizume proved reduced-SNC subadditivity for fibers with
abundant log canonical divisor \cite[Theorem~1.2]{Hashizume}.
Maehara had established logarithmic addition over a base of log
general type; Fujino gives a proof through twisted weak positivity
\cite[Theorem~1.9]{FujinoWeakPositivity}.
Campana proved addition over an orbifold base of general type, and
Wang proved klt K\"ahler subadditivity over complex tori, including
an inf-multiplicity refinement
\cite[Theorem~6.3]{CampanaOrbifolds}
\cite[Main Theorem~A and Theorem~E]{WangKahlerIitaka}.
The present argument uses general-type addition only after a relative
Iitaka reduction; the original fiber need not be of general type
or admit a good minimal model.

Earlier preprints also state general ordinary inequalities.
Tsuji states ordinary subadditivity following a direct-image generation
theorem \cite[Theorems~1.9 and~1.13]{TsujiDirectImages}; his
Remark~1.14 distinguishes an unproved variation refinement.
Maehara states an ordinary determinant inequality involving variation
\cite[\S7, Theorem~8]{MaeharaIitaka}. This later preprint claim is
distinct from his logarithmic general-type-base theorem above.
Neither preprint statement is a proof input here.

\subsection{The orbifold base and the main theorem}

Multiple fibers carry information that the ordinary base omits.
Campana's orbifold base records it through inf multiplicities, and
his birational invariant gives the corresponding form of Iitaka's
conjecture \cite{Campana2004,CampanaOrbifolds}. We specify that
invariant before stating the theorem.

\label{setup:orbifold-base}
A fibration is a proper surjective holomorphic map with connected
fibers. A rational SNC boundary on a smooth manifold is a finite
sum $\Delta=\sum_i\delta_i\Delta_i$, with
$\delta_i\in\mathbf Q\cap[0,1]$ and simple normal crossing support.
Write $\Delta_E$ for the coefficient of a source prime $E$. Set
\[
 m_\Delta(E)=\frac1{1-\Delta_E},
 \qquad \frac1{0}=\infty,
\]
where $\Delta_E=0$ off the boundary. If the base of
$f:(X,\Delta)\to Y$ is smooth, define
\begin{equation}\label{setup:inf-boundary}
\begin{gathered}
 m(f,\Delta;D)=\min_{E\mapsto D}
     \bigl\{\operatorname{ord}_E(f^*D)m_\Delta(E)\bigr\},
 \qquad
 B(f,\Delta)=\sum_D\left(1-\frac1{m(f,\Delta;D)}\right)D.
\end{gathered}
\end{equation}
Here $D$ ranges over base prime divisors, $E$ over source primes
dominating $D$, and $1/\infty=0$. The minimum is nonempty and
only finitely many coefficients of $B(f,\Delta)$ are nonzero.
There is no divisibility condition on these multiplicities.

For smooth sources and bases, an elementary change of model is a
commutative diagram of fibrations
\[
\begin{tikzcd}
 (X',\Delta') \arrow[r,"u"] \arrow[d,"f'"'] &
 (X,\Delta) \arrow[d,"f"]\\
 Y' \arrow[r,"v"'] & Y
\end{tikzcd}
\]
where $u,v$ are proper bimeromorphic holomorphic maps,
$u_*\Delta'=\Delta$, and both boundaries are rational SNC
boundaries in $[0,1]$. We require, for every prime $D$ on $X$
and prime $E$ on $X'$ with $t=\operatorname{ord}_E(u^*D)>0$,
\begin{equation}\label{setup:orbifold-morphism}
                     t\,m_{\Delta'}(E)\geq m_\Delta(D).
\end{equation}
We permit $\infty\geq\infty$ and impose no condition for $t=0$.
Write $f'\sim f$ for the equivalence relation generated by these
diagrams, allowing arrows in either direction. Define
\begin{equation}\label{setup:invariant-base}
 \kappa(f\mid\Delta)=
 \inf_{f':(X',\Delta')\to Y'\sim f}
       \kappa\bigl(Y',K_{Y'}+B(f',\Delta')\bigr).
\end{equation}
For a normal singular base $Y$, resolve $Y$, resolve the main
component of the fiber product with $X$, and give the resulting
smooth source the strict transform of $\Delta$ plus the reduced
exceptional divisor over $X$. After resolving this boundary, apply
the preceding definition. It is independent of the chosen resolution
\cite[Definitions~2.1, 2.3, 3.2, 4.7, and~4.10;
Remark~4.13 and Corollary~4.14]{CampanaOrbifolds}.

For a rational line bundle $L$, $\kappa(Z,L)$ is the maximal image
dimension of its complete linear systems in sufficiently divisible
positive degrees, and is $-\infty$ if all these systems are empty.
A point has Kodaira dimension zero, and
$(-\infty)+a=-\infty$, including $a=-\infty$.
Fujiki class $\mathcal C$ consists of compact complex spaces
bimeromorphic to compact K\"ahler spaces.

\begin{theorem}[Orbifold subadditivity]\label{main:orbifold}
Let $X$ be a smooth compact connected complex manifold in Fujiki
class $\mathcal C$, let $\Delta$ be a rational SNC boundary, and
let $f:X\to Y$ be a surjective holomorphic map with connected
fibers onto a normal compact irreducible complex space.
For a very general smooth fiber $F$, put $\Delta_F=\Delta|_F$.
Then
\begin{equation}\label{main:inequality}
 \kappa(X,K_X+\Delta)
 \geq\kappa(F,K_F+\Delta_F)+\kappa(f\mid\Delta).
\end{equation}
\end{theorem}

This is a positive resolution of Campana's orbifold Iitaka
conjecture in its rational SNC, Fujiki-class-$\mathcal C$
formulation, including coefficient one
\cite[Conjecture~6.1 and Remark~6.2]{CampanaOrbifolds}.
Here a very general fiber is chosen outside a countable union of
proper closed analytic subsets, including the critical values and
the exceptional images of boundary strata.

For reduced SNC divisors $D_X,D_Y$ satisfying
$\operatorname{Supp}(f^*D_Y)\subseteq\operatorname{Supp}(D_X)$,
the invariant base term is at least $\kappa(Y,K_Y+D_Y)$.
Theorem~\ref{main:orbifold} therefore gives
\[
 \kappa(X,K_X+D_X)
 \geq\kappa(F,K_F+D_X|_F)+\kappa(Y,K_Y+D_Y)
\]
for fibrations between smooth compact class-$\mathcal C$ manifolds.
For a smooth quasi-projective variety $U$, its logarithmic Kodaira
dimension is $\bar\kappa(U)=\kappa(X,K_X+D_X)$ on a smooth
projective compactification with reduced SNC boundary
$D_X=X\setminus U$. Compatible compactifications give logarithmic
subadditivity for dominant morphisms of smooth quasi-projective
complex varieties with geometrically connected general fiber.
Empty boundaries give ordinary subadditivity. The projective
ordinary and logarithmic statements also hold over algebraically
closed fields of characteristic zero, using geometric generic fibers.

\subsection{The reduction and the two positivity inputs}

The proof first replaces the original fiber by a general-type
object without losing its complete pluricanonical systems.
Assume the two terms on the right of Theorem~\ref{main:orbifold}
are nonnegative, and put $k=\kappa(F,K_F+\Delta_F)$.
A relative Iitaka map gives, after the model changes of Section~\ref{setup:section},
\[
 X\xrightarrow{g}W\xrightarrow{h}Y,
 \qquad \dim W_y=k,
 \qquad\kappa(G,K_G+\Delta_G)=0
\]
for a very general $g$-fiber $G$.

The canonical-bundle-formula strategy separates the contribution
of singular fibers from a moduli contribution on $W$.
Fujino--Mori developed a higher-dimensional logarithmic formula,
and Ambro organized its discriminant and moduli parts birationally
\cite{FujinoMori,AmbroBoundary,AmbroModuli}.
Our use of this strategy requires an exact comparison on the fixed
original source. Section~\ref{setup:section} obtains it from the orders of relative
pluriforms. A sufficiently divisible pluricanonical system on $G$
is one-dimensional. Taking a root of its generator on a cyclic
cover gives a differential form, and the line it spans belongs to
the cohomology of that cover. Tensor descent and extension across
the boundary produce a rational line bundle $M$ on $W$.
The orders of the same form determine a boundary $T$ such that
\[
 B(g,\Delta)\leq T\leq1,
 \qquad
 H^0\bigl(X,m(K_X+\Delta)\bigr)
 \simeq H^0\bigl(W,m(K_W+T+M)\bigr)
\]
in sufficiently divisible degrees. The comparison also holds
relatively over $Y$. Thus $K_W+T+M$ is big on the
$k$-dimensional fibers of $h$.

Two inputs turn this fiberwise bigness into the required global
lower bound. The first concerns the line $M$: the period map,
which records how the Hodge structures of the covers vary, supplies
a second fibration $p:W\to S$, with $S$ smooth projective, and
a nef rational line bundle $P$ such that
\[
                      M=p^*P,\qquad K_S+aP\ \text{is big}
\]
for some positive rational $a$. The maps $h$ and $p$ have different
jobs, as Figure~\ref{fig:iitaka-period} shows. The second input is
log-general-type addition: a log canonical divisor that is big on
the fibers of $h$ has Kodaira dimension at least their dimension
plus the orbifold contribution of $Y$.

\begin{figure}[ht]
\centering
\begin{tikzcd}[column sep=large,row sep=large]
 X \arrow[r,"g"] & W \arrow[r,"h"] \arrow[d,"p"'] & Y\\
 & S &
\end{tikzcd}
\caption{The relative Iitaka reduction factors the original fibration
through $W$. The map $h$ retains the original base and has
$k$-dimensional general fibers. The independent map $p$ expresses
$M$ as the pullback of a nef line $P$ on a projective base where
$K_S+aP$ is big. No map between $Y$ and $S$ is asserted.}
\label{fig:iitaka-period}
\end{figure}


Section~\ref{addition:section} states both inputs precisely and proves the intervening
addition principle. Write $D_W=K_W+T$. For a fixed rational $c<1$
close to one, general-type addition gives sufficiently many ratios
of sections of $D_W+cM+\eta p^*A$, where $A$ is ample on $S$
and $\eta>0$ is rational. Weak positivity then supplies one fixed
section of $D_W+aM-\lambda p^*A$ for suitable $a>1$ and
$\lambda>0$. A positive rational combination cancels the ample
terms and gives coefficient one on $M$. Multiplication by powers
of that fixed section preserves the ratios of the first system.
This is an exact section argument; it requires neither abundance
of $P$ nor a limiting assertion for Kodaira dimensions.

The full proofs of the inputs follow this calculation.
Section~\ref{period:section} proves adjoint positivity of the Hodge line. It combines
the curvature and boundary theory of Griffiths, Schmid, and
Cattani--Kaplan--Schmid with a projective period quotient
\cite{GriffithsPeriodsIII,Schmid,CKS}.
The construction retains suitable rational adjoint factors of the
ambient integral variation, associated with its monodromy, from which
a positive power of the chosen line can be recovered. Their full
periods give the quotient, using Villadsen's class-$\mathcal C$
compactification theorem and the period-quotient approach of
Bakker--Brunebarbe--Tsimerman and Sommese; curvature and a volume
argument establish its projectivity
\cite{BakkerBrunebarbeTsimerman,SommesePeriodMapping,VilladsenHodgeKahler}.

On this quotient $S$, the retained full period map is generically
immersive, but the map remembering only the chosen line may have
smaller rank. Curvature identifies the latter rank with the numerical
dimension $\nu(P)$, measured by its nonzero intersection powers.
Let $D$ be the reduced boundary where the retained variation has
nonidentity local monodromy. Brunebarbe--Cadorel gives $K_S+D$ big
\cite{BrunebarbeCadorel}. Using Andr\'e's monodromy normality and
Bakker--Tsimerman's Ax--Schanuel theorem to analyze fibers of the
line map, the boundary argument proves
$\nu(P|_{D_i})<\nu(P)$ for every component $D_i$ of $D$, including
finite nonidentity monodromy
\cite{AndreMonodromy,BakkerTsimermanAxSchanuel}.
The boundary loss in the asymptotic section count consequently has
lower degree in the coefficient of $P$ than the available sections
on $S$. Lemma~\ref{period:subtract} turns this comparison into
bigness of $K_S+aP$.

Section~\ref{generaltype:section} proves log-general-type addition in class $\mathcal C$.
Kov\'acs--Patakfalvi's stable-family positivity applies to projective
families. We construct such an auxiliary family and descend its
relative sections to the original class-$\mathcal C$ base.
Fujino's weak positivity supplies the other direct-image input
used in Section~\ref{addition:section} \cite{KovacsPatakfalvi,FujinoWeakPositivity}.
Section~\ref{application:section} then returns once to the original fibration, applies the
addition principle to $(W,T,M)$, and derives the ordinary,
logarithmic, characteristic-zero, and core consequences. The core
and the relevant specialness of Kodaira-zero fibers are Campana's
constructions and results \cite{Campana2004,CampanaOrbifolds}.

\subsection{The reduced-boundary refinement}

For a rational boundary, the root form above spans one character
subspace of the cyclic cover's highest Hodge space; the whole
highest space can be larger. For projective reduced-boundary pairs
of logarithmic Kodaira dimension zero, Section~\ref{cyc:section} proves more:
the minimal root cover has a one-dimensional entire logarithmic
top-form space. Mixed-Hodge extension and semipositivity in the
forms of Fujino--Fujisawa control the coefficient-one boundary
\cite{FujinoFujisawa,FujinoFujisawaVMHS}.
We identify the resulting line with the moduli divisor in the
projective subpair setting of Bakker--Filipazzi--Mauri--Tsimerman
\cite[Theorem~6.28]{BFMT}. This identification does not use their
separate semiampleness theorem, whose hypotheses include
effectiveness over the generic point \cite[Theorem~1.5]{BFMT}.
Section~\ref{cyc:section} also gives the exact divisorial normalization.
Appendix~\ref{p1logcbfalt:section} collects complementary projective
calculations by valuations and toroidal models, together with the
comparison on higher models.

The logarithmic lower bound, adjoint positivity, and this
reduced-boundary refinement supply the stated inputs to the
whole-fiber-variation companion \cite{IitakaVariationCompanion}. The reverse-inequality companion
\cite{IitakaAdditivityCompanion} uses the lower bound only for additivity; its upper theorem has
additional smoothness hypotheses on all boundary strata.
Neither companion is used in the proof of Theorem~\ref{main:orbifold}.
The separate article on projective Hodge lines
\cite{IitakaProjectiveCompanion} gives a complete
ordinary proof through rationally polarized integral variations,
together with the distinct integral and analytic boundary arguments.

\section{The relative Iitaka reduction and its Hodge line}
\label{setup:section}

The first task is to replace the fibers by their Iitaka images while
retaining every sufficiently divisible pluricanonical system. We first
fix the birational models on which divisorial order tests suffice. The
root form of a Kodaira-zero fiber will then provide a Hodge line and an
exact comparison of sections on the intermediate base.

We use additive notation for rational line bundles; equality of them
means equality in $\operatorname{Pic}\otimes\Q$. A rational line is
effective if a positive integral multiple has a nonzero section.
All degree assertions below are made after clearing denominators.
For a fibration $g:(Z,\Gamma)\to T$ and a very general smooth
fiber $F$, adjunction gives
\begin{equation}\label{setup:fiber-adjunction}
 (K_Z+\Gamma)|_F=K_F+\Gamma|_F.
\end{equation}

\subsection{The fixed smooth models}

A smooth-base fibration $g:(Z,\Gamma)\to T$ is \emph{neat} if there is a
proper bimeromorphic orbifold morphism
$w:(Z,\Gamma)\to(Z_0,\Gamma_0)$ to a smooth pair, with
$w_*\Gamma=\Gamma_0$, such that every prime divisor sent by $g$ into
codimension at least two is also sent by $w$ into codimension at least
two. Campana's construction gives suitable neat models, model
independence, and the formula
\begin{equation}\label{setup:neat-formula}
 \kappa(g\mid\Gamma)=\kappa\bigl(T,K_T+B(g,\Gamma)\bigr)
 \quad\text{on a neat model}
\end{equation}
\cite[Definition~3.8, Proposition~3.10, Corollary~4.11, and
Remark~6.2]{CampanaOrbifolds}.
For the proof we only need the defining inequality, valid on every
allowed smooth model:
\begin{equation}\label{setup:base-comparison}
 \kappa\bigl(T,K_T+B(g,\Gamma)\bigr)
 \ \geq\ \kappa(g\mid\Gamma).
\end{equation}

\subsection{Reduced exceptional boundaries}

\begin{lemma}[Logarithmic modifications]\label{setup:log-modification}
Let $p:Z'\to Z$ be a proper bimeromorphic holomorphic map between smooth
compact manifolds. Suppose that $\Gamma$ has simple normal crossing
support and coefficients in $\Q\cap[0,1]$, and set
\[
 \Gamma'=p_*^{-1}\Gamma+\operatorname{Exc}(p)_{\mathrm{red}}.
\]
Assume, after further resolution if necessary, that this boundary has
simple normal crossing support. There is an effective $p$-exceptional
rational divisor $A$ such that
\begin{equation}\label{setup:log-discrepancy}
 K_{Z'}+\Gamma'=p^*(K_Z+\Gamma)+A.
\end{equation}
For every sufficiently divisible positive integer $m$, the natural
identification of meromorphic pluricanonical forms induces
\begin{equation}\label{setup:section-invariance}
 p_*\mathcal O_{Z'}\bigl(m(K_{Z'}+\Gamma')\bigr)
 =\mathcal O_Z\bigl(m(K_Z+\Gamma)\bigr).
\end{equation}
These identifications respect multiplication, so the divisible section
rings and their Iitaka dimensions agree. Moreover, $p$ satisfies
\eqref{setup:orbifold-morphism}.
\end{lemma}

\begin{proof}
The difference in \eqref{setup:log-discrepancy} is exceptional, since
$p$ is an isomorphism at the generic point of every target prime.
To compute its coefficient at an exceptional prime $E$, choose
compatible local canonical forms.
Near a general point of its center on $Z$, let $z_1,\ldots,z_n$ be
coordinates in which the boundary components containing that center are
$z_i=0$ for $1\leq i\leq r$, with coefficients $\gamma_i$.
Put $t_i=\ord_E(p^*z_i)$ and let $k_E$ be the order of the Jacobian of
$p$ along $E$. The pullback of
\[
 \frac{dz_1}{z_1}\wedge\cdots\wedge\frac{dz_r}{z_r}
 \wedge dz_{r+1}\wedge\cdots\wedge dz_n
\]
has at most a simple pole along $E$. Indeed, at a general smooth point of
$E$, each logarithmic factor has the form $t_i\,du/u$ plus a holomorphic
one-form, and a nonzero wedge product contains at most one factor
$du/u$. Consequently $k_E+1\geq\sum_i t_i$. The coefficient in question
is therefore
\[
 k_E+1-\sum_i\gamma_i t_i
 \ \geq\ \sum_i(1-\gamma_i)t_i\ \geq\ 0.
\]
For an effective integral exceptional divisor $mA$, normality of $Z$
gives $p_*\mathcal O_{Z'}(mA)=\mathcal O_Z$: a meromorphic function
with poles only on $mA$ descends outside the codimension-two exceptional
image and extends by Hartogs' theorem. The projection formula proves
\eqref{setup:section-invariance}, compatibly with products and ratios.

Finally, an exceptional prime has infinite boundary multiplicity.
A nonexceptional prime is a strict transform, with pullback order one
and unchanged coefficient. These two cases prove
\eqref{setup:orbifold-morphism}.
\end{proof}

\begin{corollary}[Descent across exceptional centers]
\label{setup:exceptional-descent}
In the situation of \Cref{setup:log-modification}, a meromorphic
$m$-pluricanonical form on $Z'$ satisfying the boundary bounds of
$m\Gamma'$ at every nonexceptional prime descends to a section of
$m(K_Z+\Gamma)$ whenever $m\Gamma$ and $m\Gamma'$ are integral.
The assertion also holds after tensoring by the pullback of a line
bundle on $Z$.
\end{corollary}

\begin{proof}
On the complement of the exceptional image the form is a holomorphic
section of the prescribed line bundle on $Z$: divisorial regularity
suffices on a smooth space. A local trivialization reduces extension
across the remaining codimension-two analytic subset to Hartogs'
extension theorem. The tensor-twisted assertion has the same proof.
\end{proof}

\begin{corollary}[Fibers, composition, and a fixed base]
\label{setup:modification-compatibility}
Let $p:(Z',\Gamma')\to(Z,\Gamma)$ be a modification with the boundary
rule of \Cref{setup:log-modification}. Then:
\begin{enumerate}[label=\textup{(\roman*)}]
\item If $p$ is a modification over a fixed base, it induces a
logarithmic modification on a very general smooth fiber. The log
Kodaira dimensions of the corresponding fibers agree.
\item A finite composition of such modifications again has boundary
equal to the strict transform of the initial boundary plus the reduced
exceptional divisor of the composite.
\item If $g:(Z,\Gamma)\to T$ has smooth base and $g'=g\circ p$, then
\begin{equation}\label{setup:fixed-base-boundary}
 B(g',\Gamma')=B(g,\Gamma).
\end{equation}
\end{enumerate}
\end{corollary}

\begin{proof}
For (i), choose the parameter so that the relevant smoothness and
dimension statements hold for the finitely many exceptional components
and their images. A component dominating the base restricts to a
reduced divisor, and its exceptional image restricts to codimension at
least two. Components not dominating the base do not meet this fiber.
The restricted map has the same boundary rule, so
\Cref{setup:log-modification} applies on the fiber. For (ii), every
prime exceptional for the composite is either newly exceptional or the
transform of an earlier exceptional prime; in both cases its coefficient
is one.

For (iii), fix a prime $D\subset T$. Every old prime dominating $D$
has a strict transform, with the same pullback order and boundary
multiplicity. Every new prime in the minimum defining
$m(g',\Gamma';D)$ is exceptional for $p$, so its boundary multiplicity
is infinite. Adding these entries leaves the old minimum unchanged,
including when that minimum was already infinite.
\end{proof}

For choice independence in the singular-base definition, take a common
resolution of two chosen diagrams. Their induced source boundaries
coincide: old primes retain their coefficients and all primes exceptional
over the initial $X$ have coefficient one. The comparison maps satisfy
\Cref{setup:log-modification}, so the smooth fibrations are equivalent
and their total-space and very general fiber log Kodaira dimensions agree.

\subsection{Flattening relative to a fixed reference pair}

To compare forms over generic points of base divisors, we need control of
source divisors over higher codimension. Proper analytic flattening gives
a base modification whose strict transform is flat
\cite{HironakaFlattening,FujikiDouady}. Here the strict transform is the
closure of the original family over the isomorphism locus, with base
torsion removed. Further base change preserves flatness. We use only
the resulting equidimensionality before resolving this intermediate
source.

\begin{lemma}[Flattening and the exceptional reference]
\label{setup:flattening}
Let $g_0:(Z_0,\Gamma_0)\to T_0$ be a fibration from a smooth compact
pair with rational SNC boundary in $[0,1]$ to a smooth compact manifold.
There are a smooth-base modification
$q:T\to T_0$, a smooth-source modification $p:Z\to Z_0$, and a
fibration $g:Z\to T$ with $q\circ g=g_0\circ p$ such that, for
\[
 \Gamma=p_*^{-1}\Gamma_0+\operatorname{Exc}(p)_{\mathrm{red}},
\]
the boundary has simple normal crossing support and every prime divisor
$E\subset Z$ satisfying $\operatorname{codim}_Tg(E)\geq2$ is
$p$-exceptional. In particular every such prime has coefficient one in
$\Gamma$, and the resulting fibration is neat with reference pair
$(Z_0,\Gamma_0)$.

One may resolve prescribed proper analytic bad loci on the base during
this construction. After any finite sequence of further base
modifications, one can choose dominating smooth source models for which
the same exceptional assertion holds with respect to the original
reference $Z_0$.
\end{lemma}

\begin{proof}
Flatten first, before resolving the source. After resolving the new base
by further base change, denote the equidimensional strict transform by
$\bar g:\bar Z\to T$, and its modification to $Z_0$ by $\bar p$.
Resolve $\bar Z$ and all boundary data, and write
\[
 Z\xrightarrow{r}\bar Z\xrightarrow{\bar p}Z_0,
 \qquad p=\bar p\circ r,
 \qquad g=\bar g\circ r.
\]
Let $N=\dim Z_0$, $b=\dim T$, and $d=N-b$. If a prime $E$ on $Z$
has image contained in an analytic subset $C\subset T$ of codimension
at least two, equidimensionality gives
\[
 \dim\bar g^{-1}(C)\leq\dim C+d\leq N-2.
\]
It follows that $r(E)$, and therefore also $p(E)$, has dimension at most
$N-2$. Thus $E$ is exceptional over $Z_0$. The boundary and
orbifold-morphism assertions follow from \Cref{setup:log-modification}.
The new map has connected general fiber.
Its Stein factorization has a finite bimeromorphic map to the normal
base $T$, which is an isomorphism, so all its fibers are connected.

For any further base modification $T'\to T$, pull back the flat
intermediate family and take a common resolution dominating the previous
smooth source. The same dimension estimate proves exceptionality over
$Z_0$. This includes principalizations of prescribed bad loci and their
SNC resolutions on the base. Iterating proves the persistence assertion;
\Cref{setup:modification-compatibility}\textup{(ii)} identifies the
stepwise boundary with that computed directly over $Z_0$.
\end{proof}

The resolution need not be flat: the conclusion we retain is
exceptionality over $Z_0$, which permits
\Cref{setup:exceptional-descent}.

All constructions remain available in Fujiki class $\mathcal C$.
A smooth compact member admits a compact K\"ahler modification
\cite[Theorems~0.7 and~3.4]{DemaillyPaun2004}. Starting with one for
$Z_0$, graph resolutions and principalizations by smooth-center blowups
give a K\"ahler dominating source. K\"ahlerness is preserved by these
projective modifications, and proper meromorphic images remain in
class $\mathcal C$ \cite[Lemmas~4.4 and~4.6\textup{(3)}]{FujikiDouady}.
Thus the bases also admit K\"ahler models. The preceding dimension
argument still refers to the original $Z_0$.

\subsection{Keeping the original base fixed}

\begin{corollary}[Simultaneous model control]\label{setup:simultaneous}
Suppose a smooth pair $(X_1,\Delta_1)$ has a fibration $f_1:X_1\to Y$
to a fixed smooth base, and every prime divisor contracted by $f_1$ into
codimension at least two has coefficient one in $\Delta_1$. Put
$C=B(f_1,\Delta_1)$. Let a factorization through another smooth compact
space be given, after an allowed source modification if necessary, by
\[
 X_1\xrightarrow{g_1}W_1\xrightarrow{h_1}Y.
\]
Apply \Cref{setup:flattening} to $g_1$ and make any finite number of
subsequent modifications of $W_1$, always taking dominating smooth
source models with reduced added exceptional boundaries. The resulting
diagram $X\xrightarrow{g}W\xrightarrow{h}Y$, with source boundary
$\Delta$, can be chosen to have all the following properties:
\begin{enumerate}[label=\textup{(\roman*)}]
\item Every prime contracted by $g$ into codimension at least two is
exceptional over the chosen smooth reference source for $g_1$.
\item Every prime contracted by $f=h\circ g$ into codimension at least
two in $Y$ has coefficient one in $\Delta$.
\item $B(f,\Delta)=C$. In particular, for each prime $D\subset Y$ and
each prime $E\subset X$ dominating $D$,
\begin{equation}\label{setup:retained-multiplicity}
 \ord_E(f^*D)\,m_\Delta(E)\geq m_C(D).
\end{equation}
\item The log Kodaira dimensions of the total source and of its very
general fibers over $Y$ are unchanged by these source modifications.
\end{enumerate}
The hypothesis on contracted divisors follows from
\Cref{setup:flattening}, applied to the original map to its resolved
base before fixing $Y$.
\end{corollary}

\begin{proof}
Assertion (i) is \Cref{setup:flattening}. For (ii), a source prime
is either newly exceptional, with coefficient one, or the strict
transform of an old prime with unchanged image in $Y$. The hypothesis
handles the latter case. Assertions (iii) and (iv) are
\Cref{setup:modification-compatibility};
\eqref{setup:retained-multiplicity} is the defining minimum for $C$.
Flattening the original map supplies the initial coefficient-one
condition by making its contracted primes exceptional.
\end{proof}

In applications we include the critical values and images of boundary
strata among the base bad loci, so that $C$ has SNC support. We then
fix $Y$ and use \Cref{setup:simultaneous} for changes of intermediate
bases.

\subsection{The relative Iitaka map}

\begin{lemma}[Relative Iitaka construction]
\label{application:relative-iitaka}
Let \(f:(X,\Delta)\to Y\) be a fibration between smooth compact
manifolds in Fujiki class \(\mathcal C\), with \(\Delta\) a rational
SNC boundary in \([0,1]\). Suppose that
\[
 k=\kappa(F,K_F+\Delta_F)\geq 0
\]
on very general smooth fibers.  After modifications of the source, adding
the reduced exceptional divisor to the strict transform of the boundary,
there is a factorization
\[
 X\xrightarrow{g}W\xrightarrow{h}Y
\]
with smooth compact intermediate space in class \(\mathcal C\), both
maps proper with connected fibers, and
\[
 \dim W_y=k,\qquad
 \kappa(G,K_G+\Delta_G)=0
\]
for very general fibers of \(h\) and \(g\), respectively.  The factorization
restricts to the log Iitaka fibration on very general \(f\)-fibers.
These properties persist on higher models obtained by the same rule for
the source boundary.
\end{lemma}

\begin{proof}
Consider all degrees clearing the denominators of \(\Delta\).
Their relative adjoint direct images are coherent. First shrink to
the dense analytic open where the map and boundary strata are smooth.
There the family and its adjoint line bundles are flat over the base;
generic constancy of the fiber section dimension permits cohomology
and base change \cite[\S7, nos.~4--6, S\"atze~3--5]{GrauertDirectImages},
with its corrigendum \cite{GrauertDirectImagesCorr}.
Together with the evaluation ranks this gives a dense analytic open
for each degree;
outside the resulting countable collection of exceptional sets, the
maximum image dimension is \(k\), attained in one fixed degree.

Choose a sufficiently divisible degree giving the stabilized maps,
resolve the relative meromorphic map to the projective bundle of its
direct image, and take Stein factorization onto the image. The Iitaka
fibration theorem for line bundles on compact complex spaces identifies
the resulting map on very general fibers
\cite[Chapter~II, \S5; Chapter~III, \S7]{UenoClassification}.
No finite generation is required. The image has relative dimension \(k\); its
Stein factor is in class \(\mathcal C\), being a proper image of a space
in that class.  Resolve it and the resulting source map.  Connectedness
of the fibers of \(h\) follows from that of \(f\), since their images
under \(g\) are precisely the fibers of \(h\).

The assertion that the restriction has Kodaira dimension zero concerns
the original line bundle. Its standard proof applies here: coherent
direct images on the projective Iitaka image acquire sections after an
ample twist. If the
restriction had positive Iitaka dimension, multiplying those sections
by powers of the defining system would enlarge the Iitaka image.
A nonzero defining section restricts nontrivially to a very general
Iitaka fiber, so the restriction has dimension exactly zero.

Apply this to the resolved log adjoint on a very general \(f\)-fiber.
Its divisible section ring is unchanged by
\Cref{setup:log-modification}, and adjunction identifies the restriction
with \(K_G+\Delta_G\). Further source modifications preserve the same
section ring, hence the relative Iitaka factorization and its asserted
fiber dimensions.
\end{proof}

\subsection{The exact comparison for Kodaira-zero fibers}
\label{rankone:section}

For fibers of logarithmic Kodaira dimension zero, the relative
pluricanonical forms determine a single Hodge line on the base.  We
construct that line, compute its boundary orders, and use them to
identify the adjoint section spaces.  Coefficient-one boundary
components introduce mixed Hodge structures, but the line itself
lies in one pure weight grade.

A pure variation of Hodge structures carries a local system and a
varying Hodge filtration. The \emph{highest step} of a summand is its
highest nonzero filtration piece. The \emph{parabolic extension} of
such a line is obtained by local power substitutions making boundary
monodromy unipotent, extension of the unipotent Hodge filtration, and
descent with the resulting rational weights. We will compute those
weights using the actual relative pluricanonical form.

\begin{proposition}[The rank-one comparison]\label{rankone:comparison}
Let $g:(X,\Delta)\to W$ be a fibration with connected fibers from a
smooth compact K\"ahler manifold to a smooth compact manifold in
class $\mathcal C$.  Suppose that $\Delta$ is an SNC boundary with
rational coefficients in $[0,1]$ and that
$\kappa(G,K_G+\Delta_G)=0$ on a very general smooth fiber $G$.
After the modifications of
\Cref{setup:log-modification,setup:flattening}, fix a smooth compact
reference pair $(X_0,\Delta_0)$ and a modification
\[
 \pi:X\longrightarrow X_0,\qquad
 \Delta=\pi_*^{-1}\Delta_0+\operatorname{Exc}(\pi)_{\mathrm{red}},
\]
such that every prime divisor mapped by $g$ into codimension at least two
is $\pi$-exceptional.  Then there are a rational line bundle $M$
and a rational divisor $T$ on $W$ with the following properties.
\begin{enumerate}[label=\textup{(\roman*)}]
\item A positive integral multiple of $M$ is the parabolic highest Hodge
line of a complex direct summand of a pure, real-polarizable variation
with an integral lattice on a dense open subset of $W$.  Its parabolic extension is the line to which the adjoint-positivity
theorem, \Cref{period:adjoint}, will apply.
\item The divisor $T$ is supported on an SNC divisor and satisfies
\begin{equation}\label{rankone:boundary}
 B(g,\Delta)\leq T\leq 1.
\end{equation}
In particular, $T$ is effective.  If every prime over a prime $D$ of $W$
has $\Delta$-coefficient one, then $T_D=1$.
\item In sufficiently divisible degrees $q$, multiplication by the
relative generating forms gives linear isomorphisms
\begin{equation}\label{rankone:sections}
 H^0\bigl(W,q(K_W+T+M)\bigr)
 \simeq H^0\bigl(X,q(K_X+\Delta)\bigr)
 \simeq H^0\bigl(X_0,q(K_{X_0}+\Delta_0)\bigr).
\end{equation}
In a common divisible grading they preserve products and ratios of
sections, and hence the dimensions
of the images of their linear systems.
\item The same comparison holds relatively.  If $h:W\to Y$ and
$f_0:X_0\to Y$ are holomorphic maps with $f_0\pi=hg$, then
\begin{equation}\label{rankone:relative-sections}
 h_*\mathcal O_W\bigl(q(K_W+T+M)\bigr)
 \simeq (f_0)_*\mathcal O_{X_0}\bigl(q(K_{X_0}+\Delta_0)\bigr)
\end{equation}
for sufficiently divisible $q$.  The comparison of ratios also applies
to systems over $Y$ and to their restrictions wherever generic base
change holds. The relative identity remains valid after tensoring by a
line bundle on $Y$, or a rational line bundle in degrees clearing its
denominator.
\end{enumerate}
\end{proposition}

The Hodge construction and its divisorial calculation are independent
of the reference pair.  The condition on $\pi$ is used in the final
section comparison, to descend forms across exceptional centers.

\subsubsection{The root form and its eigenspace}

The cyclic-root construction and its distinguished eigensheaf occur in
\cite[Remark~2.6]{FujinoMori},
\cite[Lemma~5.2]{AmbroBoundary}, and
\cite[Remark~3.9]{FujinoGongyo}. Here the order calculations also
include coefficient-one boundary and keep the character line distinct
from the entire top Hodge space.

Choose a positive integer $m$ clearing the boundary denominators such
that the generic rank of
\[
 g_*\mathcal O_X\bigl(m(K_{X/W}+\Delta)\bigr)
\]
is one.  The same holds in every positive multiple of this degree:
powers of a nonzero section exist, whereas two independent sections
would contradict Kodaira dimension zero.  Its reflexive hull
$\mathcal E_m$ is a line bundle since $W$ is smooth.  A local frame $s$ of
$\mathcal E_m$ determines a meromorphic relative $m$-canonical form on
$X$ over the corresponding base open set.  Two such frames differ by
an invertible base function.
The direct image records $m$-pluriforms. Positivity will come from the
Hodge extension of their rational root; comparing the two extension
orders will supply $T$.

Resolve the horizontal zeros, poles, and boundary, then choose a dense
open set $W^\circ$ on which $g$ is smooth, the resulting divisors are
relatively SNC, and all their strata are smooth over the base.  By
\Cref{setup:log-modification}, new exceptional components have
coefficient one, and the logarithmic Kodaira dimension of the general
fiber stays zero.  On each base chart take the normalized cyclic cover
of $m$-th roots of $s$ as a relative pluriform and resolve it
functorially.  Write $\tau$ for its tautological relative top form and
$d=\dim X-\dim W$.

At a horizontal prime of boundary coefficient $\delta$, let $l$ be
the order of $s$ divided by $m$.
The inequality $l\geq-\delta\geq-1$ shows that $\tau$ has at most
logarithmic poles.  Indeed, at a covering prime of ramification index
$j$, its order is
\begin{equation}\label{rankone:root-order}
 j(1+l)-1.
\end{equation}
This number is integral.  If $l>-1$ it is nonnegative; if $l=-1$ it
equals $-1$.  On the resolved cover let $H$ be the reduced inverse
image of the horizontal primes with $l=-1$.  It includes the
exceptional components above those primes that are needed in the log
resolution.  Logarithmic pullback preserves the bound, and no
additional horizontal poles occur.  We use the cohomology of the
complement of $H$.

\begin{lemma}\label{rankone:eigenspace}
The tautological form spans the root-character part of
$F^dH^d(\widetilde G\setminus H,\C)$ on a general fiber of the cyclic
cover.  In particular, that part has dimension one.
\end{lemma}

\begin{proof}
For a compact K\"ahler SNC pair, logarithmic Hodge--de Rham degeneration
identifies this top Hodge space with its global logarithmic top
forms; it is also the smooth-fiber case of
\cite[Theorem~1.1]{FujinoFujisawaVMHS}, by duality.  Let $\eta$ be
another form of the tautological character.  Then $\eta/\tau$ is
deck-invariant and descends to a meromorphic function $v$ on $G$.

Consider a prime where the logarithmic section $s$ has no zero.
There $l=-\delta$.  For $l>-1$, the integer in
\eqref{rankone:root-order} lies between $0$ and $j-1$.  A pole of
$v$ downstairs would lower the covering order by at least $j$, which
is impossible.  For $l=-1$, the logarithmic pole allowance is already
exhausted, so the same conclusion holds.  Thus the poles of $v$ can
occur only on the positive zero divisor of $s$ as a section of
$m(K_G+\Delta_G)$.  A sufficiently high power $s^N$ absorbs those
poles.  Both $s^N$ and $v s^N$ are then sections of
$Nm(K_G+\Delta_G)$.  The assumption $\kappa=0$ forces them to be
dependent, and hence forces $v$ to be constant.

The argument also applies to a disconnected cover: we retain the whole
cover with its $\mu_m$-action, whose invariant meromorphic functions
descend to $G$.
\end{proof}

The ranks of the Hodge and weight graded bundles are locally constant
on $W^\circ$.  The conclusion of \Cref{rankone:eigenspace}, proved on
a very general fiber, therefore holds throughout this smooth log
locus.  Notice that neither the argument nor
\eqref{rankone:root-order} requires a boundary coefficient to have
the form $1-1/a$ with $a$ an integer.

\subsubsection{Pure weights and descent of the local systems}

We next place a power of the local Hodge line in a global pure
variation.  This requires compatible polarizations and descent,
because the cyclic covers need not form a global family.

First make the local normalized-cover constructions on charts of the
compact base $W$, using the meromorphic forms supplied by the frames
of $\mathcal E_m$.  They extend across the missing base divisor as
finite covers; subsequent resolutions can be projective and
functorial in the analytic category
\cite[Theorem~1.1]{BierstoneMilmanFunctorial}.
The resulting maps to the K\"ahler source are projective, since the
finite covers are projective. Thus their total spaces are K\"ahler
over relatively compact source neighborhoods
\cite[Definition~4.1 and Lemma~4.4]{FujikiDouady}. Properness allows
such a neighborhood to contain the inverse image of a sufficiently
small disk transverse to a base divisor.
The divisor of the tensor is unchanged under multiplication
by a base unit.  Thus the constructions on overlaps are identified
after taking a local root of that unit, and the resolutions and their
exceptional divisor classes are identified as well.

Fix a K\"ahler class $\xi$ on $X$.  On each resolved-cover chart
$q_i:Z_i\to X|_{U_i}$, choose the relatively ample line bundle $A_i$
as follows.  Functoriality matches the blowup center ideals and their
tautological line bundles under the root-change identifications.
Insert identity stages where a center is empty; the corresponding
tautological bundle is trivial there.  A sufficiently weighted tensor
product gives compatible $A_i$, equivariant under deck transformations,
by descent of the centers, blowups, and tautological bundles.  The classes
\[
 \theta_i=Nq_i^*\xi+c_1(A_i)
\]
are K\"ahler on the smooth compact fibers for $N$ sufficiently large;
their restrictions to the compact boundary strata are K\"ahler too.
At the intermediate finite covering stage one uses the pullback
K\"ahler \emph{class}; the literal pulled-back form can degenerate
at ramification.  The relatively ample class supplies positivity
along the projective resolution.  One may choose a single $N$:
carry out the construction over charts $U_i^+$ containing the closures
of finitely many relatively compact charts $U_i$ covering $W$.
Properness makes the inverse images of these closures compact, so
each admits a bound $N_i$ in this class construction.  Taking
$N\geq\max_iN_i$ works on all charts, since increasing $N$ adds
the semipositive class $q_i^*\xi$.  This is why the construction was
extended over the compact base before shrinking to $W^\circ$.
For the invariant K\"ahler-class construction on an equivariant
resolved cyclic cover, see also \cite[Lemma~8.1]{ChenLimits}.

The classes $\theta_i$ come from total-space cohomology, so their
restrictions to the smooth fibers and strata are flat for the
Gauss--Manin connection.  They agree under the overlap
identifications.  Lefschetz decomposition and cup product consequently
give compatible flat real polarizations on the compact-stratum
cohomology systems.  The weight spectral sequence, with its Tate
twists, gives the pure weight grades of the open-fiber cohomology as
subquotients of these systems.  They are real-polarizable: a real Hodge
subvariation of a polarized pure variation is nondegenerate for the
polarization, and a quotient can be identified with an orthogonal
complement.  These operations are compatible with the overlap maps.
The integral structures come from cohomology modulo torsion and the
saturated intersections with the rational weight filtration.  We do
not assert that the real polarizations are rational.

There is a unique weight $w$ for which
\[
 F^d\Gr^W_w H^d(\widetilde G\setminus H,\C)_\chi\ne0.
\]
Here $\chi$ is the tautological character.  Uniqueness follows from
the rank-one assertion and strictness of the weight filtration.  Let
$V_i$ be the pure integral variation given by this weight grade on
the $i$th chart; the chosen line is $F^d(V_i)_\chi$.

\begin{lemma}\label{rankone:descent}
A positive tensor power of these local Hodge lines is the highest
Hodge line of a complex direct summand of a global real-polarizable
integral pure variation on $W^\circ$.
\end{lemma}

\begin{proof}
Refine an overlap so that a root of its transition unit can be
chosen.  It gives an isomorphism between the cyclic-cover families,
and therefore an isomorphism of their integral cohomology systems.
Another choice differs by a deck transformation.  On a triple
overlap the cocycle defect is also a deck transformation.  All these
maps commute with the distinguished $\mu_m$-action.

Put
\[
 \mathbb U_i=
 \bigl((V_{i,\Z}/\text{torsion})^{\otimes m}\bigr)^{\mu_m},
\]
using the diagonal deck action.  The ambiguity and the cocycle defect
act trivially on $\mathbb U_i$.  Hence the $\mathbb U_i$ descend to
an integral local system $\mathbb U$; saturation only changes the
lattice by an isogeny.  The compatible real polarizations, tensorized
and restricted to invariants, polarize this pure variation.

Over $\C$ the character projector $e_\chi$ is a flat Hodge
endomorphism on each chart.  The projector
$e_\chi^{\otimes m}$ restricts to the invariant tensor space,
because $\chi^m=1$.  It commutes with all transition maps and
therefore defines a global complex direct summand with local fibers
$((V_i)_\chi)^{\otimes m}$.  Its highest step is
$(F^d(V_i)_\chi)^{\otimes m}$, a line.  In local rooted notation
this line is generated by $\tau^{\otimes m}$, which is identified
meromorphically with $s$.  Thus no global root of $\mathcal E_m$ and
no global family of cyclic covers is needed.
\end{proof}

\subsubsection{The two extension orders}

Let $M$ denote one $m$th of the parabolic extension of the line just
constructed.  We use $s^{1/m}$ for its local rational frame; the
notation becomes literal after the powers and finite boundary covers
above. We now compare this parabolic extension with the pole bounds
of the original pluriform.

Fix a prime divisor $D$ on $W$ and use a local ordinary base volume
$\omega$ that does not vanish at its generic point. Work on a source
model on which the divisor of the generator, the boundary, and the
inverse image of $D$ have SNC support over that generic point,
using strict transforms plus reduced exceptional boundary as before.
All primes of this prepared source dominating $D$, including the
exceptional primes introduced by its resolution, enter the following
order tests. Put
\[
 l_E=\frac1m\ord_E(s\wedge g^*\omega^m),\qquad
 a_E=\ord_E(g^*D),\qquad \delta_E=\Delta_E,
\]
and
\[
 \alpha_D=\min_{g(E)=D}\frac{l_E+\delta_E}{a_E},
 \qquad \beta_D=\ord_D^M(s^{1/m}).
\]
The notation $s\wedge g^*\omega^m$ means the relative-times-base
identification in the $m$th tensor power of the canonical bundle.
For $q$ divisible by $m$ and clearing the rational data, and a
meromorphic base function $\varphi$, the total pluriform
$\varphi\,\omega^q s^{q/m}$ has boundary-adjusted order
\[
 a_E\ord_D(\varphi)+q(l_E+\delta_E)
\]
at $E$. It therefore satisfies the source boundary bounds at all
these primes if and only if
\begin{equation}\label{rankone:base-bound}
 \ord_D(\varphi)+q\alpha_D\geq0.
\end{equation}
The Hodge frame contributes $q\beta_D$ to the order on the base,
so the candidate boundary correction is $T_D=\alpha_D-\beta_D$.
Computing $\beta_D$ will reduce the required bounds on $T_D$ to
a comparison of two minima. We first justify computing that Hodge
order by logarithmic orders on a semistable cover, without changing
the chosen pure Hodge line. Global consistency of $T$ and descent
across the untested source primes will follow afterward.

\subsubsection{The extension across a transverse disk}

We use the analytic canonical-extension theorem
\cite[Theorem~1.1]{FujinoFujisawaVMHS}.  It applies to a proper
surjective morphism of an analytic SNC pair $(Z,H)$ to a smooth base,
with $H$ reduced, every stratum K\"ahler and dominant, and all strata
smooth off a normal crossing discriminant.  It supplies a
graded-polarizable real variation on compact-support cohomology,
locally free double Hodge/weight grades on the extension, and the
dual direct-image description by $\omega_{Z/B}(H)$.  Admissibility
is established in \cite[Remark~4.8]{FujinoFujisawaVMHS}.
These conclusions apply to proper K\"ahler families; projectivity
is not required.

Take a transverse disk at a general point of a base divisor. Resolve
the cyclic-cover pair together with the central fiber, leaving its
smooth logarithmic locus over the punctured disk unchanged. Near a
central point the resulting map has the form
\[
 t=\varepsilon x_1^{a_1}\cdots x_r^{a_r},
 \qquad H\subset\{x_{r+1}\cdots x_l=0\},
\]
where $\varepsilon$ is a unit and $H$ is the horizontal boundary. Absorbing it
into a vertical coordinate gives a strict toroidal morphism. Compactness
of the central fiber and properness allow us to shrink the disk until
finitely many such charts cover its inverse image. The associated cone
complex is therefore finite. Apply the combinatorial semistable
subdivision theorem
\cite[Theorem~2.7]{AdiprasitoLiuTemkinSemistable}. Its base is one ray,
so the base lattice alteration is realized by $t=u^N$. Normalize this
base change, keeping all components, and realize the projective source
subdivision analytically
\cite[Lemmas~2.14--2.15, 2.18, and~2.31(4)]
{EnokizonoHashizumeSemistable}. Zero-image faces, which record the
horizontal boundary, are allowed in the combinatorial theorem.

The semistable local form is
\[
 u=y_1\cdots y_q,
 \qquad H\subset\{y_{q+1}\cdots y_k=0\};
\]
see \cite[Lemma~2.21(4)]{EnokizonoHashizumeSemistable}.
Here $H$ includes every horizontal exceptional component in the
transformed boundary. The total space is smooth, the central fiber is
reduced, and its union with $H$ is SNC. On every horizontal
intersection stratum the displayed monomial is nonconstant;
properness makes its image the whole disk. The composite to the
initial K\"ahler total space is projective, so the new total space and
its smooth strata are K\"ahler after shrinking the disk
\cite[Lemma~4.4]{FujikiDouady}.

These toroidal modifications preserve the open pair on the punctured
fibers. Its cohomology is therefore the pullback of the original
open-fiber cohomology, including the full direct sum if the cover
splits. Transport the original $\mu_m$-action through this
identification. The deck projector then acts on the variation and,
by functoriality, on its canonical extension; no action on a selected
component or on the chosen semistable model is needed. We apply the
extension theorem componentwise and take the direct sum. We use this
semistable pair itself: an arbitrary further resolution need not
preserve reducedness.

For application of the cited theorem, $H$ consists of the horizontal
logarithmic boundary.  The central fiber is not included in $H$:
it appears in the relative logarithmic complex, with $dt/t$ in the
base quotient.  Write $f:Z\to B$ for this degeneration over the disk,
and $Z_0=f^{-1}(0)$ for its central fiber.  The top relative
logarithmic sheaf is
\begin{equation}\label{rankone:log-volume}
 \omega_{Z/B}\bigl(H+(Z_0)_{\mathrm{red}}-f^*[0]\bigr).
\end{equation}
On the semistable model the central fiber is reduced, so this is
$\omega_{Z/B}(H)$.  For a smooth open fiber of dimension $d$,
Poincar\'e duality identifies
\[
 H^d_c(Z_t\setminus H_t)^*\simeq
 H^d(Z_t\setminus H_t)(d).
\]
Accordingly the dual degree-zero Hodge quotient in
\cite[Theorem~1.1(iv)]{FujinoFujisawaVMHS} is precisely the top
$F^d$ step of ordinary $H^d$.  The upper and lower canonical
extensions coincide after unipotentization.  This proves that the
canonical extension of the top step is the direct image of
\eqref{rankone:log-volume}.  The K\"ahler form of the
one-parameter limiting theorem is also explained in
\cite[Remarks~2.16, 4.4, and~4.8]{FujinoFujisawaVMHS}.

\begin{lemma}\label{rankone:extension}
The projection of the rank-one top eigenspace to its nonzero pure
weight grade is an isomorphism on canonical extensions after
unipotentization.  Consequently it is an isomorphism of parabolic
rational lines before that base change.
\end{lemma}

\begin{proof}
Let $\overline{\mathcal V}_\chi$ be the extended mixed eigensystem
over the disk.  The finite deck projector acts on the canonical
extension and its filtrations.  By the cited extension theorem,
each
\[
 \Gr_F^p\Gr^W_j\overline{\mathcal V}_\chi
\]
is locally free and has its rank on the punctured disk.  The
filtrations on the weight quotients are induced filtrations.
Thus the exact sequences
\[
 0\longrightarrow F^dW_{j-1}\overline{\mathcal V}_\chi
 \longrightarrow F^dW_j\overline{\mathcal V}_\chi
 \longrightarrow F^d\Gr^W_j\overline{\mathcal V}_\chi
 \longrightarrow0
\]
have locally free quotients.  There is no step above $F^d$, also
on the extension, by these constant-rank assertions.  Exactly one
of the displayed quotients has rank one, namely that for $j=w$;
all the others have rank zero.  Induction on the weight filtration
therefore gives
\[
 F^dW_{w-1}\overline{\mathcal V}_\chi=0,\qquad
 F^dW_w\overline{\mathcal V}_\chi
   =F^d\overline{\mathcal V}_\chi,
\]
and identifies the latter line with
$F^d\Gr^W_w\overline{\mathcal V}_\chi$ across the origin, without
an additional zero or pole.  Descent gives the parabolic assertion.
\end{proof}

\subsubsection{The divisorial order of the Hodge line}

With the notation and prepared source above, we claim that the
parabolic order of the frame $s^{1/m}$ in $M$ is
\begin{equation}\label{rankone:order}
 \beta_D=\ord_D^M(s^{1/m})
     =\min_{g(E)=D}\frac{l_E+1-a_E}{a_E}.
\end{equation}
Equivalently, $\beta_D$ is the largest rational number $b$ such that
$t^{-b}s^{1/m}$ extends at $D=(t=0)$ in the parabolic sense.

It suffices to work over a transverse disk, suppressing the other
base coordinates.  Write $\tau=s^{1/m}$ and
$\Omega=\tau\wedge g^*dt$.  The multivalued total form
$\Omega/t$ has order $l_E-a_E$ at $E$.  A logarithmic extension
allows order $-1$, so twisting by $t^{-b}$ imposes
\begin{equation}\label{rankone:inequality}
 l_E+1-a_E-a_Eb\geq0.
\end{equation}
To verify this directly against the canonical extension, take a
finite base change $t=u^N$ and a semistable model of the root
cover.  At a central prime over the strict transform of $E$, let
$e$ be the total ramification index of the composite map to $X$,
including the cyclic-cover ramification.  Reducedness gives
$ea_E=\ord_{E'}(t)=N\ord_{E'}(u)=N$, where $E'$ is this new prime.
Orders of a total volume form shifted by one multiply by $e$.
Dividing its pullback by
$dt/du=Nu^{N-1}$ therefore gives
\[
 \ord(\tau\wedge du)
 =e(l_E+1)-N
 =e(l_E+1-a_E).
\]
The base twist subtracts $Nb=ea_Eb$.  Regularity in
\eqref{rankone:log-volume} is exactly
\eqref{rankone:inequality}.  By \Cref{rankone:extension}, testing
the logarithmic top form tests the chosen pure Hodge line as well.
The form being tested is the pullback of the original tautological
form on the full cover. The identification of open-fiber cohomology
above and the direct-image description of its canonical extension
identify it with that same Hodge section. Keeping all normalized
components ensures that a prime above every original $E$ is present;
the proper toroidal modifications retain its strict transform.

To prove sufficiency, continue to write $X$ for the smooth source
over the transverse disk and put $n=d+1$.  In SNC coordinates write
\[
 \Omega=h\prod_i x_i^{l_i}\,dx_1\wedge\cdots\wedge dx_n,
 \qquad t=h'\prod_i x_i^{a_i},
\]
where $h,h'$ are units and $a_i=0$ for horizontal components.
Put $c_i=l_i+1-(1+b)a_i$.  The vertical inequalities
\eqref{rankone:inequality} and the horizontal logarithmic bounds
give $c_i\geq0$.  For any additional divisorial valuation $v$,
write $A_X(v)$ for its log discrepancy over this smooth model,
that is, one plus the order of the Jacobian.  Logarithmic pullback
of the SNC coordinate divisor gives
$A_X(v)\geq\sum_i v(x_i)$, exactly as in
\Cref{setup:log-modification}.  Thus the shifted order of
$t^{-(1+b)}\Omega$ is
\[
 A_X(v)+\sum_i\bigl(l_i-(1+b)a_i\bigr)v(x_i)
 \ \geq\ \sum_i c_i v(x_i)\ \geq\ 0.
\]
These inequalities persist under finite covers, since shifted orders
multiply by ramification indices.  Further exceptional primes therefore
impose no stronger condition.
Horizontally, shifted orders are strictly positive outside the
designated log boundary; on the root cover these are positive
integers, so there are no poles there.  Conversely, covering primes
above strict transforms detect every
inequality in \eqref{rankone:inequality}.  Taking their minimum
proves \eqref{rankone:order}.

\begin{remark}\label{rankone:collision}
The use of all primes on the resolved model in
\eqref{rankone:order} matters even for a simple logarithmic
family.  On $\mathbf P^1_x$ over a disk, remove the horizontal
divisor $x^2=t$ and use $\tau=dx/(x^2-t)$.  After $t=u^2$ and
$x=uz$, one has $\tau=u^{-1}dz/(z^2-1)$, so the parabolic order
is $-1/2$.  Resolving the tangency of the horizontal divisor with
$t=0$ introduces a prime with $a_E=2$ and $l_E=0$, which gives
exactly $-1/2$ in \eqref{rankone:order}.  The original central
prime alone would not detect it.
\end{remark}

\subsubsection{The boundary and the comparison of sections}

Define
\begin{equation}\label{rankone:alpha-T}
 T_D=\alpha_D-\beta_D.
\end{equation}
Changing $s$ to $v s$, for a meromorphic base function $v$, adds
$a_E\ord_D(v)/m$ to $l_E$.  Hence it adds
$\ord_D(v)/m$ to both $\alpha_D$ and $\beta_D$ and leaves
$T_D$ unchanged.  This also checks the transformation rule for
the rational Hodge frame in \eqref{rankone:order}.  Passing to
a higher divisible generating degree has the same effect: on the
generic fiber its generator is a base multiple of a power of $s$.
Thus the construction is independent of these choices.

For the boundary inequality, write
\[
 x_E=\frac{l_E+\delta_E}{a_E},\qquad
 c_E=1-\frac{1-\delta_E}{a_E}.
\]
Then $0\leq c_E\leq1$, and
\[
 T_D=\min_E x_E-\min_E(x_E-c_E).
\]
For any finite collections of real numbers $x_E,c_E$, this
difference lies between $\min_Ec_E$ and $\max_Ec_E$.
Moreover the inf-multiplicity convention gives
\[
 \min_Ec_E
 =1-\max_{g(E)=D}\frac{1-\delta_E}{a_E}
 =B(g,\Delta)_D.
\]
This proves \eqref{rankone:boundary}.  If all $\delta_E=1$,
all $c_E=1$, so $T_D=1$.

Choose the discriminant to include the vertical boundary and all
loci where the smooth log construction fails.  Off it the family
and the generator line are smooth log-relative data, and there is
no vertical boundary.  There $\alpha_D=\beta_D$; in a
nonvanishing generator frame both are zero.  Thus $T$ has finite
support in the chosen discriminant, which may be resolved to be
SNC.  On subsequent modifications one pulls back the variation and
its parabolic line and recomputes $\alpha_D$ and
\eqref{rankone:alpha-T}; the order computation is unchanged on the
resulting models.

\begin{proof}[Completion of the proof of \Cref{rankone:comparison}]
Take $q$ divisible by $m$ and by the denominators of all rational
data.  In a local base frame, a meromorphic section of
$q(K_W+T+M)$ is represented
by
\[
 \varphi\,\omega^q s^{q/m}.
\]
The order of the frame in $T+M$ is
$q(T_D+\beta_D)=q\alpha_D$.  Consequently this section is
regular at $D$ exactly when \eqref{rankone:base-bound} holds.
The source-side order test already proved identifies this with
the boundary bounds at every source prime dominating $D$.
Along horizontal primes the required bound is automatic because
$s^{q/m}$ is a logarithmic
pluricanonical section on the general fiber.  We have proved the
exact section comparison in codimension one on the base, and at
all source primes except those contracted by $g$ into codimension
at least two.

Conversely, a section of $q(K_X+\Delta)$ restricts on a general
$g$-fiber to a multiple of $s^{q/m}$.  Their ratio is a
meromorphic function constant on the connected general fibers,
and therefore descends meromorphically to $W$.  The same
divisorial inequalities make it a section of $q(K_W+T+M)$
outside codimension two; smoothness of $W$ extends it uniquely
across that set.  This proves the converse comparison in all
divisible degrees under consideration.

For the forward direction, a section on $W$ gives a meromorphic
adjoint form on $X$ whose only possible remaining poles lie on
$g$-contracted prime divisors.  By hypothesis these divisors are
$\pi$-exceptional.  Push the form to $X_0$.  At every prime of
$X_0$ its transform was among the tested primes of $X$, so the
pushed form satisfies the $\Delta_0$ pole bound in codimension
one.  Since $X_0$ is smooth and the adjoint multiple is a line
bundle, it extends across codimension two.  Finally
\Cref{setup:log-modification} pulls it back to a section of
$q(K_X+\Delta)$, proving \eqref{rankone:sections}.

If the diagram is over $Y$, apply this argument over an arbitrary
open subset $U\subset Y$.  Its inverse image in $X$ is
$\pi^{-1}(f_0^{-1}U)=g^{-1}(h^{-1}U)$, and $\pi$ is still
proper over the smooth open space $f_0^{-1}U$.  The exceptional
images still have codimension at least two.  The construction and
the extension of sections commute with restriction to $U$,
giving \eqref{rankone:relative-sections}.  Relative canonical
versions, when the auxiliary base is smooth, follow by the
projection formula.  All comparisons multiply by the same relative
generator, with its powers in the common divisible grading, so they
preserve products and ratios of sections. The projection formula also
gives the identity after any pulled-back base twist, in degrees clearing
the denominator for a rational twist. This completes
\Cref{rankone:comparison}.
\end{proof}

\subsection{Retaining the original base multiplicities}

\begin{lemma}[Composition of inf-multiplicity bounds]
\label{application:composition}
Suppose that \(X\xrightarrow{g}W\xrightarrow{h}Y\) is a diagram of
surjective maps between smooth spaces, and \(\Delta,T,C\) are rational
boundaries on these spaces.  If
\[
 T\geq B(g,\Delta),\qquad C\leq B(h\circ g,\Delta),
\]
then \(B(h,T)\geq C\).  If, in addition, every prime divisor of \(X\)
mapped into codimension at least two in \(Y\) has coefficient one in
\(\Delta\), and \(T\leq1\), then every prime divisor of \(W\) mapped
into codimension at least two in \(Y\) has coefficient one in \(T\).
\end{lemma}

This is the inf-multiplicity composition comparison of
\cite[Proposition~3.13(1)]{CampanaOrbifolds}; we include its order
calculation to track coefficient-one exceptional divisors.

\begin{proof}
Let \(D\subset Y\) and \(V\subset W\) be prime divisors with
\(h(V)=D\).  Write \(b_V=\ord_V(h^*D)\).  If \(E\subset X\) is a
prime divisor dominating \(V\), put \(a_E=\ord_E(g^*V)\).  At their
generic points orders multiply, so
\[
 \ord_E((h\circ g)^*D)=b_Va_E.
\]
The hypothesis on \(C\) implies
\(b_Va_E m_\Delta(E)\geq m_C(D)\) for every such \(E\).  Taking the
minimum over them and then using the hypothesis on \(T\) gives
\[
 b_Vm_T(V)\ \geq\ b_Vm(g,\Delta;V)\ \geq\ m_C(D).
\]
Taking the minimum over \(V\) proves the first assertion.  All these
inequalities are valid with the prescribed convention for infinity.

For the second assertion, let \(V\) be mapped into codimension at least
two in \(Y\).  Every prime \(E\) dominating it has
\(\Delta_E=1\), so \(m(g,\Delta;V)=\infty\).  It follows that
\(B(g,\Delta)_V=1\), and hence \(T_V=1\).
\end{proof}

\section{The adjoint-addition principle}\label{addition:section}

The section comparison reduces subadditivity to a statement about
$K_W+T+M$. We prove that statement here from the two positivity
inputs specified below. Its two maps have different purposes:
$h$ is the fibration whose fiber dimension is to be added, whereas
$p$ expresses the nef term as a pullback from a projective base.

The first input is the precise adjoint-positivity theorem for the
line produced by the construction. We use the highest-step and
parabolic-extension conventions of Section~\ref{rankone:section}.

\begin{theorem}[Adjoint-positivity input]\label{period:adjoint}
Let $B$ be a smooth compact connected manifold in class $\mathcal C$
and $B^\circ\subset B$ a dense Zariski open set with SNC complement.
Let $\mathbb V$ be a pure real-polarizable variation of Hodge
structures on $B^\circ$ with an integral ambient lattice. Suppose
a complex direct summand of $\mathbb V_{\mathbf C}$ has a
rank-one highest nonzero Hodge step, with parabolic rational
extension $L$. There are a modification $\tau:B'\to B$, a
surjective morphism $p:B'\to S$ with connected fibers to a smooth
projective variety, and a nef rational line bundle $P$ on $S$ such that
\begin{equation}\label{period:conclusion}
 \tau^*L=p^*P,\qquad K_S+aP\text{ is big for some }a\in\mathbf Q_{>0}.
\end{equation}
Positive rational rescaling of $L$ is allowed, and the line identity
persists under further modifications. If $S$ is a point, $\tau^*L$
is rationally trivial.
\end{theorem}

The full proof is in Section~\ref{period:section}. The lattice belongs to the ambient
variation: the complex summand need not itself be defined over
$\mathbf Q$, and the real polarization need not be rational.

The second input supplies addition when the fiber already has a big
log canonical divisor. It will be proved in Section~\ref{generaltype:section}.

\begin{proposition}\label{generaltype:addition}
Let \(h\colon V\to Y\) be a surjective holomorphic map with connected fibers
between smooth compact complex manifolds in Fujiki class \(\mathcal C\).
Let \(T\) and \(C\) be effective rational divisors with simple normal crossing
support and coefficients in \([0,1]\) on \(V\) and \(Y\), respectively.
Suppose that
\begin{enumerate}[label=\textup{(\roman*)}]
\item \(K_F+T_F\) is big for a very general smooth fiber \(F\) of \(h\);
\item \(C\leq B(h,T)\);
\item \(T_E=1\) for every prime divisor \(E\) of \(V\) whose image in \(Y\)
has codimension at least two.
\end{enumerate}
Then
\begin{equation}\label{generaltype:inequality}
  \kappa(V,K_V+T)\ \geq\
  \dim F+\kappa(Y,K_Y+C).
\end{equation}
As usual, the assertion is automatic if the last term is \(-\infty\).
\end{proposition}


General-type addition will give sections after a positive ample
twist from $S$. To remove that twist, we will need one fixed section
with a negative ample twist. The next lemma produces it from a big
base divisor, once nonvanishing on the general $p$-fiber is known.
It uses weak positivity with exceptional poles and
descent to a fixed smooth reference, following the method of
\cite[Remark~6.5(3) and Proposition~6.6]{CampanaOrbifolds}.
We include the descent step explicitly.

\begin{lemma}[Effectivity after a big base twist]
\label{application:weak-effectivity}
Let $p:(W,T)\to S$ be surjective, with $W$ smooth compact in
class $\mathcal C$, $S$ smooth projective, and $T$ a rational SNC
boundary. Suppose one fixed divisible relative log pluricanonical
system is nonzero on very general fibers. For every big rational
line bundle $H$ on $S$, the rational line bundle
$K_{W/S}+T+p^*H$ has a nonzero section in some positive degree.
\end{lemma}

\begin{proof}
Choose a K\"ahler source modification, flatten over a smooth
projective modification $\pi:S'\to S$, and resolve the flat main
transform. Write $\mu:W'\to W$ and $p':W'\to S'$ for the
resulting maps. Give $W'$ the strict transform of $T$ plus the
reduced $\mu$-exceptional divisor, denoted $T'$. Then
\begin{equation}\label{application:source-discrepancy}
 K_{W'}+T'=\mu^*(K_W+T)+E,
 \qquad E\geq0\ \text{and $\mu$-exceptional}.
\end{equation}
Every source prime over a subset of codimension at least two in
$S'$ is exceptional over the fixed $W$: before resolving the
source, flatness bounds that inverse image by codimension at least
two. Put $R=K_{S'/S}\geq0$ and $H'=\pi^*H+R$. The exact identity
\begin{equation}\label{application:twist-cancellation}
 K_{W'/S'}+T'+p'^*H'
 =\mu^*(K_{W/S}+T+p^*H)+E
\end{equation}
will allow descent.

Choose a divisible $m>0$ such that
$\mathcal E=p'_*\mathcal O_{W'}(m(K_{W'/S'}+T'))$ has positive
rank. Generic base change and the hypothesis provide such an $m$.
Fujino's weak positivity theorem applies to this log canonical pair
and its projective base \cite[Theorem~1.1]{FujinoWeakPositivity}.
Choose an ample Cartier divisor $A_0$ and an integer $\alpha>0$
such that $H'-A_0/(m\alpha)$ is big. For some integer $\beta>0$,
\[
 (\operatorname{Sym}^{\alpha\beta}\mathcal E)^{**}
        \otimes\mathcal O_{S'}(\beta A_0)
\]
is generated by global sections at the generic point.

On the locally free locus of $\mathcal E$, multiplication of
relative sections gives an evaluation map to
\[
 \mathcal O_{W'}\bigl(m\alpha\beta(K_{W'/S'}+T')
                         +\beta p'^*A_0\bigr).
\]
It is nonzero generically, because a nonzero relative section has
nonzero powers. Generic generation supplies a global section whose
image is nonzero. The complement of the locally free locus has
codimension at least two. Multiply a sufficiently divisible power
of the evaluated section by a section of a sufficiently divisible
multiple of
$m\alpha\beta\,p'^*(H'-A_0/(m\alpha))$.
The latter exists by bigness. This produces a nonzero section of a
multiple of the left side of \eqref{application:twist-cancellation}, except possibly
at source primes above that codimension-two set.

Locally, a meromorphic frame of $\mathcal E$ expresses the evaluated
section as a meromorphic combination of products of relative
sections. Thus the resulting form is meromorphic, and its only
possible divisorial poles lie above that omitted codimension-two set.
Those primes, and $E$, are exceptional over $W$. At every prime
of the original smooth $W$, the resulting meromorphic form therefore
has the required order for a multiple of $K_{W/S}+T+p^*H$.
Divisorial regularity and extension across codimension two give a
global section on $W$. If $S$ is a point, the same conclusion is
the assumed nonvanishing.
\end{proof}

\begin{proposition}[Adjoint addition]
\label{addition:principle}
Let $h:W\to Y$ be a fibration between smooth compact manifolds in
class $\mathcal C$. Let $T,C$ be rational SNC boundaries with
$C\leq B(h,T)$, and suppose every prime of $W$ whose image has
codimension at least two in $Y$ has coefficient one in $T$.
Let $p:W\to S$ be a surjective morphism with connected fibers to
a smooth projective variety. Suppose $P$ is a nef rational line
bundle on $S$, $M=p^*P$, and $K_S+a_0P$ is big for some rational
$a_0>0$. If $K_W+T+M$ is big on very general $h$-fibers, then
\[
 \kappa(W,K_W+T+M)
 \geq\dim(W/Y)+\kappa(Y,K_Y+C).
\]
\end{proposition}

\begin{proof}
The assertion is automatic if the last term is $-\infty$.
Assume it is nonnegative, put $D=K_W+T$, and write
$k=\dim(W/Y)$. We may first replace $W$ by a K\"ahler
modification, adding the reduced exceptional divisor to its strict
boundary. The log section spaces, also after tensoring by the
pullback of $M$, are unchanged. Fiberwise bigness and the
inf-multiplicity boundary are preserved; newly exceptional primes
have coefficient one. Thus all hypotheses remain valid.

If $S$ is a point, $M$ is rationally trivial and
\Cref{generaltype:addition} proves the proposition. Suppose $\dim S>0$,
and choose an ample Cartier divisor $A$ on $S$. A very general
$h$-fiber is K\"ahler and carries the big rational line
$(D+M)|_{W_y}$, so it is projective. Openness of its big cone
gives a rational $c$ with $0<c<1$ for which
$(D+cM)|_{W_y}$ is still big. This choice works on very general
fibers: choose the first fiber outside the countably many
cohomology and evaluation-rank exceptional loci for rational $c$
and divisible degrees. The one full-rank system just found then
has the same rank generically. For zero-dimensional fibers this
argument is unnecessary.

For every positive rational $\eta$, $cP+\eta A$ is ample.
Choose a sufficiently divisible large $N$ and a general member
of the basepoint-free system $|Np^*(cP+\eta A)|$. Dividing by
$N$ gives an effective rational divisor $Q_\eta$ such that
\[
 Q_\eta\sim_{\mathbf Q}cM+\eta p^*A,
 \qquad T+Q_\eta\ \text{is an SNC boundary}.
\]
Bertini is applied simultaneously to the finitely many strata of
$T$; the member contains no old boundary component, and increasing
$N$ bounds its new coefficients by one. Existing coefficient-one
components remain unchanged, and increasing the source boundary
can only increase $B(h,T)$. Thus all boundary hypotheses of
\Cref{generaltype:addition} hold for $T+Q_\eta$. Its fiber adjoint is big
by the choice of $c$ and nefness of the added ample pullback.
Consequently
\begin{equation}\label{application:perturbed-bound}
 \kappa(W,D+cM+\eta p^*A)
       \geq k+\kappa(Y,K_Y+C)
       \qquad(\eta\in\mathbf Q_{>0}).
\end{equation}

The systems in \eqref{application:perturbed-bound} already have the required
image dimension. It remains to remove their ample term without
losing these ratios of sections. In particular one such divisor has
a nonzero section in a fixed
positive degree. Restriction of that section to a very general
$p$-fiber is nonzero; both pulled-back twists become trivial
there. Hence
\begin{equation}\label{application:p-fiber-nonvanishing}
 \kappa(W_s,(K_W+T)|_{W_s})\geq0
 \quad\text{for very general }s\in S.
\end{equation}
One fixed relative log pluricanonical direct image for
$p:(W,T)\to S$ therefore has positive rank.

Choose rational $a>\max\{1,a_0\}$. Since $P$ is nef,
$K_S+aP$ is big. For sufficiently small rational $\lambda>0$,
$H=K_S+aP-\lambda A$ is still big.
Lemma~\ref{application:weak-effectivity} now gives a nonzero section of a
positive multiple of
\begin{equation}\label{application:negative-twist}
 E_-=D+aM-\lambda p^*A
     =K_{W/S}+T+p^*H.
\end{equation}
Put
\[
 \theta=\frac{1-c}{a-c},
 \qquad \eta=\frac{(1-c)\lambda}{a-1},
 \qquad E_+=D+cM+\eta p^*A.
\]
Then $0<\theta<1$, $\eta>0$, and direct calculation gives
\begin{equation}\label{addition:cancel}
 D+M=(1-\theta)E_++\theta E_-.
\end{equation}
Fix a nonzero section $e$ of $rE_-$ for one denominator-clearing
$r>0$. For sufficiently divisible $n$, multiplication by
$e^{n\theta/r}$ maps
\[
 H^0(W,n(1-\theta)E_+)
       \lhook\joinrel\longrightarrow H^0(W,n(D+M)).
\]
The exponent is an integer, and the source degrees range through a
Veronese subsequence. On the complement of the zero divisor of $e$,
every ratio of two sections is unchanged. These systems therefore
have the same image dimensions before and after multiplication.
Equation~\eqref{application:perturbed-bound} for the displayed value of $\eta$
proves the asserted lower bound.
\end{proof}

\section{Adjoint positivity of an extreme Hodge line}
\label{period:section}

The relative canonical bundle calculation gives a line in the highest
nonzero Hodge filtration step of one complex summand.  Its curvature
is semipositive, but this alone does not compare it with a canonical
bundle.  We construct a projective quotient on which the line descends
and prove that adding a positive multiple of it makes the canonical
class big.

Here an \emph{extreme Hodge line} means the rank-one highest nonzero
filtration step of a complex Hodge direct summand.  An integral
structure is a lattice in the ambient real local system.  The chosen
summand need not be rational, and the ambient real polarization need
not be rational either.

The \emph{parabolic extension} is obtained by local power substitutions
that make boundary monodromy unipotent, followed by Schmid extension
and descent with rational weights.  The monodromy theorem gives
rational weights in the integral case.  When the local monodromy has
a finite part, these weights are essential.  We first resolve any
compactification whose boundary is not simple normal crossing.
Unipotent extension commutes with pullback, so the same is true of
the parabolic rational line after the power substitutions; see
\cite{Schmid,CKS} and \cite[Lemma~5.6]{BFMT}.  We use version~2
of \cite{BFMT} throughout.


We prove the adjoint-positivity theorem stated as
\Cref{period:adjoint}. The line in that statement is allowed to be a
complex character line inside an integral ambient variation.

There are two maps in the proof.  In a flat trivialization the
\emph{line map} records only the chosen highest line; the \emph{period
map} records all Hodge filtration steps.  Their ranks can differ.
Curvature identifies the numerical dimension of the line with the
rank of the first map.  We use the full periods of suitable rational
adjoint factors to construct the projective quotient.  On that
quotient, logarithmic general type gives a big canonical class with
boundary.  The remaining argument proves that the line loses numerical
rank on each marked boundary component, so a section count removes
the boundary after adding a multiple of the line.

\subsection{Curvature, numerical dimension, and the boundary}

We first relate the curvature on the open set to intersection numbers
on the compactification.  For a nef rational line bundle $A$ on a compact
Kähler $n$-fold, we use
\[
 \nu(A)=\max\{k:c_1(A)^k\cdot[\omega]^{n-k}>0\},
\]
where $\omega$ is any Kähler form.  On a projective manifold this is
the usual numerical dimension.

\begin{lemma}[Curvature and numerical dimension]
\label{period:numerics}
Let $T$ be a smooth compact Kähler manifold, let $E\subset T$ be a
simple normal crossing divisor, and let $A$ be the parabolic extension
of the rank-one highest step of a polarizable pure complex variation on
$T\setminus E$ with quasi-unipotent local monodromy.  Then $A$ is nef.
Its numerical dimension equals the generic rank of the map that
remembers this line in a flat trivialization.

The corresponding assertion for the Griffiths line, the tensor product
of the determinants of the Hodge filtration steps, uses the full period
map.  In particular, its rational extension is big if the full period
map is generically immersive.
\end{lemma}

The curvature-and-volume argument is closely related to the proofs of
\cite[Lemmas~6.15 and~6.17]{BakkerBrunebarbeTsimerman}; here we need
all mixed intersection numbers, since the chosen line need not be big.

\begin{proof}
We may check the assertions after a finite cover and resolution on
which local monodromy is unipotent.  This does not require an integral
lattice.  Indeed, $\pi_1(T\setminus E)$ is finitely generated.  Put
the entries of finitely many monodromy generators and their inverses
in a finitely generated $\Z$-algebra $R\subset\C$.  Adjoin the
finite group of roots of unity generated by the eigenvalues of the
finitely many boundary monodromies, and invert every difference
$\zeta-1$ for a nonidentity element of this group.  Reduction at a
maximal ideal with finite residue field gives a finite-index normal
congruence subgroup.  Every element of a local boundary monodromy
group that lies in this subgroup is unipotent: its eigenvalues
belong to that finite group and reduce to one, so the inverted
differences force them to equal one.  This also applies to products
of commuting boundary monodromies at crossings.  The corresponding finite cover extends
over the compactification, and a Kähler resolution gives the desired
unipotent model; see \cite[Lemma~2.21]{VilladsenHodgeKahler}.
These covers preserve nefness, numerical dimension, and bigness of
the rational line.  Let $h$ be its Hodge metric and $\Theta$ its curvature
on $T\setminus E$, normalized to represent $c_1(A)$.  Griffiths'
curvature formula gives $\Theta\geq0$; its rank is the rank of the
line map, since its value on a tangent vector is the squared norm of
the second fundamental form of the highest step.  The same formula
for the Griffiths line detects the differential of the full period
map.

The Hodge norm estimates in extending frames give, in boundary
coordinates $t_1,\ldots,t_k$, weights whose absolute values are bounded
by
\begin{equation}
\label{period:loglog}
              C\left(1+\sum_{j=1}^k
                   \log(-\log|t_j|)\right).
\end{equation}
The semipositive metric therefore extends as a positive current in
$c_1(A)$ with zero Lelong numbers.  Regularization of positive
currents then gives nefness; see
\cite[Corollary~6.4]{DemaillyRegularization}.
This extends the nefness argument of \cite[Lemma~2.16(1)]{BFMT}
from a proper log smooth algebraic space to the present compact
Kähler setting.

To identify numerical dimension, we must also show that no
intersection mass is lost at $E$.  Write $\alpha=c_1(A)$ and fix
$\varepsilon>0$.  Since $\alpha$ is nef, $\alpha+\varepsilon[\omega]$
has a smooth Kähler representative
$\theta_\varepsilon$.  Write
\[
 \Theta+\varepsilon\omega=\theta_\varepsilon+\ddc u,
 \qquad \ddc=\frac{\sqrt{-1}}{2\pi}\partial\bar\partial.
\]
The global potential $u$ satisfies the local bounds
\eqref{period:loglog}. Put $n=\dim T$. If $E$ is empty, the potential
is smooth and there is no lost mass. Otherwise write
$E=E_1+\cdots+E_N$ and choose smooth divisor metrics, with defining
sections $s_j$, such that
\[
 u_j=-\log|s_j|^2,\qquad
 \ddc(-u_j)=[E_j]-\eta_j,
\]
where the $\eta_j$ are smooth curvature forms. Choose $c>0$ with
$c\eta_j\leq\theta_\varepsilon$ for all $j$, and rescale the metrics
so that $cu_j\geq1$. Then $\phi_j=-cu_j$ is
$\theta_\varepsilon$-plurisubharmonic and $\phi_j\leq-1$.
For any fixed $0<b<1$, the attenuation theorem
\cite[Example~2.14]{GuedjZeriahiEnergy} gives
\[
 \psi_j=-(-\phi_j)^b=-c^bu_j^b
\]
with full Monge--Amp\`ere mass in the K\"ahler class
$[\theta_\varepsilon]$. The full-mass class is convex and contains
the zero potential \cite[Proposition~1.6]{GuedjZeriahiEnergy}.
Thus, for $0<\delta\leq c^b/N$, the potential
\begin{equation}\label{period:barrier}
 v=-\delta\sum_{j=1}^N u_j^b
   =\sum_{j=1}^N\frac{\delta}{c^b}\psi_j
      +\left(1-\frac{N\delta}{c^b}\right)0
\end{equation}
is $\theta_\varepsilon$-plurisubharmonic and has full mass. A positive power grows
faster than a logarithm, so \eqref{period:loglog} implies
$v-C\leq u$ globally for some constant $C$.
Full mass is preserved on passing to less singular potentials
\cite[Propositions~2.11(i) and~2.14(i)]{BEGZ}. Hence $u$ has full mass.
The non-pluripolar measure gives no mass to the analytic set $E$ and
agrees with the smooth wedge product on its complement. Consequently
\[
 \int_{T\setminus E}(\Theta+\varepsilon\omega)^n
                  =(\alpha+\varepsilon[\omega])^n.
\]
Expand both sides as polynomials in $\varepsilon$.  Their coefficients
are finite, since all summands on the left are nonnegative.  Equality
for all $\varepsilon>0$ gives
\begin{equation}
\label{period:masses}
 \int_{T\setminus E}\Theta^k\wedge\omega^{n-k}
                         =\alpha^k\cdot[\omega]^{n-k}
                         \qquad(0\leq k\leq n).
\end{equation}
If the generic curvature rank is $r$, the left side is positive for
$k\leq r$ and zero for $k>r$.  This proves the numerical assertion.
For an immersive Griffiths line the top integral is positive, and
the volume criterion for a positive current on a compact Kähler
manifold gives bigness; see \cite{BoucksomVolume}.
\end{proof}

\begin{lemma}[Restriction to a boundary component]
\label{period:boundary}
In the unipotent case of \Cref{period:numerics}, let $E_i$ be a smooth
boundary component and $N_i$ its monodromy logarithm.  On the open
stratum of $E_i$, the limiting highest line occurs in a unique weight
grade of the limiting mixed variation.  The restriction $A|_{E_i}$
is the Schmid extension, across the remaining boundary of $E_i$, of
that graded Hodge line.  Its numerical dimension is therefore the
generic rank of the graded line map.
\end{lemma}

\begin{proof}
The extended filtration and $W(N_i)$ give the limiting mixed
variation along the open stratum.  Since the highest step has rank
one, exactly one weight grade contains it.  Changes of the normal
parameter act by $\exp(cN_i)$ and act trivially on the associated
graded.  At a further crossing the relative monodromy filtration
identifies the extension of this graded variation with the graded
of the original extending filtration.  Thus the identification holds
on the entire component.
For log smooth algebraic spaces it is stated in
\cite[\S2.5.3 and Lemma~2.16(3)]{BFMT}; its local proof uses the
multivariable nilpotent orbit theorem \cite{CKS} and applies here
as well.  The filtrations and strictness also restrict to a complex
summand.  Apply \Cref{period:numerics} on $E_i$.
\end{proof}

\subsection{Rational adjoint periods and descent of the line}

The period-quotient construction requires discrete monodromy.  A
projection to one complex factor need not retain it.  We therefore
construct a rationally polarized adjoint variation from the ambient
integral local system and keep whole rational factors.  A tensor
construction then recovers a positive power of the original line,
including its scalar monodromy.

\begin{lemma}[Rational adjoint reduction]
\label{period:adjoint-reduction}
In the setting of \Cref{period:adjoint}, let $G$ be the identity
component of the rational Zariski closure of monodromy.  Then $G$ is
semisimple.  Its adjoint local system is a rationally polarized
integral variation of weight zero.

There is a monodromy-invariant collection of rational simple adjoint
factors, each containing a complex factor acting nontrivially on the
highest-line orbit, with the following properties.  The resulting
adjoint variation $\mathbb A$ has discrete integral monodromy.  After a
positive power and up to a constant one-dimensional complex Hodge
shift, the original line is the highest line of a complex
subvariation of a tensor construction on $\mathbb A$.  On a finite
cover preserving the factors, that line splits into the tensor
product of the highest lines belonging to the individual complex
factors.  All these highest steps have rank one.
\end{lemma}

\begin{proof}
Semisimplicity and the theorem of the fixed part hold for polarizable
real and complex variations on a Zariski open subset of a compact
Kähler manifold.  We may use a Kähler modification of $B$ to apply
them; see \cite[Lemma~2.3 and its proof]{PopaSchnellKahler} and
\cite[Theorem~3 and Remark~4]{BakkerMatsushita}.  Thus the algebraic
monodromy group is reductive.  Pass temporarily to a finite cover
on which it is connected.  The irreducible complex constituents of
the local system carry polarizable complex variations.  They can be
defined over a number field: the representation is rational and its
algebraic monodromy is reductive.  The determinant of each constituent
is unitary, as are all its algebraic conjugates, which are again
constituents of the rational representation.  Its monodromy values
are algebraic units by the integral lattice.  Kronecker's theorem
makes these values roots of unity in a fixed number field, hence
makes the determinant character finite.  It is consequently trivial
on $G$.  This is Deligne's finite-determinant argument; its classical
algebraic formulation is \cite[Corollary~4.2.8(iii)(b)]{DeligneHodgeII}.
The connected center of $G$ acts by scalar characters on
the irreducible constituents; their determinants detect these
characters.  Since the original representation is faithful, this
connected center is trivial.  Therefore $G$ is semisimple.

We next equip the monodromy Lie algebra with a rational polarization.
In every tensor construction,
the subspace of $G$-invariant tensors is the fixed part and has a
constant Hodge structure.  In a flat trivialization the Hodge circles
$h_x$ and $h_{x_0}$ therefore induce exactly the same linear operator
on each such subspace at each parameter $z\in S^1$.  Reductivity
identifies $G$ with the pointwise stabilizer of all its invariant
tensors.  The circles preserve these subspaces, so normalize $G$,
and the equality of their actions gives
\[
                    h_x(z)h_{x_0}(z)^{-1}\in G(\R).
\]
Here is a direct local conjugacy argument.  Along a smooth path in
the base, write $a_t=\operatorname{Ad}\circ h_t:S^1\to G^{\rm ad}(\R)$
and $u_t(z)=(\partial_t a_t(z))a_t(z)^{-1}$.  Differentiating the
homomorphism identity gives
\[
 u_t(zz')=u_t(z)+\operatorname{Ad}(a_t(z))u_t(z').
\]
Normalized Haar averaging, with $X_t=\int_{S^1}u_t(z')\,dz'$, yields
$u_t(z)=X_t-\operatorname{Ad}(a_t(z))X_t$.  The solution of
$g_t'g_t^{-1}=X_t$, $g_0=1$, therefore gives
$a_t=g_ta_0g_t^{-1}$.  Lift this path through the central isogeny
$G(\R)^+\to G^{\rm ad}(\R)^+$.  The ratio of $h_t$ and
$g_th_0g_t^{-1}$ lies in $G(\R)$ by the invariant-tensor equality,
and centralizes $G$ by the equality of their adjoint actions.
It lies in the finite center of $G$; continuity in the circle
parameter makes it the identity.  Thus $h_t=g_th_0g_t^{-1}$,
proving that the lifted period map lies in a single
$G(\R)^+$-orbit.

It follows that $\mathfrak g_{\Q}=\operatorname{Lie}(G)$ is a Hodge
sub-local system of $\End(\mathbb V_{\Q})$, with its weight-zero
grading.  The Weil involution $\theta$ preserves $\mathfrak g_{\R}$
and is the restriction of the Cartan involution of the real
polarization isometry group.  To check the restriction, use the
positive Hodge metric: on its Lie algebra $\theta(X)=-X^*$.
If $\mathfrak g_{\R}=\mathfrak k\oplus\mathfrak p$ is its
eigenspace decomposition, then $\mathfrak k\oplus i\mathfrak p$
is a compact real form of $\mathfrak g_{\C}$, represented by
skew-Hermitian endomorphisms.  Thus
$-B_{\mathfrak g}(X,\theta X)$ is positive definite, proving that
$\theta$ is Cartan on $\mathfrak g_{\R}$.  The appropriately
signed Killing form therefore polarizes this weight-zero variation.
The form is rational.  Moreover
\[
          \mathfrak g_{\Q}\cap\End(\mathbb V_{\Z})
\]
is a monodromy-invariant lattice.  This gives the rationally polarized
integral adjoint variation.

Let $\ell$ be the highest line at a reference point of the chosen
complex summand.  The non-lowering parabolic of the Hodge grading
preserves $\ell$, so a Borel subgroup preserves it.  Its closed orbit
\[
                       J=G_{\C}\ell
\]
is a flag variety, a product of flag varieties for the complex simple
factors.  The span $W$ of this orbit is irreducible.  Indeed, in a
semisimple decomposition of the representation, the nonzero
components of a Borel-eigenline all have the same highest weight;
their span is a single diagonal copy of that irreducible
representation.  The line is its unique highest Hodge step.  Indeed, at every point
the Hodge circle normalizes $G$ and preserves the highest line there;
it consequently preserves the span of its $G_{\C}$-orbit.  Thus $W$
is a Hodge subvariation.  The lifted line map takes values in $J$,
and the original line is the pullback of the tautological line there,
equivariantly for monodromy.

Call a complex simple factor \emph{visible} if it acts nontrivially
on $J$.  Retain every rational simple adjoint factor containing a
visible factor.
This collection is invariant under the possibly disconnected
monodromy group.  On its Lie algebra we obtain the variation
$\mathbb A$; its monodromy is contained in the automorphisms of a
lattice and is discrete.  Notice that an entire rational factor is
kept.  Projection to only one real factor would not ensure
discreteness.

We now recover the line from this adjoint system.  Choose a positive
integer $N$ for which the highest weight $N\lambda$ of $W$ belongs
to the root lattice.  The Cartan component
$W_N\subset W^{\otimes N}$ is generated by the orbit of
$\ell^{\otimes N}$.  It factors through the connected adjoint group,
and its highest Hodge line is $\ell^{\otimes N}$.  The same
orbit-span argument just used shows that $W_N$ is a Hodge
subvariation.

We must also account for disconnected monodromy before applying
tensor generation.  Deck transformations take the lifted line to
another point of the same connected orbit, so they preserve $W$ and
its Cartan components.  The kernel of the action on the retained
adjoint factors centralizes the $G$-action on $W$, and hence acts
scalarly by irreducibility.  Its connected part is semisimple and
has no nontrivial scalar character.  Its component group is finite.
Increasing $N$ therefore kills the entire kernel action on $W_N$.
Consequently $W_N$ is a representation of the full projected
algebraic monodromy group, not just of its identity component.
This step rules out a discarded nontorsion scalar local system.

The adjoint representation of that projected group is faithful and
self-dual.  By tensor generation and complete reducibility,
$W_N$ embeds as a flat direct summand in a tensor construction
$\mathbb T(\mathbb A)_{\C}$; see
\cite[Theorem~4.14]{MilneAlgebraicGroups}.  To make the embedding
compatible with Hodge structures, apply the theorem of the fixed part
to
\[
 H^0\bigl(B^\circ,
       \mathcal Hom(W_N,\mathbb T(\mathbb A)_{\C})\bigr).
\]
This nonzero space of flat maps has a constant complex Hodge
structure.  A nonzero Hodge-homogeneous component of any nonzero
flat map is still flat.  Its kernel is $G$-invariant, so the map is
injective because $W_N$ is irreducible.  A constant one-dimensional
complex Hodge shift makes it a morphism of variations; that shift
changes neither the underlying holomorphic line nor its boundary
extension.  This is the constituent and fixed-part argument of
\cite[Theorem~3]{BakkerMatsushita}.

The complex summand carries a flat Hermitian polarization.  In a
real tensor variation of weight $w$ polarized by $Q$, it is induced
by $i^wQ_{\C}(v,\bar u)$, up to the conventional overall sign.
Its restriction to a Hodge subvariation is nondegenerate because
it is definite with the prescribed sign on each Hodge component.
There is no need for the real bilinear form itself to restrict
nondegenerately to the chosen complex summand: it may pair that
summand with its conjugate.

Finally, on a factor-preserving finite cover an irreducible
representation is a tensor product over the simple factors.  The
dimension of its highest step is the product of the corresponding
dimensions.  Since that dimension is one, each factor has a unique
rank-one highest step.  This gives the asserted factor lines.  Their
extensions are nef and functorial by
\Cref{period:numerics,period:boundary}.
\end{proof}

\begin{lemma}[Generic stabilizers of adjoint periods]
\label{period:stabilizer}
Let $G$ be a connected semisimple adjoint group over $\Q$, and let
$\mathbb A_{\Q}$ be a rationally polarized integral variation whose
fiber is $\mathfrak g_{\Q}=\operatorname{Lie}(G)$ and whose connected
algebraic monodromy is $G$, acting by its adjoint representation.
No nonidentity rational Lie algebra automorphism of
$\mathfrak g_{\Q}$ fixes all generic lifted periods of $\mathbb A$.
The automorphism need not belong to the monodromy group.
\end{lemma}

\begin{proof}
Suppose $\gamma$ fixes those periods.  By equivariance, every
monodromy conjugate of $\gamma$ is everywhere of Hodge type $(0,0)$
in the endomorphism variation.  Their rational span is therefore a
rational subvariation of type $(0,0)$.  Its induced polarization is
positive definite.  Intersection with the endomorphism lattice is
a full monodromy-invariant lattice: clear one denominator of
$\gamma$; integral monodromy conjugation preserves that denominator.
A positive definite integral monodromy group is finite.  Consequently
the connected group $G$ fixes $\gamma$, so $\gamma$ centralizes every
inner automorphism of $\mathfrak g$.  The centralizer of the inner
group in $\operatorname{Aut}(\mathfrak g)$ is trivial: on each simple
ideal commutation with all adjoint operators makes an endomorphism
scalar, and compatibility with the Lie bracket makes this scalar
one; permutations of distinct ideals do not commute with their
separate inner actions.  Thus $\gamma=1$.
\end{proof}

\subsection{A projective quotient for the retained periods}

We now quotient by connected fibers of the retained full period map.
The preceding reduction supplies both discrete integral monodromy
and a tensor description of a power of the line.  The first property
allows us to compactify the quotient; the second will descend the
line with its exact parabolic extension.  No bigness assertion about
$K_S+aP$ is used in this construction.

\begin{proposition}[The projective period quotient]
\label{period:quotient}
In \Cref{period:adjoint}, after modifications there is a map
$p:B'\to S$ with connected fibers to a smooth projective variety
and a nef rational line $P$ with $\tau^*L=p^*P$.  On a dense open
subset $S^\circ$, the retained adjoint variation of
\Cref{period:adjoint-reduction} descends, has the same connected
algebraic monodromy, and has a generically immersive period map.
The line $P$ is its parabolically extended highest-weight line,
up to the positive scaling already described.  If the retained
collection is empty, $L$ is rationally trivial and $S$ may be a
point.
\end{proposition}

The construction uses the compact-period-quotient approach of
\cite[Theorems~7.6--7.7]{BakkerBrunebarbeTsimerman}, whose
compactification input goes back to
\cite[Proposition~III and Remark~III-C]{SommesePeriodMapping}.
We use the precise class-$\mathcal C$ formulation of
\cite[Theorem~2.19]{VilladsenHodgeKahler} below.

\begin{proof}
Choose a compact K\"ahler modification of the class-$\mathcal C$
source. Then pass to a finite level at which the projected monodromy is
neat and preserves the factors.  Finite covers of the open locus
extend to finite ramified covers of compactifications and then to
Kähler resolutions; see \cite[Lemma~2.21]{VilladsenHodgeKahler}.
Resolve the boundary to simple normal crossings.  Extend the period
map over all finite-local-monodromy strata.  At neat level their
monodromy is trivial and the nilpotent orbit theorem gives this
extension.

The extended period map on this locus is proper; see
\cite[Proposition~9.11 and the proof of Theorem~9.6]{GriffithsPeriodsIII}.
Indeed, if a sequence tending to the boundary has period images in
a compact subset of the quotient, translate period lifts into a
compact set.
The punctured-disk Schwarz estimate makes the displacement of each
vanishing-coordinate monodromy tend to zero.  The discrete group
acts properly, since its isotropy in the real period domain is
compact.  The local monodromies must therefore be trivial at neat
level, so such a sequence cannot escape the extended locus.

Take the Stein factorization of this proper map.  Its first map has
compact connected fibers and its base is finite over the period
image.  It extends after modifications to a map of compact
class-$\mathcal C$ spaces with Kähler resolutions by
\cite[Theorem~2.19]{VilladsenHodgeKahler} (version~1).
That theorem applies to proper connected-fiber maps from Zariski
open subsets of compact class-$\mathcal C$ spaces.  In terms of
cycles, one first flattens the proper map, sends its irreducible general fibers
to a component of the Douady space of the compact source, and takes
the closure of that family.  The component is compact and in class
$\mathcal C$.  Near a general fiber the map to the cycle space is an
open embedding, since an irreducible compact cycle of the same
dimension in a sufficiently small neighborhood maps to a compact
analytic subset of a small Stein neighborhood in the base, hence
to a point, and must be the whole irreducible fiber.
The incidence map is consequently a modification of the source.
This produces the required compact quotient and also shows its
uniqueness as the quotient by the connected general period fibers.

On a dense open subset of this quotient the period and its integral
variation descend.  At neat level a constant period has trivial
stabilizer; local slices to the connected-fiber map give the descent
data.  After shrinking to a proper smooth fibration, the map on
fundamental groups is surjective, so the descended variation has
the same connected algebraic monodromy.  Its period differential is
generically injective by construction.  On a Kähler compactification
its Griffiths line is big by \Cref{period:numerics}; equivalently,
one can apply \cite[Corollary~1.3]{BrunebarbeCadorel}.  The quotient
therefore has a smooth projective model.  This is also the argument
in \cite[proof of Corollary~2.22]{VilladsenHodgeKahler}.

The finite deck group acts on the connected-period-fiber quotient.
Take its finite quotient and resolve.  This again has a smooth
projective model.  A general fiber before this finite quotient is
the union of connected level fibers indexed by one deck-group orbit.
The group acts transitively on these fibers, so the quotient is the
continuous image of any one of them and is connected.  Thus the
induced map from the original source has connected general fibers,
hence connected fibers after its Stein factorization.  The latter
does not change its function field.
All maps can be resolved using modifications of the original
source and of this projective model.

There is an actual adjoint variation on a dense Zariski open subset
of the quotient before changing level.  Write $\Gamma'$ for the
normal level subgroup of $\Gamma$.  The induced action of
$\Gamma/\Gamma'$ on the level quotient is faithful.  Indeed, if
$\delta\in\Gamma$ acts trivially there, discreteness shows locally
that its action on lifted periods agrees with a fixed
$\eta\in\Gamma'$.  Then $\eta^{-1}\delta$ fixes a period germ,
hence the whole generic lifted period image by analytic
continuation.  By \Cref{period:stabilizer} it is the identity,
so $\delta\in\Gamma'$.  Resolve the finite group action
equivariantly and remove its proper algebraic fixed loci and the
branch locus.  On the resulting dense Zariski open subset the level
quotient is free and unramified.  Its equivariant integral adjoint
local system and Hodge filtration descend there.  In particular,
this descent uses an open subset, not merely very general points.

The highest-weight construction in
\Cref{period:adjoint-reduction} gives a line on that open subset
whose pullback is a fixed positive power of $L$.  Take its
parabolic extension on a log smooth projective compactification
$S$ and divide by that power to obtain $P$.  Nefness follows from
\Cref{period:numerics}.  Pullback of the unipotent extensions,
followed by rational descent, gives
$\tau^*L=p^*P$ on the compact models, including the boundary.
Equivalently, this can be checked at level and then descended by
the norm: a line whose finite pullback is trivial is torsion in
the rational Picard group.  The identity therefore includes the
boundary and persists under further modifications.

If no factor is retained, the orbit is a point and the only possible
scalar monodromy is the finite scalar group treated in
\Cref{period:adjoint-reduction}.  Its parabolic line is rationally
trivial, as asserted.
\end{proof}

\subsection{Loss of numerical rank on the marked boundary}

Fix the projective quotient of \Cref{period:quotient} and a log smooth
compactification $(S,E)$ of its descended variation.  Let $D\leq E$
be the reduced sum of components with nonidentity local monodromy
on the retained adjoint system, including finite monodromy.
Logarithmic general type will give bigness of $K_S+D$.  The next
lemma is what permits us to remove $D$.

\begin{lemma}[Strict loss of rank at the marked boundary]
\label{period:rankloss}
For every component $D_i$ of $D$,
\begin{equation}
\label{period:strictdrop}
                       \nu(P|_{D_i})<\nu(P).
\end{equation}
\end{lemma}

\begin{proof}
Write $n=\dim S$ and $r=\nu(P)$, the generic rank of the line map
$j$ by \Cref{period:numerics}. The full retained period map has rank
$n$; the argument does not require $r=n$.
Work first on a finite Galois cover of the
open locus with connected algebraic monodromy and no factor
permutations.  Its connected adjoint group is denoted by $G$.

\smallskip\noindent
\emph{The Ax--Schanuel decomposition.}
Choose a Hodge-generic point, also general for the ranks of the
factor period maps and their joint differentials.  There are only
finitely many rational normal-factor decompositions of $G$, so
these requirements can be imposed simultaneously.  Let $T$ be a
local fiber germ of $j$ there, and $Z$ its irreducible Zariski
closure in the source of the variation.  Thus $\dim T=n-r$.
Move within the germ to a smooth point of $Z$ and use a smooth
open subset of a resolution when applying Hodge theory.

The restricted variation on $Z$ has the same generic Mumford--Tate
group as the ambient variation, since $Z$ contains the chosen
Hodge-generic point.  The monodromy normality theorem therefore
makes its connected monodromy $H$ normal in $G$; see
\cite[\S5, Theorem~1]{AndreMonodromy}.  In adjoint notation write
\[
                            G=H\times Q.
\]
The $Q$-period is constant on the lifted $Z$, by the fixed-part
theorem after killing finite monodromy.  The varying monodromy
domain on $Z$ is therefore the domain for $H$.  Write $J_H$ for
its line-orbit projection and $d_H=\dim J_H$.
In local flat coordinates write $j=(j_H,j_Q)$, with targets $J_H$
and $J_Q$, the products of line orbits for the two groups.  We will
show that the $H$-line map has full rank $d_H$ on $Z$, whereas the
$Q$-line map and the full $Q$-period map have equal generic rank
on the ambient base.

For the first assertion, put
$m=\dim Z$ and $d=\dim\check{\mathcal D}_H$, where
$\check{\mathcal D}_H$ is the compact dual of the connected-monodromy
period domain.  The projection
$\check{\mathcal D}_H\to J_H$ is a surjective morphism of
homogeneous projective flag varieties.  All its fibers have
dimension $d-d_H$.  Let
$V\subset Z\times\check{\mathcal D}_H$ be the algebraic locus
where the projected line equals its value on $T$.  It follows that
$\dim V=m+d-d_H$.  The period graph has dimension $m$ in an
ambient variety of dimension $m+d$, so its expected intersection
dimension with $V$ is $m-d_H$.

By \cite[Theorem~1.1]{BakkerTsimermanAxSchanuel}, an intersection
component of larger dimension projects into a proper algebraic weak
Mumford--Tate subvariety of $Z$.  This is impossible for the component
containing the lifted period graph over $T$, since $T$ is Zariski
dense in $Z$.  Conversely, a map
to $J_H$ has fibers of dimension at least $m-d_H$.  At the chosen
smooth point the restricted fiber is exactly $T$, because the
ambient local fiber is $T\subset Z$.  Equality follows, proving the
claim.  The cited theorem applies to the rationally polarized
integral adjoint variation on a smooth algebraic base and uses
precisely its connected-monodromy domain.

Let $q$ denote the full $Q$-period map and set
$s_Q=\rank(dq)$. Choose the smooth point $x\in T\cap Z_{\mathrm{reg}}$
inside the ambient constant-rank neighborhood fixed above. Since
$T$ is Zariski dense in $Z$, such a point exists. The constant-rank
theorem gives $\ker(dj_x)=T_xT\subset T_xZ$, and hence
\[
 \rank(dj_x|_{T_xZ})=\dim Z-\dim T=d_H.
\]
The $Q$-period is constant on $Z$, so $T_xZ\subseteq\ker(dq_x)$.
Thus, at this same point,
\[
 \ker(dj)\subseteq\ker(dq),\qquad
             dj_H(\ker(dq_x))=T_{j_H(x)}J_H.
\]
The first inclusion uses $T_xT\subset T_xZ$; the second is the
just-established rank calculation on $T_xZ$.
The differential $dj_Q$ factors through $dq$.  If $dj_Q(v)=0$,
choose $w\in\ker(dq)$ with $dj_H(w)=dj_H(v)$.  Then
$v-w\in\ker(dj)$, and therefore $dq(v)=0$.  We conclude
\begin{equation}
\label{period:rankdecomposition}
       \rank(dj_Q)=s_Q,\qquad r=d_H+s_Q.
\end{equation}
These are generic identities by our choice of the point.

The factor-line construction, after a common positive rescaling
and a further finite cover, now gives nef lines with
\begin{equation}
\label{period:linedecomposition}
       \pi^*P=P_H+P_Q,\qquad
       \nu(P_H)=d_H,
       \qquad P_Q=q^*B_Q.
\end{equation}
Here $q:\widetilde S\to S_Q$ is a resolved projective quotient
for the full $Q$-period, constructed by
\Cref{period:quotient}, and $\dim S_Q=s_Q$.  The highest-line
rank on $S_Q$ is $s_Q$ by
\eqref{period:rankdecomposition}.  Only the quotient construction,
not adjoint positivity, has been used to construct $S_Q$.

\smallskip\noindent
\emph{Reduction for a boundary component.}
For a fixed marked $D_i$, use a common finite Galois cover that
preserves the factors and makes boundary monodromy unipotent.
Resolve it and the map $q$ equivariantly, and let $E_i$ be the
strict transform of a component above $D_i$ that is generically
finite over $D_i$.  The inertia group at its generic point fixes
$E_i$ pointwise.  Further rescaling does not change any numerical
dimension.  Nefness gives
\[
                     \nu(P_H|_{E_i})\leq d_H.
\]
For example, every intersection involving $E_i$ is bounded by
the corresponding intersection with a sufficiently ample
effective divisor; the vanishing of higher powers of $P_H$
then gives this inequality.  If $E_i$ does not dominate $S_Q$,
\[
                     \nu(P_Q|_{E_i})\leq s_Q-1,
\]
and the binomial expansion for the two nef lines in
\eqref{period:linedecomposition} proves the strict drop.
We may therefore suppose $E_i$ dominates $S_Q$.  At its generic
point the $Q$-variation is pulled back from the interior of
$S_Q$, so its local unipotent logarithm vanishes.

\smallskip\noindent
\emph{Infinite local monodromy.}
If the original local monodromy is infinite, its logarithm at
unipotent level is a nonzero rational element
$N\in\mathfrak h$.  Every nonzero rational element of a rational
simple factor has nonzero projection to every conjugate complex
factor.  Each retained rational factor has a visible complex
factor.  Hence $N$ acts nontrivially on $J_H$.

Let $w$ be the pure weight of the representation supplying $P_H$,
and let $F^a=\C v$ be its limiting highest line at a general
point of $E_i$.  This line has a unique weight contribution,
say of weight $w+k$.  It is primitive and $k\geq0$: a
nonprimitive contribution, or a contribution below the middle
weight, would come by monodromy Lefschetz from a nonzero
filtration piece above $F^a$.
Moreover
\begin{equation}
\label{period:nilpotenttop}
                    N^{k+1}v=0,\qquad N^kv\ne0.
\end{equation}
To see that this statement holds for the actual vector rather
than only its weight-graded image, use the Deligne splitting
of the limiting mixed Hodge structure.  The one-dimensional
highest step occupies a single summand, and $N$ has degree
$(-1,-1)$ for this functorial splitting.  Primitive vanishing
on the graded piece therefore gives the first equality in
\eqref{period:nilpotenttop}; the Lefschetz isomorphism gives
the second.  For a complex summand the same argument is made
in its conjugate-paired real variation and restricted to the
summand.

In nilpotent-orbit coordinates $[v]\in J_H$.  Since this orbit
is closed in projective space and invariant under $\exp(tN)$,
\begin{equation}
\label{period:fixedlocus}
       [N^kv]=\lim_{t\to\infty}\exp(tN)[v]
                \in J_H\cap\mathbb P(\ker N).
\end{equation}
This fixed locus is a proper closed subset of the irreducible
variety $J_H$, because $N$ acts nontrivially on it.  Its dimension
is less than $d_H$.  Restrict to a simply connected open part of the boundary stratum
on which $k$ is constant.  In a flat trivialization, both $N$ and
$W(N)$ are constant.  Project $N^kv$ to
$\Gr^{W(N)}_{w-k}$ and apply the inverse of the fixed Lefschetz
isomorphism
\[
 N^k:\Gr^{W(N)}_{w+k}\longrightarrow\Gr^{W(N)}_{w-k}.
\]
This recovers the weight-graded limiting highest line.  Every actual
vector $N^kv$ lies in $W(N)_{w-k}$, and its image in that graded
piece is nonzero.  Thus the graded line map factors through the
rational map on
\[
 J_H\cap\mathbb P(\ker N)\cap\mathbb P(W(N)_{w-k})
\]
induced by this fixed linear quotient and inverse.  This locus is
contained in the proper fixed locus in \eqref{period:fixedlocus},
and the rational map is defined at every point under consideration.
Its rank is less than $d_H$.
By \Cref{period:boundary},
\[
                       \nu(P_H|_{E_i})<d_H.
\]
Together with $\nu(P_Q|_{E_i})\leq s_Q$ this proves the
required inequality in the infinite-monodromy case.

\smallskip\noindent
\emph{Finite nonidentity local monodromy.}
Let $\gamma$ be the original finite local monodromy.  The
disconnected-monodromy argument in \Cref{period:adjoint-reduction}
shows that $\gamma$ preserves the orbit span $W$.  It therefore acts
on the highest-line orbit $J$, permuting its nontrivial simple-factor
flag varieties. Equivariance preserves the kernel of
$G_{\C}\to\operatorname{Aut}(J)$, which is the product of the
invisible simple factors. Thus visible and invisible factors cannot
be exchanged. At unipotent level its logarithm is zero, so the variation extends
with pure limiting flags.  The inertia generator fixes the generic
boundary point; equivariance therefore makes these flags fixed by
$\gamma$.  Suppose that the
boundary line retained rank $r=d_H+s_Q$.  Since $E_i$ dominates $S_Q$,
its full $Q$-period and its $Q$-line map both have rank $s_Q$.
Along a generic fiber of this $Q$-period, the $H$-line map
must consequently have rank $d_H$.  Its coordinates fill an
open subset of $J_H$ while the full $Q$-flags remain fixed.

The fixed-point condition for $\gamma$ rules out permutations
involving any visible $H$-factor, including a permutation
with a factor whose coordinate is held fixed.  On each visible
$H$-factor of the highest-line orbit it forces the identity action.
The action of a simple adjoint group on a nontrivial flag variety is faithful;
equivariance then makes the induced automorphism of that
simple group the identity as well.  Rationality forces the
same conclusion on all conjugate factors.  Thus $\gamma$
acts identically on $H$ and preserves $Q$.

Its $Q$-action fixes every generic $Q$-period, by the domination
of $S_Q$ and analytic continuation from the resulting open set of
lifted periods. Apply \Cref{period:stabilizer} to the descended
$Q$-variation.  This action is the identity.  Hence $\gamma$
is the identity on the whole retained adjoint system,
contradicting the marking.

In both cases nef binomial expansion and
\eqref{period:rankdecomposition} give
$\nu(\pi^*P|_{E_i})<r$.  Numerical dimension of a nef line is
unchanged by a generically finite pullback, so this is
\eqref{period:strictdrop}.  If $r=0$, the local line fiber is the whole source germ, so its
Zariski closure is the whole source and $H=G$, $Q=1$.  The
Ax--Schanuel calculation gives $d_H=0$.  Every retained rational
factor contains a factor acting nontrivially on the line orbit;
hence there can be no retained factor.  Generic immersivity of the
full retained period map then forces $S$ to be a point, with no
marked component.
\end{proof}

\subsection{Removing a boundary by counting sections}

The period argument must pass from a big logarithmic canonical
divisor to a big adjoint divisor without its boundary. The following
lemma isolates the last step: sections lost on the boundary grow
more slowly in the nef direction than sections on the whole space.
For a nef rational divisor $J$ on a smooth projective $n$-fold, set
\[
 \nu(J)=\max\{q\in\{0,\ldots,n\}:J^q\cdot H^{n-q}>0\},
\]
where $H$ is ample. This is independent of $H$ and equals zero
on a point. The proof below fixes the nef coefficient before
taking the asymptotic section limit.

\begin{lemma}\label{period:subtract}
Let $Z$ be a smooth projective variety, $J$ a nef rational divisor,
$A$ a rational divisor, and
$D=\sum_{i=1}^s d_iE_i$ an effective rational divisor with smooth
prime components. Suppose $A+D$ is big and
\[
 \nu(J|_{E_i})<r:=\nu(J)
 \quad\text{for every component of }D.
\]
Then $A+tJ$ is big for every sufficiently large rational $t$.
\end{lemma}

\begin{proof}
If $r=0$, the condition forces $D=0$, so $A$ is already big;
adding a nef divisor preserves bigness. This also covers dimension
zero with the usual convention. Assume $r>0$ and put $n=\dim Z$.
Choose an ample rational divisor $H_1$ with
$A+D-H_1$ rationally effective. For fixed rational $t\geq0$
and sufficiently divisible $m$,
\[
 h^0(Z,m(A+D+tJ))\geq h^0(Z,m(H_1+tJ)).
\]
Since $H_1+tJ$ is ample, asymptotic Riemann--Roch makes the
leading coefficient on the right
$(H_1+tJ)^n/n!$. As a polynomial in $t$, this intersection has
degree $r$ and positive leading coefficient.

Remove $mD$ one component at a time by divisor exact sequences.
Every lost quotient is a line bundle on an $E_i$ of the form
\[
 \mathcal O_{E_i}\left(m(A+D+tJ)-\sum_h k_hE_h\right),
 \qquad 0\leq k_h\leq md_h.
\]
Choose an ample integral divisor $B_0$ so large that a fixed
denominator-clearing multiple of $B_0-A-D$ is globally generated.
For each $h$, choose an ample integral divisor $B_h$ such that
both $B_h$ and $B_h+E_h$ are globally generated. Put
$H_2=B_0+\sum_h d_hB_h$. The difference between
$m(H_2+tJ)$ and the displayed quotient divisor is
\[
 m(B_0-A-D)+
 \sum_h\bigl((md_h-k_h)B_h+k_h(B_h+E_h)\bigr).
\]
It is globally generated in these divisible degrees. A section
not vanishing identically on $E_i$ gives an injection into
$\mathcal O_{E_i}(m(H_2+tJ)|_{E_i})$. There are exactly
$md_i$ quotient terms on $E_i$. Hence
\[
 h^0(Z,m(A+tJ))
 \geq h^0(Z,m(H_1+tJ))
       -m\sum_i d_i h^0(E_i,m(H_2+tJ)|_{E_i}).
\]
For each fixed $t$, divide by $m^n/n!$ and let $m$ tend to
infinity through the specified multiples. All the comparison
divisors are ample, so the resulting lower bound is
\[
 (H_1+tJ)^n
     -n\sum_i d_i(H_2+tJ)^{n-1}\cdot E_i.
\]
Each restriction polynomial has degree at most
$\nu(J|_{E_i})\leq r-1$, because the restricted divisor is nef.
The positive degree-$r$ term of the first polynomial therefore
dominates the loss for every sufficiently large $t$. The section
growth of $A+tJ$ is then positive in dimension $n$, which is
exactly bigness.
\end{proof}

Taking $A=K_Z$ and $D$ reduced gives the required removal of the
period boundary once its restriction ranks have been established.
The weights also allow simultaneous removal of other effective
divisors satisfying the same rank inequality, for example a
divisorial ramification term on a generically finite cover.


\begin{proof}[Proof of \Cref{period:adjoint}]
Apply \Cref{period:quotient}.  If its target is a point there
is nothing further to prove.  Otherwise mark the divisor $D$
as above.  The adjoint variation extends across every unmarked
boundary component, including their intersections away from
$D$, because its local monodromies there are trivial.  Its
period differential is generically injective.  Consequently
\cite[Theorem~1.1]{BrunebarbeCadorel} gives
\[
                              K_S+D\text{ big}.
\]
If an initial choice omitted codimension-two loci, apply the
theorem on a log resolution and push the sections down; the
resulting logarithmic canonical divisor pushes forward to
$K_S+D$.  By \Cref{period:rankloss}, each marked component
has strictly smaller numerical rank.  Apply \Cref{period:subtract} with $A=K_S$ and the marked reduced
divisor $D$. It gives $K_S+aP$ big.  Together with the
rational line identity in \Cref{period:quotient}, this proves
\eqref{period:conclusion}.  Since $P$ is nef, the same bigness
holds for every larger rational coefficient $a$.
\end{proof}

\section{Addition with a log-general-type fiber}\label{generaltype:section}

We prove subadditivity when the fiber is of log general type.  The base
may be nonprojective: the projective positivity theorem will be applied
to a stable family over the image of its classifying map.

We prove \Cref{generaltype:addition}.

The proof constructs finitely many relative pluricanonical forms whose
ratios give a generically finite map on a general fiber.  Their poles
must be controlled at every vertical divisor, including divisors whose
image on the original base has codimension at least two.  Multiplying
these forms by base pluricanonical sections then retains both the fiber
coordinates and the base Iitaka coordinates.

We first construct a projective stable family over the image of a
classifying map.  Positivity is applied to that family.  Its sections
are pulled back, descended through a finite cover, and tested at
vertical valuations.  All degrees below clear the denominators of the
relevant divisors.

\subsection{Comparison with a stable family}

We use the stable-family convention of
\cite[Notation~3.7, Remark~3.8, and Definitions~3.9 and~3.11]
{KovacsPatakfalvi}. A stable pair has semi-log-canonical singularities
and an ample rational log canonical divisor. A stable family over a
normal base is flat, proper and surjective with connected stable
fibers; its boundary support avoids the generic points and the
codimension-one singular points of every fiber. Its relative log
canonical divisor is rational Cartier and relatively ample, so the
family is projective. Boundaries are pulled
back by divisorial restriction, and Cartier multiples of the relative
log canonical divisor commute with base change between normal bases.
Semi-log-canonicity
allows nonnormal fibers and tests log canonicity on the normalization
with the conductor boundary. The general fibers in our application
are additionally Kawamata log terminal (klt): they are normal and
all their log discrepancies are positive. This extra condition will
be checked before applying positivity.
A stable family has \emph{maximal variation} when its classifying map
to the moduli space of stable pairs is generically finite onto its image.

Fiberwise finite generation gives an embedding degree separately on
each fiber.  The next lemma selects one degree on an analytic base.
As throughout, \emph{very general} excludes a countable union of proper
closed analytic subsets.

\begin{lemma}\label{generaltype:selection}
Let \(Y\) be an irreducible compact complex space.  For each positive integer
\(j\), suppose that \(q_j\colon Q_j\to Y\) is a proper surjective map from an
irreducible compact complex space and that, over a dense analytic open
\(U_j\subset Y\), its fibers are irreducible of a fixed dimension \(r_j\).
Let
\[
  \phi_j\colon Q_j\dashrightarrow Z_j
\]
be a meromorphic map with irreducible projective image \(Z_j\).
Suppose that a dense open part of \(q_j^{-1}(U_j)\), called the genuine
frame locus, meets every fiber densely and that \(\phi_j\) is holomorphic
there.  Let \(G_j\subset Z_j\) be a countable union of algebraically
constructible subsets.

Let \(\Lambda_j\subset U_j\) be sets whose union contains a very general
subset of \(Y\), such that for \(y\in\Lambda_j\) all genuine frames above
\(y\) map into \(G_j\).  Then some \(\Lambda_j\) is not contained in a
countable union of proper closed analytic subsets of \(Y\), and for
that index one of the constructible subsets making up \(G_j\) is dense
in \(Z_j\).  In particular \(G_j\) contains a nonempty algebraic open
subset of \(Z_j\).
\end{lemma}

\begin{proof}
Some \(\Lambda_j\) is not contained in a countable union of proper
closed analytic subsets; otherwise their union would be so contained.
Resolve the graph of \(\phi_j\) and shrink \(U_j\), if necessary.
This discards only a proper analytic subset of \(Y\) and preserves
the stated properties of the general fibers.
Write \(G_j=\bigcup_iG_{ji}\).  If no \(G_{ji}\) is dense, the inverse
images
\[
  E_{ji}=\phi_j^{-1}(\overline{G_{ji}})
\]
are proper closed analytic subsets of the resolved \(Q_j\).
For \(y\in\Lambda_j\), the genuine frame fiber is covered by these
subsets.  Baire's theorem implies that one \(E_{ji}\) contains a
relatively open part of the fiber.  Irreducibility then implies
that \(E_{ji}\) contains the whole compact fiber.

The locus
\[
  A_{ji}=\{y\in U_j:
            \dim(E_{ji}\cap q_j^{-1}(y))\geq r_j\}
\]
is analytic, by the fiber-dimension theorem for proper maps.  It is
proper: otherwise \(E_{ji}\) would contain the general fibers and
hence all of \(Q_j\).  The same theorem before restriction to \(U_j\)
shows that \(A_{ji}\), together with \(Y\setminus U_j\), is contained
in a proper closed analytic subset of \(Y\).
This would place \(\Lambda_j\) in countably many proper closed
analytic subsets, a contradiction.  Finally, a dense constructible
subset of an irreducible algebraic variety contains a nonempty
open subset.
\end{proof}
\begin{lemma}[Analytic comparison with a stable family]
\label{generaltype:stable-comparison}
Let \(h\colon V\to Y\) be a surjective holomorphic map with connected
fibers, where \(V\) is a smooth compact Kähler manifold and \(Y\) is
a smooth compact manifold in class \(\mathcal C\).  Let \(A\) be an
effective rational SNC divisor on \(V\), horizontal over \(Y\), with
coefficients strictly less than one.  Suppose \(K_F+A_F\) is big on
very general fibers \(F\) of positive dimension.  There are
\begin{enumerate}[label=\textup{(\roman*)}]
\item a proper generically finite map \(\tau\colon Y'\to Y\), with \(Y'\)
smooth, compact, and in class \(\mathcal C\);
\item a surjective holomorphic map \(\rho\colon Y'\to R\), with \(R\)
smooth and projective;
\item a projective stable family
\(\pi\colon(\mathcal V,\mathcal A)\to R\) of maximal variation, with klt
general fiber;
\end{enumerate}
such that the main component of \(V\times_Y Y'\) is bimeromorphic
over \(Y'\) to \(\mathcal V\times_RY'\).  Over a dense analytic open,
this meromorphic comparison is the fiberwise log canonical model map
for \(K_F+A_F\), with the horizontal pushforward of \(A\) as boundary.
The space \(R\) may be a point.
\end{lemma}

\begin{proof}
The construction has three outputs.  Relative pluricanonical images
produce algebraic parameter data even though the base is analytic.
A classifying map then supplies a projective stable family after a
finite base cover.  Finally, an isomorphism cover matches that family
with the relative image; fiberwise canonical models determine the
horizontal discrepancy signs.

\emph{Relative images and frames.}
Set \(\mathcal L=\mathcal O_V(K_{V/Y}+A)\), understood as a rational
line bundle.  For the factorial multiples \(m\) clearing its denominator,
the sheaves
\[
  \mathcal E_m=h_*\mathcal L^{\otimes m}
\]
are coherent. First restrict to a dense analytic open where \(h\) and
the horizontal boundary strata are smooth, so the proper family and
its line bundles are flat over the base. Generic constancy of the
fiber section dimension and cohomology and base change then identify
the fibers of \(\mathcal E_m\) with the complete systems, separately
for each \(m\); see \cite[\S7, nos.~4--6, S\"atze~3--5]{GrauertDirectImages}
and its corrigendum \cite{GrauertDirectImagesCorr}.
After a modification of \(Y\),
the strict transform of \(\mathcal E_m\), modulo torsion, is locally free.
The evaluation map gives a relative meromorphic map into its projective
bundle.  Resolve this map and take its proper analytic image
\(\mathcal Z_m\).  We also take the proper images of the horizontal
boundary components and retain their divisorial cycle parts.
Analytic flattening makes these embedded image and cycle data flat on a
dense analytic open, with fixed Hilbert and cycle polynomials.
These uses of coherent direct images and analytic flattening are
independent of any minimal-model theorem; see \cite{HironakaFlattening}.

Write \(s_m=\rank\mathcal E_m\) on this open; indices with \(s_m=0\)
are omitted.  Its projective frame bundle has a natural compactification
\[
 \overline Q_m=
 \mathbb P\!\left(\Hom(\mathcal O_Y^{\,s_m},\mathcal E_m)\right).
\]
Here \(Y\) denotes the modification chosen for this degree; its map to
the original base remains proper and bimeromorphic.  The locus of
invertible matrices modulo scalars is a principal
\(\operatorname{PGL}(s_m)\)-bundle.
A genuine frame embeds the corresponding fiber of \(\mathcal Z_m\)
in a fixed \(\mathbb P^{s_m-1}\).  There is therefore a holomorphic
Hilbert--Chow parameter map on the genuine frame locus.
The Hilbert spaces here are the algebraic Hilbert schemes of
projective space, viewed analytically; equivalently one may use
their Douady realization \cite{Douady}.

To extend this map meromorphically, use the existing compact relative
image and the projective linear action in local trivializations.
The graph of that action extends over the projective matrix
space: inversion is given by the adjugate matrix.  Apply this construction
to the proper family \(\mathcal Z_m\) and its marked cycles, take proper
analytic images, and flatten by strict transform.  One obtains a proper flat family
of embedded subspaces over a modification of \(\overline Q_m\).
Its Hilbert--Chow map is holomorphic.  The local graph constructions
agree on the genuine frames and hence glue.  Thus the original
parameter map is meromorphic.  After resolving it, its image is
compact analytic in a fixed projective Hilbert--Chow space, and is
consequently algebraic.  Denote the irreducible image by \(Z_m\).
This construction uses a frame bundle, not a meromorphic
trivialization of \(\mathcal E_m\).

\emph{The good parameter locus.}
In \(Z_m\), call an embedded pair \((Z,D)\) good if it is klt and stable,
has the prescribed boundary weights, has a complete nondegenerate
embedding, and satisfies
\begin{equation}\label{generaltype:embedded-adjoint}
  \mathcal O_Z(1)\simeq
       \mathcal O_Z\bigl(m(K_Z+D)\bigr).
\end{equation}
One may also impose the generic Hilbert and marking data.
This is a countable union of algebraically constructible loci.
Here is the parameter-space justification.  Stratify the universal
embedded families by their Hilbert polynomials, Cartier indices,
and base-change loci.  On each such stratum,
\Cref{generaltype:embedded-adjoint} can be imposed by parameterizing
the map
\[
  \omega_Z^{[m]}\longrightarrow\mathcal O_Z(1)
\]
with its prescribed divisor.  These parameter spaces are of finite
type.  The stable-pair and klt conditions are constructible there,
and their images are constructible by Chevalley's theorem.
This is the embedded construction used for stable log-variety
moduli in \cite[Section~6.A, especially Propositions~6.8 and~6.11
and Lemma~6.12]{KovacsPatakfalvi}.  In particular no assertion
about analytic minimal models is hidden in the good locus.

Every very general fiber is projective: bigness makes it Moishezon,
and it is Kähler. We use Moishezon's projectivity criterion, which
is also the smooth case of \cite[Theorem~6]{NamikawaProjectivity}.
The projective klt finite-generation theorem
\cite[Theorem~1.1 and its consequences]{BCHM} therefore gives its log
canonical model.  Every sufficiently large divisible pluricanonical
degree embeds that model and satisfies
\Cref{generaltype:embedded-adjoint}.  Exclude at once the countably
many generic base-change loci, relative-rank loci, and generic-degree
loci for the maps constructed above.  Let \(\Lambda_m\) consist of
the remaining parameters for which the degree-\(m\) map actually
constructs the fiber's log canonical model.
Their union contains a very general subset of \(Y\), and all genuine
frames over a point of \(\Lambda_m\) give good data.
Apply \Cref{generaltype:selection}.  We obtain one fixed \(m\) for which
\(\Lambda_m\) is not contained in a countable union of proper analytic
subsets and good data hold on an algebraic open of \(Z_m\).

For this degree the relative image map is generically birational.
Indeed its generic finite degree agrees with its fiberwise degree
outside the generic-degree exceptional set already excluded.
Choose a point of \(\Lambda_m\) there; that degree is one.
Retain \(\Lambda_m\): its actual canonical models will determine the
horizontal discrepancy signs at the end of the construction.

\emph{The invariant classifying map.}
Fix the dimension, the volume read from
\Cref{generaltype:embedded-adjoint}, and a finite coefficient set
containing the coefficients of \(A\) and closed under sums at most one.
The good algebraic parameter locus has a classifying map to the
projective coarse space \(M\) of stable pairs with these data.
If an additional parameterization of the map in
\Cref{generaltype:embedded-adjoint} was used, its moduli graph descends:
the embedded pair determines a unique isomorphism class, so the
closure of that graph is generically a singleton over the good
parameter locus.  Restricting a dense algebraic open therefore gives
the same rational classifying map on \(Z_m\).
It is invariant under change of projective frame.

Compose with the meromorphic frame parameter map.
The resulting map \(\overline Q_m\dashrightarrow M\) is constant
on general frame fibers.  Its graph has a proper analytic image
in \(Y\times M\).  This image is generically a singleton over \(Y\),
so its projection to \(Y\) is bimeromorphic.  It is the graph of
a meromorphic map
\[
  \mu\colon Y\dashrightarrow M.
\]
Let \(M_0\) be the closure of its image; it is projective.

\emph{A stable parameter cover and the matching cover.}
Stable moduli supplies a finite parameter cover
\(S\to M\) from a normal scheme carrying a stable family
\cite[Corollary~6.19]{KovacsPatakfalvi}.
The coarse space \(M\) is projective
\cite[Corollary~7.3]{KovacsPatakfalvi}, so its finite cover \(S\)
is projective as well. The moduli construction for the specified
coefficient set is given in
\cite[Definition~6.2, Proposition~6.4, and
Proposition~6.14]{KovacsPatakfalvi}.
Resolve \(\mu\), and normalize a component of the pullback of
\(S\to M\) dominating \(Y\).  Let \(R_0\) be its image in \(S\).
The restricted stable family over \(R_0\) has maximal variation,
because \(R_0\to M_0\) is finite.

At this point the projective family exists over \(R_0\).
Equality of coarse moduli points identifies its general fibers with
those of the original analytic image, but does not identify the
families. We construct a further generically finite cover carrying
such an identification.

On a dense analytic open of the comparison base, embed the original
flat image by the complete system of
\(\mathcal O_{\mathcal Z_m}(k)\), and the pulled-back stable family
by its \(km\)-adjoint system, for a common sufficiently large and
divisible \(k\). The two sheaves of sections are coherent direct
images of compact families; generic cohomology and base change make
them locally free after shrinking. On a common compact modification
of the base, make their torsion-free strict transforms locally free,
resolve the evaluation maps, and flatten the embedded image families
and their weighted boundary cycles. This gives holomorphic relative
Hilbert--Chow data. The modification concerns the base itself, not a
space of choices of frames.

On every general good fiber, \Cref{generaltype:embedded-adjoint}
identifies the two polarizations. Thus the projective linear maps
identifying the two embedded varieties and their weighted boundary
cycles form a nonempty Isom set. This set is finite, since stable
pairs have finite automorphism groups
\cite[Propositions~6.5 and~6.8]{KovacsPatakfalvi}; completeness and
nondegeneracy of the embeddings exclude additional ambient
stabilizers. At this stage we use only the embedded analytic image
and its polarization. Its \(\Q\)-Gorenstein family structure will
be transported from the stable family by the identification.

To compactify these Isom sets, work in local trivializations of the
two section bundles. In the product of their projective parameter
spaces with the projective matrix space, take the algebraic closure
of the incidence relation for invertible matrices identifying the
embedded varieties and weighted boundary cycles, after clearing the
fixed coefficient denominators. Pull back this closure by the two
Hilbert--Chow maps. Equivariance under changes of both frames glues
the local incidence spaces into a proper analytic space in the
projective bundle of homomorphisms between the section bundles; no
meromorphic trivialization is used. Retain its compact irreducible
components that dominate the base and meet the invertible locus,
and normalize them. The invertible locus has the finite Isom fibers
just described, so the resulting cover is proper and generically
finite. Its tautological isomorphism identifies the generic families,
including the stable family's \(\Q\)-Gorenstein structure; its graph
therefore extends meromorphically over the compact base.

Finally resolve \(R_0\) projectively and resolve the graph of the
map from the comparison cover to it.  Denote the resulting spaces
by \(R\) and \(Y'\).  Stable families and their divisible relative
adjoint line bundles commute with base change, with boundary understood
by divisorial restriction \cite[Remark~3.8]{KovacsPatakfalvi}.
Thus the pullback family over \(R\) remains stable; it has klt general
fibers and maximal variation.  All covers and modifications used above stay
in class \(\mathcal C\).

On a common resolution we can now compare the source adjoint
divisor with the pullback of the stable relative adjoint divisor,
which is rational Cartier.  There are only finitely many horizontal
prime coefficients in their difference.  On fibers above
\(\Lambda_m\), away from the finitely many generic comparison
exceptions, they are the log canonical model discrepancy coefficients.
They are therefore nonnegative and supported on divisors exceptional
over the stable model.  The same statements hold horizontally in
the family.  Remove the images of the remaining vertical discrepancy
divisors.  On the resulting dense analytic open this is the log
canonical model comparison, with its divisorial pushforward boundary.
This proves the asserted meromorphic comparison without relative
analytic finite generation.  If \(M_0\) is a point, the construction
uses the fixed stable pair and \(R\) is a point.
\end{proof}

\subsection{Relative forms and vertical valuations}

For a meromorphic section \(\sigma\) of
\(\mathcal O_V(bK_{V/Y})\), its horizontal pole bound \(bA\)
means that \(\ord_E(\sigma)\geq-bA_E\) at every horizontal prime,
using a local nonvanishing base volume to interpret relative forms.
At a vertical prime over a divisor \(D\subset Y\), write
\(\omega_D\) for a local nonvanishing logarithmic top form on the
base with a simple pole along \(D\), and no other pole at its generic
point.  We say that \(\sigma\) has the relative logarithmic bound at
this prime if
\begin{equation}\label{generaltype:relative-log-bound}
  \ord_E\bigl(\sigma\wedge h^*\omega_D^{\,b}\bigr)\geq-b.
\end{equation}
The notation \(\wedge\) denotes the relative-times-base identification
of the corresponding pluricanonical line bundles.

\begin{lemma}\label{generaltype:stable-orders}
In the situation of \Cref{generaltype:stable-comparison}, pull back a
holomorphic section of a divisible multiple of
\(K_{\mathcal V/R}+\mathcal A\), and compare it with relative pluriforms
on the main component of
\(V\times_Y Y'\).
It satisfies the horizontal bound given by the pullback of \(A\).
Moreover, on a common resolution it satisfies
\Cref{generaltype:relative-log-bound} at every prime lying over a
divisor of any smooth modification of the comparison base.
\end{lemma}

\begin{proof}
On general fibers, pullback from the log canonical model lies in the
original log pluricanonical system.  The effective canonical-model
discrepancy on a common resolution therefore proves the horizontal bound.

For the vertical assertion, first fix a smooth modification of the
comparison base and a prime divisor on it.  The stable family and its
divisible relative adjoint line bundles commute with this base change.
Work transversely to a general point of that divisor.  The resulting
one-parameter family is still stable.  This order of operations is
what permits the same bound at divisors exposed only on a higher base
model.  Its total pair with the reduced
central fiber added is log canonical.  More explicitly, choose the
disk with klt fibers away from its origin.  Its total space is normal:
flatness over the smooth disk and the \(S_2\) property of the fibers
give \(S_2\).  A codimension-one point on the central fiber is a
generic point of one of its reduced components.  Its local ring is
one-dimensional, and its quotient by the disk parameter is a field;
its maximal ideal is consequently generated by that parameter, so
this local ring is regular.  Away from the central fiber the total
space is normal, because its fibers are klt.  Serre's criterion now
gives normality.
Compatibility of the relative adjoint divisor with the stable fiber
identifies the adjunction different on the normalization of the central
fiber with its stable boundary and conductor.  The resulting pair is
log canonical because the central fiber is stable; see
\cite[Lemma~2.10 and Corollary~2.11]{PatakfalviFibered} for this
normalization-and-conductor identity.
The analytic Cartier-divisor inversion theorem
\cite[Theorem~1.1]{FujinoAnalyticAdjunction} consequently applies.
The argument also applies after any finite ramified disk change.

Let \(t\) be the disk parameter and let \(\xi\) denote the
pulled-back relative section in degree \(b\).  Then
\(\xi\wedge(dt/t)^b\) is a section of the total log canonical
multiple with the central fiber added.  On a log resolution the
log canonical discrepancy bound says that this form has pole
order at most \(b\) at every vertical prime.  A further common
resolution with the original family preserves this bound.
This is \Cref{generaltype:relative-log-bound}.  Smooth transverse
parameters do not change any of these orders.
\end{proof}

\begin{lemma}\label{generaltype:relative-system}
Under the hypotheses of \Cref{generaltype:stable-comparison}, there
is a positive integer \(b_0\) and a finite-dimensional space
\[
  \mathcal S\subset
  H^0_{\mathrm{mer}}\bigl(V,\mathcal O_V(b_0K_{V/Y})\bigr)
\]
whose associated meromorphic map has rank \(\dim F\) on general
\(h\)-fibers, with the following properties:
\begin{enumerate}[label=\textup{(\roman*)}]
\item every member has horizontal pole bound \(b_0A\);
\item every member satisfies the relative logarithmic bound
at primes over base divisors;
\item the same bound may be tested after any modification of
the base, at a divisor exposed on that modification.
\end{enumerate}
The same assertions hold for the systems of products of members
of \(\mathcal S\).
\end{lemma}

\begin{proof}
Use \Cref{generaltype:stable-comparison}.  If \(\dim R>0\), the
stable-family bigness theorem
\cite[Theorem~8.1 and Corollary~8.3]{KovacsPatakfalvi} applies:
the base is normal projective, the family is stable and has
maximal variation, and its general fiber is klt.
Thus \(K_{\mathcal V/R}+\mathcal A\) is big.
Choose finitely many sections of one sufficiently divisible
multiple whose map is generically finite.  Its image has dimension
\(\dim F+\dim R\), and its restriction to a general \(\pi\)-fiber
has rank \(\dim F\).  If \(R\) is a point, choose such
sections of the ample log canonical divisor of the fixed stable
pair instead.  Pullback and the meromorphic comparison preserve these
image dimensions and give the pole bounds of
\Cref{generaltype:stable-orders}.

Replace the comparison cover by a Galois cover and resolve the
diagrams equivariantly.  To do this analytically, first take its Stein
factorization over \(Y\); its finite factor has the same cover over a
dense open set.  Over the finite étale locus, the Galois closure is a
component of a sufficiently large fiber product of this finite factor.
Take the corresponding compact component, normalize, and then resolve
the graph of its map to the comparison cover.  Equivariant resolution
gives the required diagram; write \(G\) for its finite group.
This construction uses finite analytic maps and does not require the
meromorphic function field of \(Y\) to have transcendence degree
\(\dim Y\).

Let \(\mathcal B\) be the graded \(\C\)-algebra generated by the
chosen finite set of relative sections and all its \(G\)-conjugates.
Every homogeneous element has the required horizontal and vertical
linear pole bounds: products add the orders and sums cannot
decrease the lower bounds.  The algebra is integral and finite over
\(\mathcal B^G\).  Indeed each homogeneous generator satisfies its
orbit polynomial with invariant coefficients, and invariants of
a finite group acting on a finitely generated algebra are finitely
generated.  The corresponding map
\(\operatorname{Proj}\mathcal B\to\operatorname{Proj}\mathcal B^G\)
is finite.  Before taking invariants, the system contains the sections
pulled back from \(\mathcal V\), whose restriction to a general source
fiber has image dimension \(\dim F\).  A finite map preserves the
dimension of the image of that fiber, even though \(G\) may permute
the fibers above one point of \(Y\).  Choose a Veronese degree for
\(\mathcal B^G\) generated by finitely many elements of that single
degree.  Their ratios therefore still have rank \(\dim F\) on general
fibers.  Denote their pluricanonical degree by \(b_0\).
The same finite quotient also preserves the total image dimension,
which is at least \(\dim F+\dim R\).

Use the natural identification of relative canonical bundles on the
common smooth locus for every base change or modification.  On the
dense étale locus, invariant relative
tensors therefore descend to \(V\).  They extend meromorphically
across the remaining locus by finite analytic descent.  The
horizontal bounds descend because the comparison cover is
generically étale along each horizontal divisor.

For completeness, the vertical order calculation is exact.
Let \(E'\) lie over \(E\), with ramification index \(e\).
For a degree-\(b_0\) top pluriform \(\alpha\),
\begin{equation}\label{generaltype:ramified-order}
  \ord_{E'}(\alpha_{\mathrm{pullback}})+b_0
       =e\bigl(\ord_E(\alpha)+b_0\bigr).
\end{equation}
A logarithmic top form on the base pulls back to a nonvanishing
logarithmic top form at a prime dominating its divisor.
Apply \Cref{generaltype:ramified-order} to the relative form
wedged with this base logarithmic form.
The bound upstairs is therefore equivalent to
\Cref{generaltype:relative-log-bound} downstairs.
For (iii), pull the chosen Galois comparison cover and its original
sections to a common refinement over the modified base.
Lemma~\ref{generaltype:stable-orders} applies to these same pullbacks
at primes above the newly exposed base divisor; its bounds pass to
their conjugates, products, and invariant combinations.  The
ramification calculation therefore gives the bound for the pullback
of each fixed descended member of $\mathcal S$.
Products add the order inequalities, and taking a full product system
is a Veronese operation on its image, so products preserve the relative
rank as well.
\end{proof}

\begin{lemma}\label{generaltype:expose-valuation}
Let \(h\colon V\to Y\) be a proper surjective holomorphic map
of irreducible compact complex spaces, with \(Y\) smooth.
If a prime divisor \(E\) of a smooth \(V\) does not dominate \(Y\),
there is a modification of the base and a dominating main-component
model on which the strict transform of \(E\) maps onto a prime
divisor of the modified base.
\end{lemma}

\begin{proof}
Apply analytic flattening to \(h\), by strict transform, and resolve
the modified base.  Pullback preserves the flatness of the intermediate
family, whose fibers have the generic relative dimension \(d\).
This intermediate source is proper and bimeromorphic over \(V\).
On a resolution dominating it, the strict transform of \(E\) remains
a divisor, because a proper modification of the smooth \(V\) is an
isomorphism at the generic point of \(E\).  Its image in the flat
intermediate source has dimension \(\dim V-1\): it still dominates
\(E\).  The fiber-dimension bound consequently gives an image on the
new base of dimension at least
\(\dim V-1-d=\dim Y-1\).
That image is proper, since its projection to \(Y\) is contained in
\(h(E)\), and properness makes it a closed analytic set.  It is
therefore a prime divisor.  Further source resolutions retain the
same generic image and do not change the valuation of \(E\).
\end{proof}

\subsection{Proof of the addition proposition}

\begin{proof}[Proof of \Cref{generaltype:addition}]
We may suppose \(\kappa(Y,K_Y+C)\geq0\).
Choose a Kähler modification \(\nu\colon\widetilde V\to V\),
resolve the boundary,
and put
\[
  \widetilde T=\nu_*^{-1}T+\operatorname{Exc}(\nu)_{\mathrm{red}}.
\]
By \Cref{setup:log-modification,setup:modification-compatibility},
pullback and descent identify the divisible log pluricanonical rings
on the total spaces and on their very general fibers.  In particular,
fiberwise bigness is preserved, and
\(B(h\circ\nu,\widetilde T)=B(h,T)\).
Every prime over base codimension at least two still has coefficient
one: it is either a strict transform of such a prime or is newly
exceptional.  Write \((V,T)\) for this Kähler model; its sections
descend to the original pair.

If \(\dim F=0\), the map is bimeromorphic.  Pull back sections
of \(K_Y+C\).  At a prime dominating a base divisor the
multiplicity is one and the condition \(C\leq B(h,T)\) gives
the required pole bound.  At an exceptional prime the
log canonical discrepancy bound for \((Y,C)\), together with
\(T_E=1\), gives the same conclusion.  Thus pullback preserves
all the base systems and proves the assertion.  If \(Y\) is a
point, the assertion is simply the assumed bigness.  We may
therefore assume both dimensions are positive.

Write \(T^{\mathrm h}\) for the horizontal part of \(T\).
Choose a rational \(\epsilon>0\) sufficiently small and set
\[
  A=(1-\epsilon)T^{\mathrm h}.
\]
Its coefficients are strictly less than one.
To make \(K_F+A_F\) big uniformly on very general fibers, consider
\(\epsilon=1/j\) and all denominator-clearing degrees.  Choose a fiber
outside their countably many generic base-change and relative-rank
exceptional sets.
Openness of the big cone on that projective fiber gives one
choice of \(\epsilon\) and one degree with full image dimension.
The generic-rank choice makes the same degree have full rank
on general fibers.  Notice that no vertical coefficient has
been lowered, and \(C\) has not been changed.

Apply \Cref{generaltype:relative-system} to \(A\), and write
\(\sigma\) for a member of its system in degree \(b\), including
any product degree.  Let
\[
  \eta\in H^0\bigl(Y,\mathcal O_Y(b(K_Y+C))\bigr).
\]
The relative-times-base product
\(\sigma\wedge h^*\eta\) is a meromorphic section of
\(\mathcal O_V(bK_V)\).  We check that its poles are bounded by
\(bT\) at every prime.

There is nothing further to check at horizontal primes, since
\(A\leq T^{\mathrm h}\).  Suppose next that \(E\) maps onto a
prime divisor \(D\subset Y\), and put
\[
  a_E=\ord_E(h^*D).
\]
Relative to a local nonvanishing logarithmic base volume
\(\omega_D\), the section \(\eta\) vanishes to order at least
\(b(1-C_D)\).  The relative logarithmic bound gives
\begin{equation}\label{generaltype:orbifold-order}
  \ord_E(\sigma\wedge h^*\eta)
        \geq -b+b\,a_E(1-C_D).
\end{equation}
The inf-multiplicity convention is precisely
\[
 B(h,T)_D
   =1-\max_{E\mapsto D}\frac{1-T_E}{a_E}.
\]
Consequently \(C_D\leq B(h,T)_D\) implies
\[
 a_E(1-C_D)\geq1-T_E,
\]
and \Cref{generaltype:orbifold-order} is at least \(-bT_E\).
This uses no integrality condition on the boundary multiplicities.

Finally, suppose \(\operatorname{codim}_Y h(E)\geq2\).
Expose a base divisor \(D'\) below the strict transform of \(E\)
by \Cref{generaltype:expose-valuation}, and resolve the diagrams.
The pullback of \(\eta\) has at most a logarithmic pole along
\(D'\): this is the divisorial discrepancy inequality for the
log canonical smooth pair \((Y,C)\).
Property (iii) of \Cref{generaltype:relative-system} therefore
gives order at least \(-b\) at the strict transform of \(E\).
The order is unchanged, since source modifications are isomorphisms
at the generic point of \(E\).
The assumption \(T_E=1\) supplies exactly this allowance.
We have proved, at all prime divisors, that
\begin{equation}\label{generaltype:product-sections}
 \sigma\wedge h^*\eta
       \in H^0\bigl(V,\mathcal O_V(b(K_V+T))\bigr).
\end{equation}
On a smooth space these divisorial bounds imply holomorphic
extension of the corresponding log pluricanonical section.

Choose a degree in which the complete base system has image
dimension \(\kappa(Y,K_Y+C)\), and pass to a common multiple
with \(b_0\) in \Cref{generaltype:relative-system}.
The product system of relative sections still has image dimension
\(d=\dim F\) on a general fiber.  The corresponding product of the
chosen base system has image dimension \(\kappa(Y,K_Y+C)\), since a
Veronese embedding preserves its image.  In this common degree write
the nonzero relative generators as \(\sigma_i\) and the base
generators as \(\eta_j\).
The total system in \Cref{generaltype:product-sections} contains
all \(\sigma_i\wedge h^*\eta_j\).
On the open set where \(\sigma_0\) and \(\eta_0\) are nonzero,
ratios in this system include
\[
  \frac{\sigma_i}{\sigma_0},
  \qquad
  h^*\!\left(\frac{\eta_j}{\eta_0}\right).
\]
Its image therefore maps onto the base Iitaka image, and over
a general point of that image it contains the \(d\)-dimensional
image of a general \(h\)-fiber.
The dimension-of-fibers inequality gives image dimension at
least \(d+\kappa(Y,K_Y+C)\).  This proves
\Cref{generaltype:inequality} on the Kähler model.
The initial log pluricanonical descent then proves it on the
given smooth source.
\end{proof}

\section{Return to the original fibration and consequences}
\label{application:section}

The previous sections supply the two positivity statements and the
exact section comparison. We now apply them to the original fibration,
keeping its base fixed while changing the intermediate Iitaka space.

\begin{proof}[Proof of \Cref{main:orbifold}]
If either summand on the right is \(-\infty\), the asserted inequality
is automatic.  Suppose henceforth that both are nonnegative, and put
\[
 k=\kappa(F,K_F+\Delta_F),\qquad b=\kappa(f\mid\Delta).
\]

\smallskip\noindent
\emph{Fixing the base.}
Use Section~\ref{setup:section} and \Cref{setup:flattening} to resolve the diagram,
take K\"ahler models, and flatten before resolving the source. Use
strict transforms plus reduced exceptional divisors at every source
modification. Resolve the discriminant and boundary-stratum images
on the base, then fix this smooth K\"ahler base \(Y\). The source
and very general fiber log section rings are unchanged, and every
source prime mapped into codimension at least two in \(Y\) has
coefficient one. Put
\[
 C=B(f,\Delta)
\]
on this model. It has SNC support, and the definition of the invariant
gives
\begin{equation}
 \kappa(Y,K_Y+C)\geq b.
 \label{application:invariant-bound}
\end{equation}
\Cref{setup:simultaneous} preserves \(C=B(f,\Delta)\) and the
coefficient-one condition through all further changes over \(Y\).

\smallskip\noindent
\emph{The relative Iitaka map and its Hodge line.}
Apply \Cref{application:relative-iitaka}, obtaining
\[
 X\xrightarrow{g}W\xrightarrow{h}Y,
 \qquad \dim W_y=k,
 \qquad \kappa(G,K_G+\Delta_G)=0.
\]
Make \(W\) K\"ahler, resolve the diagram, and fix the current smooth
source pair \((X_0,\Delta_0)\) as reference. Flatten over \(W\)
and resolve the generator and boundary data as required by
\Cref{rankone:comparison}. Thus every prime mapped by \(g\) into
codimension at least two is exceptional over \(X_0\), and on a dense
smooth log open the relative
systems have rank one in all sufficiently divisible degrees.

Apply \Cref{rankone:comparison,period:adjoint}. Their further
modifications remain over \(Y\); \Cref{setup:simultaneous} preserves
the exceptional reference, while the variation and its parabolic line
pull back from the common open. The relative Iitaka property persists
by \Cref{application:relative-iitaka}. On the resulting common model
we obtain a rational line \(M\), a rational SNC boundary \(T\) in
\([0,1]\), and a surjective map to a smooth projective variety:
\begin{equation}
 p:W\longrightarrow S,
 \qquad M=p^*P,
 \qquad P\text{ nef},
 \qquad K_S+a_0P\text{ big for some }a_0>0,
 \label{application:period-data}
\end{equation}
such that
\begin{equation}
 B(g,\Delta)\leq T\leq1.
 \label{application:rankone-boundary}
\end{equation}
Compute \(T\) by the rank-one order formula on this final model.
Parabolic pullback and \Cref{rankone:comparison} identify the divisible
systems of \(K_W+T+M\) with those on \((X_0,\Delta_0)\), also
relatively over \(Y\), preserving all ratios.

\smallskip\noindent
\emph{Relative bigness and the base boundary.}
Set \(D_W=K_W+T\).  The systems defining the relative Iitaka image
have image dimension \(k\) on very general \(f\)-fibers. Push these systems
into \(D_W+M\) using \Cref{rankone:comparison}.  Their ratios are
unchanged, so the resulting systems have image dimension \(k\) on
very general \(h\)-fibers, whose dimension is \(k\).  Consequently
\begin{equation}
 (D_W+M)|_{W_y}\text{ is big for very general }y.
 \label{application:relative-big}
\end{equation}
Here the systems are obtained by generic base change over an open patch
of \(Y\) and extend on the smooth \(W_y\) by the relative section
comparison. Adjunction identifies the restricted adjoint bundles.

By \Cref{application:composition},
\begin{equation}
 B(h,T)\geq C.
 \label{application:composed-boundary}
\end{equation}
The same lemma gives coefficient one in \(T\) for every prime of
\(W\) mapped into codimension at least two in \(Y\). Thus both
boundary hypotheses of \Cref{generaltype:addition} hold.

The hypotheses of \Cref{addition:principle} have now all been
verified. Its conclusion is
\begin{equation}\label{application:final-intermediate-bound}
 \kappa(W,K_W+T+M)\geq k+\kappa(Y,K_Y+C).
\end{equation}
The proof of that proposition uses a single nonzero section to cancel
the ample perturbation and preserves the ratios defining this bound.

Finally, \Cref{rankone:comparison} lifts these systems by the relative
generators, preserving their ratios and image dimensions. Untested
\(g\)-contracted primes are exceptional over \((X_0,\Delta_0)\),
so the forms descend and extend there. By
\Cref{setup:log-modification}, they then descend to the original source.
Combining \eqref{application:invariant-bound} and
\eqref{application:final-intermediate-bound} proves
\[
 \kappa(X,K_X+\Delta)
     \geq \kappa(F,K_F+\Delta_F)+\kappa(f\mid\Delta).
\]

Bigness on a point means Iitaka dimension zero, and a zero-dimensional
connected fiber has trivial relative line. If \(Y\) is a point the
inequality is an identity; if the relative Iitaka fiber is a point its
Hodge line is trivial. Together with the point-period-base case above
and the initial \(-\infty\) convention, these observations cover all
dimensions and values of the two summands.
\end{proof}

\subsection{The logarithmic and ordinary inequalities}
\label{spec:section}

We first compare the invariant orbifold base with a prescribed
logarithmic base.  This makes the two classical specializations of
Theorem~\ref{main:orbifold} immediate, while retaining control of the
base under modification.

\begin{lemma}[Comparison with a logarithmic base]
\label{spec:base-comparison}
Let $f:(X,\Delta)\to Y$ be a fibration of smooth compact complex
manifolds, where $\Delta$ has rational coefficients in $[0,1]$ and
simple normal crossing support.  Let $D_Y$ be a reduced simple normal
crossing divisor on $Y$.  Suppose that every prime component of
$f^*D_Y$ has coefficient one in $\Delta$.  Then
\[
  \kappa(f\mid\Delta)\geq\kappa(Y,K_Y+D_Y).
\]
\end{lemma}

\begin{proof}
Choose a neat model as in Lemma~\ref{setup:flattening}, with
$p:X'\to X$, $q:Y'\to Y$, $f':X'\to Y'$, and
$\Delta'=p_*^{-1}\Delta+\operatorname{Exc}(p)_{\mathrm{red}}$.
We may arrange that
$E_Y=(q^{-1}\operatorname{Supp}D_Y)_{\mathrm{red}}$ has simple
normal crossings.  Every prime of $X'$ dominating a component of
$E_Y$ has coefficient one in $\Delta'$: an exceptional prime has
coefficient one by definition, and a nonexceptional prime is the strict
transform of a component of $f^*D_Y$.  Thus its weighted multiplicity
is infinite, and the inf-multiplicity rule gives
\[
 B(f',\Delta')\geq E_Y.
\]
Pullback by $q$ takes logarithmic top forms with poles along $D_Y$
to logarithmic top forms with poles along $E_Y$.  In local coordinates
this follows by pulling back each factor $dz_i/z_i$; the same assertion
holds for all tensor powers.  Consequently, pullback embeds every
plurilogarithmic system on $(Y,D_Y)$ into the corresponding system of
$K_{Y'}+B(f',\Delta')$, preserving all ratios of sections.
The neat-model identity~\eqref{setup:neat-formula} now proves the
claim.
\end{proof}

\begin{corollary}[Logarithmic subadditivity]
\label{spec:logarithmic}
Let $f:X\to Y$ be a surjective morphism with connected fibers between
smooth connected projective complex varieties.  Let $D_X,D_Y$ be
reduced simple normal crossing divisors, with the zero divisor allowed,
such that
\[
 \operatorname{Supp}(f^*D_Y)\subseteq\operatorname{Supp}(D_X).
\]
For a very general fiber $F$, put $D_F=D_X|_F$.  Then
\begin{equation}\label{spec:log-inequality}
 \kappa(X,K_X+D_X)
 \geq\kappa(F,K_F+D_F)+\kappa(Y,K_Y+D_Y).
\end{equation}
The same statement holds for compact manifolds in Fujiki class
$\mathcal C$ and holomorphic fibrations between them.
\end{corollary}

\begin{proof}
Theorem~\ref{main:orbifold} applies to $(X,D_X)$, and
Lemma~\ref{spec:base-comparison} bounds its base term from below.
If either term on the right of~\eqref{spec:log-inequality} is
$-\infty$, the assertion follows from our convention for that value.
\end{proof}

In the projective setting, this is the usual compactification form of
the logarithmic Iitaka conjecture.  Indeed, for a smooth
quasi-projective variety $U$ with a smooth projective compactification
$X$ and reduced simple normal crossing boundary $D_X=X\setminus U$,
its logarithmic Kodaira dimension is
$\overline\kappa(U)=\kappa(X,K_X+D_X)$
\cite{IitakaLog}.  A dominant morphism $U\to V$ between smooth
connected quasi-projective varieties, with geometrically connected
general fiber, admits compatible such compactifications:
compactify the graph and resolve the graph and boundaries.  The
Stein factor of the resulting morphism $X\to Y$ is finite
birational over the normal variety $Y$, hence equals $Y$.
Thus $f$ has connected fibers, and
$\operatorname{Supp}(f^*D_Y)\subseteq\operatorname{Supp}(D_X)$.
For very general $v\in V$, $(X_v,D_X|_{X_v})$ compactifies $U_v$.
Thus Corollary~\ref{spec:logarithmic} gives
$\overline\kappa(U)\geq\overline\kappa(U_v)+\overline\kappa(V)$.

Ordinary subadditivity also enters the Albanese reduction of
\cite{OpenAILogAbundance2026}. That work combines
it with separate nonvanishing and good-model arguments to prove log
abundance in characteristic zero.

\begin{corollary}[Ordinary subadditivity in characteristic zero]
\label{spec:ordinary}
Let $k$ be an algebraically closed field of characteristic zero, and
let $f:X\to Y$ be a surjective morphism with connected fibers between
smooth connected projective $k$-varieties.  If $F$ is its geometric
generic fiber, then
\begin{equation}\label{spec:ordinary-inequality}
 \kappa(X)\geq\kappa(F)+\kappa(Y).
\end{equation}
More generally, the logarithmic inequality of
Corollary~\ref{spec:logarithmic}, with the same boundary hypotheses,
holds over $k$, with the geometric generic fiber in place of a very
general complex fiber.
\end{corollary}

\begin{proof}
Over $\mathbb C$, the ordinary assertion follows by taking
$D_X=D_Y=0$ in Corollary~\ref{spec:logarithmic}.  In a projective
family over $\mathbb C$, the Iitaka dimension of the geometric
generic fiber equals that of a very general fiber.  To see this,
apply cohomology and base change in each positive divisible
pluricanonical degree.  After shrinking the base separately in each
degree, the corresponding relative rational map also has constant
fiberwise image dimension.  Excluding the resulting countable
union of proper closed subsets identifies all these image
dimensions with those of the generic fiber.

For the field extension, let $K\subseteq K'$ be fields and let $V$
be a smooth projective geometrically integral $K$-variety with a
rational divisor $D$.  For every $m$ for which $mD$ is integral,
flat base change gives
\[
 H^0(V,\mathcal O_V(mD))\otimes_K K'
   \simeq H^0(V_{K'},\mathcal O_{V_{K'}}(mD_{K'})).
\]
The complete linear system on the right is the scalar extension of
the one on the left, so nonvanishing and image dimension are
unchanged.  Hence Iitaka dimension is invariant under extension
of the ground field.  Applying this observation to the generic
fiber over the function field of the base proves the corresponding
invariance for geometric generic fibers.

Descend $f$, its projective embeddings, and its boundary divisors
to a finitely generated subfield $K\subset k$ over $\mathbb Q$.
The descended varieties are smooth and geometrically integral;
the boundary conditions descend as well.  The equality
$f_*\mathcal O_X=\mathcal O_Y$ descends by proper flat base change
and faithful flatness, and therefore also holds after an embedding
$K\hookrightarrow\mathbb C$.  Apply the complex logarithmic
inequality to this base change.  The preceding invariance identifies
its three Iitaka dimensions with those over $k$, proving both
assertions.
\end{proof}

\subsection{Kodaira dimension along the core}
\label{spec:core}

A smooth compact orbifold pair $(X,\Delta)$ in class $\mathcal C$ is
\emph{special} if it has no meromorphic fibration onto a
positive-dimensional base of orbifold general type; here general
type means that the invariant base dimension equals the dimension
of the base.  Campana constructed its \emph{core}
\[
 c_{(X,\Delta)}:(X,\Delta)\dashrightarrow C(X,\Delta),
\]
whose very general orbifold fiber is special and whose orbifold base is of general
type, unless it is a point
\cite[Theorem~9.1]{CampanaOrbifolds}.  In particular, the core has
a point base exactly when $(X,\Delta)$ is special. Fiber properties
of the core and the tower below refer to Campana's stable orbifold
fibers: they are computed on the suitable equivalent holomorphic
representatives furnished by these constructions
\cite[Definition~2.49 and \S7.1]{CampanaOrbifolds}.
A neat representative $g:(Z,\Gamma)\to T$ in the sense of
Section~\ref{setup:section} is \emph{strictly neat} if its
inf-multiplicity boundary on the smooth base is SNC. It is
\emph{high} if $g$ is an orbifold morphism to this base: for every
base prime $D$ and source prime $E$ with
$a_E=\ord_E(g^*D)>0$, including primes contracted by $g$,
\[
              a_E m_\Gamma(E)\geq m_{B(g,\Gamma)}(D).
\]
These are the model conditions of
\cite[Definition~3.8 and Proposition~3.10]{CampanaOrbifolds}.
The high condition is not automatic: the inf-multiplicity base tests
divisors dominating base primes, whereas an orbifold morphism must
also control primes contracted into higher codimension
\cite[Remark~3.3(2)]{CampanaOrbifolds}.
For the core and its tower we choose compatible strictly neat and
high representatives supplied by Campana's constructions
\cite[Proposition~3.15 and Corollary~4.14]{CampanaOrbifolds}.
The section comparisons of Section~\ref{setup:section} do not
assert that arbitrary model changes preserve every stable-fiber
property.

\begin{corollary}[Kodaira dimension along the core]
\label{spec:core-dimension}
Let $(X,\Delta)$ be a smooth compact orbifold pair in Fujiki class
$\mathcal C$, with rational simple normal crossing boundary
coefficients in $[0,1]$.  Let $C=C(X,\Delta)$ be its core base and
$(G,\Delta_G)$ a very general orbifold fiber.  Then
\begin{equation}\label{spec:core-equality}
 \kappa(X,K_X+\Delta)
   =\kappa(G,K_G+\Delta_G)+\dim C.
\end{equation}
If $\kappa(X,K_X+\Delta)\geq0$, then
\[
 0\leq\dim C\leq\kappa(X,K_X+\Delta),
\]
and equality in the second inequality holds precisely when
$\kappa(G,K_G+\Delta_G)=0$.
In particular, $\kappa(X,K_X+\Delta)=0$ implies that
$(X,\Delta)$ is special.
\end{corollary}

\begin{proof}
Resolve the core map, using the boundary convention of
Lemma~\ref{setup:log-modification}.  The total and fiber Iitaka
dimensions are preserved.  Theorem~\ref{main:orbifold} gives
the lower bound in~\eqref{spec:core-equality}, since the invariant
base dimension is $\dim C$.

For completeness, the reverse bound is the elementary upper
addition inequality.  For each divisible $m$, let $\phi_m$ be
the meromorphic map defined by $|m(K_X+\Delta)|$, when this
system is nonempty.  The map $(c,\phi_m)$ has image dimension
at most $\dim C+\kappa(G,K_G+\Delta_G)$: on a very general
$c$-fiber its second component is given by a subspace of
$H^0(G,m(K_G+\Delta_G))$.  Since projection from this image
onto the image of $\phi_m$ is surjective, the same upper bound
holds for $\dim\phi_m(X)$.  Taking the maximum over $m$ proves
the desired inequality.  If the fiber Iitaka dimension is
$-\infty$, restriction to a very general fiber shows that no
global pluricanonical section can be nonzero, so both sides of
\eqref{spec:core-equality} are $-\infty$.

Finally, if the total Iitaka dimension is nonnegative, a nonzero
section restricts nontrivially to a very general core fiber.
Its Iitaka dimension is therefore nonnegative.  The dimension
bound and its equality case follow from~\eqref{spec:core-equality}.
When the total dimension is zero, the core base is a point,
which is equivalent to specialness.
\end{proof}

These statements recover, for the core, Campana's general-type-base addition
theorem and its consequences
\cite[Theorems~6.3, 6.7, and~6.8]{CampanaOrbifolds}.
The existence of the core and the implication from Kodaira
dimension zero to specialness are established results of his
theory.

\subsection{The core tower}
\label{spec:core-tower-section}

For a smooth orbifold pair $(Z,\Theta)$, the notation
$\kappa_+(Z,\Theta)=-\infty$ means that every meromorphic
fibration onto a positive-dimensional base has invariant
orbifold Kodaira dimension $-\infty$; the condition is vacuous
for a point.  Equivalently, one may test all dominant meromorphic
maps: in Stein factorization, the finite orbifold canonical-divisor
inequality shows that passing to the connected-fiber factor cannot
decrease the invariant base dimension.  Indeed, over a divisor
with multiplicity $m_D$, a prime with ramification index $e$ and
orbifold multiplicity $m_E$ satisfies $m_D\leq e m_E$; the
canonical-divisor difference has coefficient
$e/m_D-1/m_E\geq0$.
The \emph{$\kappa$-rational quotient} is a meromorphic fibration
$r:(Z,\Theta)\dashrightarrow R$ with $\kappa(r\mid\Theta)\geq0$
whose very general stable orbifold fiber has $\kappa_+=-\infty$
\cite[Definition~5.23 and Corollary~6.14]{CampanaOrbifolds}.
Its \emph{stable orbifold base} is the inf-multiplicity base on the
compatible representatives just described
\cite[Proposition~3.15, Corollary~4.14, and \S7.1]{CampanaOrbifolds}.

\begin{corollary}[Campana's decomposition of the core]
\label{spec:core-tower}
Let $(X,\Delta)$ be as in Corollary~\ref{spec:core-dimension},
and set $n=\dim X$.  Starting from $(X,\Delta)$, successively
take the $\kappa$-rational quotient $r$ and the
Moishezon--Iitaka fibration $M$ of its stable orbifold base.
At each iteration choose such compatible representatives and equip
the new base with the stable inf-multiplicity orbifold structure
of the composite.
The composite after at most $n$
iterations is the core, up to bimeromorphic equivalence of
fibrations:
\[
 c_{(X,\Delta)}=(M\circ r)^n.
\]
The positive-dimensional very general stable orbifold fibers of
the $r$-steps have $\kappa_+=-\infty$, and those of the
$M$-steps have Kodaira dimension zero.  In particular, for
a special pair this tower ends at a point.
\end{corollary}

\begin{proof}
Theorem~\ref{main:orbifold} establishes the hypothesis of
Campana's Corollary~6.14, which supplies the
$\kappa$-rational quotient.  Its scope is exactly that used
here: smooth compact class-$\mathcal C$ pairs with rational
coefficients in $[0,1]$
\cite[Conjecture~6.1 and Remark~6.2]{CampanaOrbifolds}.
Its stable base has nonnegative Kodaira dimension, so $M$
is defined and has very general orbifold fiber of Kodaira
dimension zero.  The new smooth base remains in class
$\mathcal C$, with rational boundary coefficients in $[0,1]$.
Thus the construction can be repeated.

Campana's Lemma~10.1 states that each composite $M\circ r$
has special general orbifold fiber and decreases dimension
unless its source pair is of general type.  His
Theorems~10.2--10.3 identify the resulting stabilized
composite with the core.  Their sole conjectural input is
the existence of the $\kappa$-rational quotient, now supplied
above.  Strict dimension decrease allows at most $n$
nontrivial iterations; identity steps give the displayed
formula.  A point is fixed throughout.
These are precisely the decomposition and fiber assertions of
\cite[Lemma~10.1, Theorems~10.2--10.3, and Corollary~10.4]
{CampanaOrbifolds}.
\end{proof}

\section{Minimal root covers for reduced boundaries}\label{cyc:section}

For a smooth projective pair of logarithmic Kodaira dimension zero
with reduced SNC boundary, the minimal cyclic cover has a
one-dimensional space of logarithmic pluriforms in every positive
degree. This strengthens
the character-line statement of Proposition~\ref{rankone:comparison}.
We prove the refinement together with its exact Hodge and divisorial
normalization. It gives a pure integral realization of the entire
highest line and fixes the normalization used in further
relative-Iitaka constructions.
The normalization lemmas below compare the minimum over actual fiber components
with the threshold over all divisorial valuations. They show that a
zero coefficient of the normalized base boundary yields a
multiplicity-one component, and that later permitted modifications
preserve the complete section spaces. These are the additional
normalization statements used by the variation companion.

\begin{lemma}[The cyclic cover of a logarithmic Kodaira-zero pair]
\label{lem:reduced-cyclic-top}
Let $(J,D)$ be a smooth geometrically integral projective pair over a
field $K$ of characteristic zero. Assume that $D$ is reduced and has
simple normal crossings and that
$\kappa(J_{\overline K},K_{J_{\overline K}}+D_{\overline K})=0$.
Choose the least positive integer $p$ for which
$H^0(J,p(K_J+D))\ne0$, a nonzero section $\omega$ in that space,
and put
\[
 E=\operatorname{div}(\omega)+pD,
 \qquad \Delta=D-\frac1pE.
\]
Write $\omega=b\eta^{\otimes p}$ in a rational canonical frame,
and let $\tau:N\to J$ be the normalization in
$K(J)[t]/(t^p-b)$. Let $\rho:\widetilde N\to N$ be a log
resolution, and set
\[
 P=\operatorname{Supp}(\Delta^{=1}),\qquad
 D_{\widetilde N}=\bigl((\tau\circ\rho)^{-1}P\bigr)_{\mathrm{red}}.
\]
Then $N$ is geometrically integral, the rational top form
$\Omega=t\tau^*\eta$ extends logarithmically to
$(\widetilde N,D_{\widetilde N})$, and
\[
 H^0\bigl(\widetilde N,m(K_{\widetilde N}+D_{\widetilde N})\bigr)
       =K\,\Omega^m\qquad(m\ge1).
\]
\end{lemma}

\begin{proof}
Every nonzero plurilogarithmic section space on $J$ is
one-dimensional: two independent sections of one degree would
define a nonconstant rational map. Consequently $\kappa(E)=0$,
and every effective integral divisor supported on $E$ also has
exactly one section, since it is bounded above by a multiple of $E$.
These assertions and the minimal index $p$ are unchanged by field
extension, by flat base change for global sections.

Over $\overline K(J)$, reducibility of $T^p-b$ would make $b$
an $\ell$th power for some prime $\ell\mid p$. The corresponding
root of $\omega$ would be a nonzero logarithmic pluriform of
degree $p/\ell$, since its logarithmic zero divisor is $E/\ell$.
This contradicts minimality. Thus $N$ is geometrically integral.

At a prime $Q$ of $N$ above a prime $R$ of $J$, write
\[
 a=\operatorname{coeff}_R E,\qquad
 d=\operatorname{coeff}_R D\in\{0,1\},\qquad
 e=\operatorname{ord}_Q(\tau^*R).
\]
The ramification formula for top differentials gives
\begin{equation}\label{eq:cyclic-root-order}
 \operatorname{ord}_Q\Omega=\frac{ea}{p}-ed+e-1.
\end{equation}
If $a=0$, the cover is unramified at $R$, and this order is
$-d$. If $a>0$, then $q=ea/p$ is a positive integer: the
integrality follows from $t^p=b$ and $d\in\mathbf Z$.
The order is $q-1\ge0$ when $d=1$, and $q+e-1\ge0$ when
$d=0$. Thus the divisorial poles are exactly the reduced inverse
image $D_N=(\tau^{-1}P)_{\mathrm{red}}$.

We also check extension across the exceptional divisors of $\rho$.
The subpair $(J,\Delta)$ is sub-log-canonical and is sub-klt
outside $P$. Indeed, its positive part is an SNC boundary bounded
by $D$, with all coefficients strictly below one outside $P$;
subtracting an effective divisor cannot worsen discrepancies.
Define the crepant finite pullback by
$K_N+\Delta_N=\tau^*(K_J+\Delta)$. With the compatible
canonical trivialization, $\Delta_N=-\operatorname{div}(\Omega)$.
For a prime divisor $V$ over $N$, let its restricted valuation
on $K(J)$ be $e_Vv$. The same ramification formula gives
\[
 1+\operatorname{ord}_V\Omega
       =e_V A_{J,\Delta}(v),
\]
where $A$ denotes log discrepancy. The right-hand side is
nonnegative, and is positive if the center lies outside $P$.
Since the left-hand side is integral, $\Omega$ has at most a
simple pole above $P$ and is regular elsewhere. This proves the
asserted logarithmic extension.

Finally let $H_N=\operatorname{div}_N(\Omega)+D_N$. Equation
\eqref{eq:cyclic-root-order} gives the sharper bound
\[
             0\le H_N\le\tau^*E.
\]
For completeness, when $a>0$ and $d=1$ its coefficient is
$q-1\le ea$; when $d=0$, it is
$q+e-1\le q+p-1\le pq=ea$, because $e\le p$ and $q\ge1$.
At $a=0$ its coefficient is zero.
Writing a logarithmic $m$-canonical form as $u\Omega^m$ and
testing its orders at the primes of the finite normalization gives
\[
 \operatorname{div}_N(u)+mH_N\ge0,
 \qquad
 u\in H^0(N,\mathcal O_N(m\tau^*E)).
\]
Finite pullback preserves Iitaka dimension, so
$\kappa(N,\tau^*E)=\kappa(J,E)=0$. Since $\tau^*E$ is
effective and $N$ is geometrically integral, the last space is
$K$. Conversely $\Omega^m$ is a logarithmic pluriform for every
$m\ge1$, proving the equality.
\end{proof}

\begin{corollary}[An integral pure realization of the moduli line]
\label{cor:reduced-integral-moduli}
Let $x:X\to W$ be a projective fibration between smooth complex
projective varieties, with geometrically integral generic fiber
$J$. Let $D_X$ be a reduced SNC boundary whose restriction $D_J$
satisfies the hypotheses of Lemma~\ref{lem:reduced-cyclic-top}.
Let $M_W$ be the rational Hodge divisor of the rank-one reduction,
with its rational frame and parabolic orders as in
Equation~\eqref{rankone:order}. On sufficiently high smooth models,
there is a polarizable integral pure variation
of Hodge structures on a dense smooth open $W^\circ\subset W$
whose top nonzero Hodge filtration piece is a line, and
\[
 M_W\sim_{\mathbf Q}
 \text{the rational canonical extension of that line}.
\]
In the canonical-bundle-formula normalization, this is the moduli
divisor of the auxiliary generic subpair
$\Delta_J=D_J-E_J/p$, which may have negative coefficients.
\end{corollary}

\begin{proof}
Let $s$ be the minimal-degree logarithmic generator on the generic
fiber, and choose a rational relative canonical frame $\eta$ on
$X$. The rational function $s/\eta^{\otimes p}$ defines one
cyclic cover over a dense open of $W$, not merely a collection
of local character covers. Spread its resolution and pole boundary
over a smaller dense open $W^\circ$ so that they form a smooth
projective family of SNC pairs. Write $r=\dim J$, and let
$u:\widetilde X^\circ\setminus D_{\widetilde X}^\circ\to W^\circ$
be the complementary open family. Its geometric variation
\[
 \mathbb V_{\mathbf Z}=R^ru_*\mathbf Z/\mathrm{torsion}
\]
is admissible and graded-polarizable. Logarithmic Hodge theory
identifies
\[
 F^{r+1}\mathbb V=0,\qquad
 F^r\mathbb V_{\mathcal O}
   =\widetilde x_*\omega_{\widetilde X^\circ/W^\circ}
                           (D_{\widetilde X}^\circ).
\]
Lemma~\ref{lem:reduced-cyclic-top} makes this entire top piece a
line. Strictness for the weight filtration therefore singles out
one weight $w$ with
\[
 F^rW_{w-1}\mathbb V=0,
 \qquad F^rW_w\mathbb V=F^r\mathbb V,
 \qquad
 F^r\operatorname{Gr}^W_w\mathbb V\simeq F^r\mathbb V.
\]
The rational variation $\operatorname{Gr}^W_w\mathbb V$ is
polarizable. Intersecting its rational weight filtration with
$\mathbb V_{\mathbf Z}$ and removing torsion supplies an integral
lattice; after rescaling, its rational polarization is integral.
Thus this pure variation has precisely the required top line.

After a finite cover making local monodromies unipotent, the
canonical extensions have locally free double gradeds for the
Hodge and weight filtrations. The displayed identification of top
lines consequently extends across the boundary; see
\cite[Proposition~3.12]{FujinoFujisawa}. Descent gives the same
identification for their rational canonical extensions.

The global rational frame of this entire top line is
$s^{1/p}$. Equation~\eqref{rankone:order} computes its parabolic
order at every prime divisor $P$ as $\beta_P$, the order of the
same rational frame in $M_W$. Equality of these divisorial
orders identifies the rational extensions. In particular, it
identifies $M_W$ itself, with no additional factor of $p$.
We check the rank condition for the canonical-bundle-formula
comparison. On a log resolution $\mu:J'\to J$, let
$R'=\lceil A^*(J,\Delta_J)_{J'}\rceil$, where $A^*$ omits the
log-discrepancy-zero components of the discrepancy divisor.
The subpair is sub-log-canonical, so $R'\geq0$. At the strict
transform of each original prime its coefficient is zero if
$0\leq\Delta_J\leq1$, and is $\lceil-\Delta_J\rceil$ otherwise.
Thus $0\leq\mu_*R'\leq E_J$. Testing original primes gives
$H^0(J',\cO_{J'}(R'))\subseteq H^0(J,\cO_J(E_J))=\C(W)$,
and constants give equality. This is the required rank-one condition.
Extending the generic subpair by
$-p^{-1}\operatorname{div}_{K_{X/W}}(s)$ makes its log canonical
divisor rationally equivalent to $x^*K_W$. The generic sub-lc,
relative triviality, and rank conditions of
\cite[Definition~6.18]{BFMT} therefore hold. Adding a base pullback
to normalize vertical orders changes the total base term and the
discriminant equally, leaving the moduli divisor unchanged.

The least positive index trivializing $K_J+\Delta_J$ is $p$.
Indeed, a smaller such index $q$ would produce a rational
$q$-canonical form with logarithmic zero divisor $qE_J/p\geq0$,
contradicting the minimality of $p$. Thus the cyclic field extension
in Construction-Definition~6.23 of \cite{BFMT} is the one used above.

For its boundary comparison, put
$P_J=\operatorname{Supp}(\Delta_J^{=1})$ and first make a common
crepant log resolution of the auxiliary subpair and its root cover. The
coefficient-one support on the crepant model is contained in the
total inverse image of $P_J$: outside $P_J$ the auxiliary subpair is
sub-klt. Every additional component allowed by our pole boundary
has positive log discrepancy, so the ramification discrepancy
formula in Lemma~\ref{lem:reduced-cyclic-top} makes $\Omega$
regular there. The top-form space for the BFMT pole boundary
therefore injects into the top-form space for our boundary and
contains $\Omega$; both are the line generated by $\Omega$.
After spreading over a common smooth base open, inclusion of the
open complements induces a morphism of their admissible geometric
mixed variations. This morphism is strict for the Hodge and weight
filtrations, including their limiting filtrations after a finite
cover making boundary monodromy unipotent. It therefore identifies
the top lines in the same pure weight grade. Functoriality
of canonical extension gives the same identification across the
boundary and hence for the parabolic rational lines. Theorem~6.28 of
\cite{BFMT} now identifies this Hodge divisor with the moduli
divisor using one common positive multiple on both sides. This
comparison permits negative coefficients in the generic subpair;
it does not use the separate semiampleness theorem requiring
generic effectiveness.
\end{proof}

\begin{remark}
The hypothesis that $D$ is reduced is essential to
Lemma~\ref{lem:reduced-cyclic-top}. For five distinct points on
$\mathbf P^1$, the rational boundary
$D=\frac25(P_1+\cdots+P_5)$ satisfies $K_{\mathbf P^1}+D\sim_{\mathbf Q}0$.
Its index-five root cover has genus six and empty pole boundary.
Thus its entire top-form space has dimension six. The refinement therefore applies to reduced boundaries; the
character-line construction of Section~\ref{rankone:section} is used
for general rational orbifold boundaries.
\end{remark}

\subsection{The original relative orders}


\begin{lemma}[Normalization of the rank-one boundary]
\label{lem:rankone-normalization}
Let $x:X'\to W$ be a projective fibration on smooth models, let $D'$ be
an effective reduced SNC divisor, and suppose that its geometric generic
fiber $(J,D_J)$ has logarithmic Kodaira dimension zero. Fix a nonzero
relative rational $p$-canonical form $s$ generating
$H^0(J,p(K_J+D_J))$. At a prime $P$ of $W$ set
\[
 a_E=\operatorname{ord}_E(x^*P),\qquad
 r_E=p^{-1}\operatorname{ord}^{pK_{X'/W}}_E(s),\qquad
 d_E=\operatorname{coeff}_E(D'),\qquad
 t_P=\min_{E\mapsto P}\frac{r_E+d_E}{a_E}.
\]
Here the minimum is over actual components on the chosen source model.
Choose that model so that the generator divisor, boundary and inverse
image of $P$ have SNC support over the generic point of $P$; later
source modifications carry strict boundary plus reduced exceptional
boundary. The coefficients in the rank-one comparison satisfy
\[
 \alpha_P=t_P,\qquad
 \beta_P=\inf_{E\mapsto P}\frac{1+r_E}{a_E}-1,\qquad
 T_P=1-c_P,
 \quad
 c_P:=\inf_{E\mapsto P}\frac{1+r_E}{a_E}-t_P.
\]
The infima allow divisorial valuations on higher smooth source models.
Equivalently, $c_P$ is the log canonical threshold of $x^*P$ for the
normalized auxiliary subpair
\[
 \Delta^{\mathrm{aux}}
 =-p^{-1}\operatorname{div}_{K_{X'/W}}(s)+x^*\sum_Qt_QQ.
\]
If $T_P=0$, an actual component attaining $t_P$ has
$a_E=1$, $d_E=0$, and attains the threshold infimum as well.
\end{lemma}

\begin{proof}
Choose a local frame $\eta$ of $K_W$ that does not vanish at the generic
point of $P$. The canonical line-bundle identification
\[
 \mathcal O_{X'}(pK_{X'/W})\otimes x^*\mathcal O_W(pK_W)
       \simeq\mathcal O_{X'}(pK_{X'})
\]
identifies the order of $s$ with the order of
$s\wedge x^*\eta^{\otimes p}$. Thus the quantity $l_E$ in the
rank-one order formula is exactly $r_E$. This uses the actual relative
canonical line bundle: the Jacobian contribution is already present in
$r_E$, so no ramification term is added or subtracted here.
Consequently its two order formulas give
\[
 \alpha_P=t_P,
 \qquad
 \beta_P=\min_E\frac{r_E+1-a_E}{a_E}.
\]
Put
\[
 \Delta_0=-p^{-1}\operatorname{div}_{K_{X'/W}}(s),
 \qquad \lambda_P=\min_{E\mapsto P}\frac{1+r_E}{a_E},
\]
where the minimum is over actual vertical components. The indicated
SNC preparation and the generic logarithmic bound give $\Delta_0$ SNC
support and horizontal coefficients at most one. In
$\Delta_0+\lambda_Px^*P$, each vertical coefficient is
$-r_E+\lambda_Pa_E\leq1$. Thus this subpair is sub-log-canonical
over the generic point of $P$, and every higher divisorial valuation
satisfies
\[
                     1+r_E-\lambda_Pa_E\geq0.
\]
An actual component attaining $\lambda_P$ gives equality. Hence the
infimum over all such valuations equals $\lambda_P$, and
$\beta_P=\lambda_P-1$.

For completeness, the first minimum also remains unchanged on higher
source models. A coefficient realizing $t_P$ makes
$t_Pa_E-r_E$ no greater than the original boundary coefficient at every
actual component over $P$. Equivalently, after clearing denominators,
the normalized generating form satisfies every logarithmic order test
over the generic point of $P$. Its pullback remains logarithmic with
strict plus reduced exceptional boundary. No new valuation can lower
the first minimum, while an old minimizing valuation persists.

At any divisorial valuation above the generic point of $P$, the
crepant coefficient of $\Delta^{\mathrm{aux}}$ is $-r_E+a_Et_P$.
The threshold formula therefore gives
\[
 \operatorname{lct}_{\eta_P}(X',\Delta^{\mathrm{aux}};x^*P)
 =\inf_E\frac{1+r_E-a_Et_P}{a_E}=c_P.
\]
Since $T_P=\alpha_P-\beta_P$, we obtain $T_P=1-c_P$.
Finally, if $T_P=0$ and $E$ attains the first minimum, then
\[
 1=c_P\le\frac{1+r_E}{a_E}-t_P
       =\frac{1-d_E}{a_E}\le1.
\]
All terms are equal. Since $a_E$ is a positive integer and
$d_E\in\{0,1\}$, this forces $a_E=1$ and $d_E=0$, as claimed.
\end{proof}

\begin{lemma}[Permitted models after normalization]
\label{lem:normalized-models}
Suppose that $(W,T)$ is smooth with effective SNC rational boundary and
that the moduli divisor $M$ is a rational Cartier divisor on $W$.
For a smooth modification $\mu:W'\to W$, put
\[
 T'=\mu_*^{-1}T+\operatorname{Exc}(\mu)_{\mathrm{red}},
 \qquad M'=\mu^*M,
\]
and assume the resulting boundary is SNC. Then pullback identifies the
section spaces of every sufficiently divisible multiple of
$K_W+T+M$ and $K_{W'}+T'+M'$. At unchanged codimension-one points the
normalization in Lemma~\ref{lem:rankone-normalization} is preserved.
Every new exceptional prime has coefficient one in $T'$.
\end{lemma}

\begin{proof}
Log canonicity of $(W,T)$ gives
\[
 K_{W'}+T'+M'=\mu^*(K_W+T+M)+E,
\]
where $E$ is effective and $\mu$-exceptional. For divisible $m$,
normality gives
\[
 \mu_*\mathcal O_{W'}(mE)=\mathcal O_W.
\]
The projection formula therefore identifies the section spaces.
The modification is an isomorphism at each old generic divisorial
point, so the old local order calculations are unchanged. The final
assertion is the definition of $T'$.
\end{proof}

Appendix~\ref{p1logcbfalt:section} gives complementary projective
proofs of the section comparisons and further calculations with the
auxiliary subpair on higher models.

\appendix
\section{Supplementary projective calculations}
\label{p1logcbfalt:section}

The reduced-boundary construction of Section~\ref{cyc:section} admits
complementary projective proofs and calculations. We first give
divisorial proofs of the section comparisons and recover the
fixed-source descent by extracting valuations on the base. The fixed
logarithmic zero divisor then gives another discrepancy-sheaf test,
and a toroidal calculation proves logarithmic extension of the root
form. The final subsection tracks the auxiliary subpair when the
minimum-rule divisor is recomputed on a higher base model. This last
operation is distinct from the prescribed boundary modification of
Lemma~\ref{lem:normalized-models}: its discriminant is recomputed,
rather than defined as the strict boundary plus the reduced exceptional
divisor. Throughout these calculations, modifications of the original
source carry strict boundary plus reduced exceptional boundary.

\subsection{Divisorial sections and relative Iitaka systems}

The following elementary comparisons fix the complete systems and the
original models used in the normalization. All varieties in this
subsection are projective over $\C$, all rational divisor degrees are
sufficiently divisible, and the very general fibers are taken integral.
The logarithmic modification identity, also with a pulled-back rational
twist, is \Cref{setup:log-modification}.

\begin{lemma}[Divisorial section test]\label{p1logcbfalt:section-test}
Let \(V\) be normal and \(D\) an integral Weil divisor. A rational section
belongs to \(H^0(V,D)\) if and only if it satisfies the corresponding order
inequality at every prime divisor of \(V\). If \(\pi:V'\to V\) is finite
surjective with \(V'\) normal and \(D\) is Cartier, this test on \(V'\) needs
only primes lying over prime divisors of \(V\). If \(\pi\) is Galois, then
\[
H^0(V',\pi^*D)^{\operatorname{Gal}(V'/V)}=H^0(V,D).
\]
\end{lemma}

\begin{proof}
The first assertion is the description of a rank-one reflexive sheaf as
the intersection of its lattices at codimension-one points. For a finite
morphism between normal varieties, a prime divisor upstairs lies over a
prime divisor downstairs. The last assertion follows from
\((\pi_*\cO_{V'})^{\operatorname{Gal}(V'/V)}=\cO_V\) and the projection
formula.
\end{proof}

\begin{lemma}[Finite pullback]\label{p1logcbfalt:finite}
If \(\pi:V'\to V\) is a dominant generically finite morphism of normal
projective varieties and \(D\) is a rational Cartier divisor on \(V\),
then
\[
\kappa(V',\pi^*D)=\kappa(V,D).
\]
If \(\kappa(V,D)=0\), every nonzero space \(H^0(V,mD)\) has dimension
one.
\end{lemma}

\begin{proof}
Stein factorization reduces the first assertion to a finite morphism,
since a proper birational morphism onto a normal variety has direct image
\(\cO_V\). After replacing \(D\) by a multiple, use the natural inclusion
of graded section rings
\[
R(V,D)\ \subseteq\ R(V',\pi^*D).
\]
Each homogeneous section upstairs is integral over the ring downstairs.
Indeed its monic characteristic polynomial in the finite function-field
extension has coefficients which are rational sections of the appropriate
multiples of \(D\); they are regular in the required lattices at every
prime divisor, and hence are global sections by
Lemma~\ref{p1logcbfalt:section-test}. Thus the two nonzero graded rings have the
same transcendence degree, which gives equality of Iitaka dimensions.
If the ring downstairs has no positive-degree section, the same norm
argument excludes one upstairs. Finally, two independent sections in one
degree have a nonconstant ratio and define a positive-dimensional
rational image. This proves the last assertion.
\end{proof}

\begin{lemma}[Logarithmic pullback]\label{p1logcbfalt:log-pullback}
Let \(\rho:V\to Y\) be a dominant generically finite morphism of smooth projective
varieties. Let \(D_Y\) be a reduced snc divisor, and let \(B\) be a
boundary on \(V\) satisfying \(B\ge(\rho^*D_Y)_{\red}\). Assume that the
boundary supports on the chosen models are snc. Then pullback of rational
pluricanonical forms gives
\[
H^0(Y,m(K_Y+D_Y))\ \longrightarrow\
H^0(V,m(K_V+B))
\]
injectively, for divisible \(m\).
\end{lemma}

\begin{proof}
A logarithmic differential on a smooth snc pair pulls back to a
logarithmic differential along the reduced inverse image. In local
coordinates this follows from
\(d(\rho^*z)/\rho^*z\), whose divisorial poles are simple. Taking top
exterior powers and then tensor powers proves the order inequalities.
Equivalently, at a prime of ramification index \(e\) over a boundary
component, the Jacobian contributes \(e-1\), and the coefficient-one
boundary supplies the remaining one. At exceptional-image primes the
same logarithmic differential calculation applies. Injectivity follows
from dominance and separability.
\end{proof}

\begin{lemma}[Relative Iitaka fibration]\label{p1logcbfalt:iitaka}
Let \(p:V\to B\) be a dominant projective morphism, with \(V\) smooth
projective, and let \(D\) be a rational divisor whose restriction to the
geometric generic fiber has Iitaka dimension \(k\ge0\). After birational
modifications there is a relative Iitaka fibration
\[
V'\xrightarrow{q}U\longrightarrow B
\]
in which \(q\) has geometrically connected generic fiber, such that
\[
\dim U-\dim B=k,\qquad
\kappa(G,D'|_G)=0
\]
for a very general fiber \(G\) of \(q\). Here \(D'=\mu^*D\), where
\(\mu:V'\to V\) is the modification. When \(D=K_V+B_V\), one may
instead take the full log divisor from Lemma~\ref{setup:log-modification}.

If \(D'=K_{V'}+B_{V'}\) and the models are log smooth, then
\(D'|_G\sim_{\Q}K_G+B_G\).
\end{lemma}

\begin{proof}
Apply the Iitaka fibration theorem over the function field of \(B\)
\cite[Section~6, Remark~1]{IitakaOriginal},
take the relative algebraic closure of its section field in \(\C(V)\),
and spread out. Resolve the resulting rational maps and take their
Stein factorizations. The dimension of the relative section field is
\(k\); the restriction of the full divisor to the Iitaka fiber has
Iitaka dimension zero by the Iitaka fibration theorem. For the logarithmic
replacement, the proof of Lemma~\ref{setup:log-modification} gives the sheaf identity
\[
 \mu_*\cO_{V'}\bigl(m(K_{V'}+B_{V'})\bigr)
 =\cO_V\bigl(m(K_V+B_V)\bigr)
\]
in every sufficiently divisible degree. Pushing to the original base
identifies the relative section systems, even if there are no global
sections. Apply the same full-divisor Iitaka theorem to the replacement
divisor; its relative section field identifies its fibration with the
one already constructed and gives the asserted Iitaka dimension zero
on the fiber.
Generic smoothness and adjunction give the last assertion.
The dimensions at very general complex fibers agree with the geometric
generic values: for each divisible degree use cohomology and base
change, and remove the union of the exceptional closed sets over the
countably many degrees.
\end{proof}

\begin{lemma}[Easy addition]\label{p1logcbfalt:easy-addition}
Let \(p:V\to B\) be dominant and projective, with \(V\) smooth projective
and integral very general fiber \(G\), and let \(D\) be a rational
divisor on \(V\). Then
\[
\kappa(V,D)\le \dim B+\kappa(G,D|_G).
\]
If \(D\) has a nonzero global section, its restriction to a general
fiber is nonzero.
\end{lemma}

\begin{proof}
A nonzero section cannot vanish identically on all fibers over a dense
open subset. For each divisible \(m\) with sections, let \(\phi_m\)
be its rational map. The image of \((p,\phi_m)\) has dimension at most
\(\dim B+\kappa(G,D|_G)\), and dominates the image of \(\phi_m\).
Taking the maximum over \(m\) proves the inequality. If the fiber term
is \(-\infty\), the last assertion implies that the total term is
\(-\infty\) as well.
\end{proof}

\subsection{Restricted valuations and the fixed original source}

Let $p:H\to I$ be a dominant morphism of smooth projective varieties,
and fix the model $H$. A divisorial valuation $v$ of $\C(H)$ whose
restriction to $\C(I)$ is nontrivial restricts to a positive integral
multiple of a divisorial valuation of $\C(I)$. To see this, put
$n=\dim H$, $m=\dim I$, and denote the restricted valuation by $w$.
The value groups have rational rank one with finite-index inclusion.
The relative residue-transcendence-degree inequality gives
\[
 \operatorname{trdeg}_{\kappa(w)}\kappa(v)\le n-m.
\]
As $\operatorname{trdeg}_{\C}\kappa(v)=n-1$, this implies
$\operatorname{trdeg}_{\C}\kappa(w)\ge m-1$. The Abhyankar
inequality gives the reverse inequality. Thus $w$ is divisorial.
Finitely many such restrictions can be extracted on one smooth
projective birational model of $I$.

Only finitely many prime divisors of the fixed $H$ have image of
codimension at least two in $I$: they lie in the complement of a
locus where $p$ is smooth and are components of its divisorial part.
Extract their restrictions as above. Original horizontal primes
remain horizontal, while original vertical primes then have
divisorial centers on the prepared base. These centers remain
divisorial on higher base models. Hence on any subsequent commuting
birational diagram, a source prime whose base image still has
codimension at least two is exceptional over the fixed $H$.

For the rank-one section comparison, fix the original smooth pair
$(X,D_X)$, and let $X'\to X$ carry strict plus reduced exceptional
boundary $D'$. Let $x:X'\to W$ be the relative Iitaka fibration
on smooth projective models; its geometric generic fiber $(J,D_J)$
has logarithmic Kodaira dimension zero. Let $s$ be a nonzero
relative rational $p$-canonical form generating the degree-$p$
generic section space. For a source prime $E$, put
$d_E=\operatorname{coeff}_E D'$, and for a base prime $P$ define
\[
 t_P=\min_{E\mapsto P}
 \frac{p^{-1}\operatorname{ord}^{pK_{X'/W}}_E(s)+d_E}
 {\operatorname{ord}_E(x^*P)},\qquad
 T_W=\sum_Pt_PP.
\]
Here $T_W$ is the entire minimum-rule base term, before its division
into discriminant and moduli parts. In divisible degrees, division by
$s^{m/p}$ gives
\begin{equation}\label{p1logcbfalt:source-descent}
 H^0\bigl(X',m(K_{X'}+D')\bigr)
 \lhook\joinrel\longrightarrow
 H^0\bigl(W,m(K_W+T_W)\bigr).
\end{equation}
Indeed, on the geometric generic fiber the section is a constant
multiple of $s^{m/p}$; connectedness identifies that coefficient
with a rational $m$-canonical form $\sigma$ on $W$. At a component
$E$ dominating $P$, write $\sigma$ in a nonvanishing local frame of
$mK_W$, and put $a_E=\operatorname{ord}_E(x^*P)$. Its logarithmic
regularity upstairs is equivalent to
\[
 a_E\operatorname{ord}_P(\sigma)
 +(m/p)\operatorname{ord}^{pK_{X'/W}}_E(s)+md_E\ge0.
\]
The inequalities for all components above $P$ are exactly
$\operatorname{ord}_P(\sigma)+mt_P\ge0$. Horizontal primes impose
no further condition because $s$ is an allowed generic logarithmic
pluriform. The same argument works over the generic point of the
original base $Y$. If the relative Iitaka dimension is $k$, then
$\dim(W/Y)=k$, so preservation of section ratios makes
$K_W+T_W$ big over $Y$.

For the converse, extract on $W$ the restrictions of all source
divisors with exceptional image that survive as divisors on the
original $X$. Resolve the induced diagram and recompute the minimum
rule. Every divisor of the original $X$ is now tested either by the
generic horizontal condition or by a base prime. A section on the
right of \eqref{p1logcbfalt:source-descent}, multiplied by $s^{m/p}$,
therefore satisfies the logarithmic order inequality at every
original source prime. Any remaining untested pole lies on a
divisor exceptional over $X$. The codimension-one criterion on the
normal original source gives a section of $m(K_X+D_X)$; later
models preserve this criterion. Conversely the logarithmic
modification identity recovers its section on $X'$. The
identifications respect powers and multiplication. Under these
identifications, corresponding section ratios agree after pullback by
$x:X'\to W$. Their section fields and rational-map image dimensions
therefore agree.

\subsection{The fixed zero divisor and the discrepancy-sheaf test}

In the notation of Lemma~\ref{lem:reduced-cyclic-top}, the rational
divisor $Z=E/p$ is independent of the chosen nonzero logarithmic
pluriform and its degree. Indeed, if $\omega_m$ is a nonzero section
in another degree $m$, the degree-$mp$ sections $\omega_m^p$ and
$\omega^m$ are proportional, since that section space is
one-dimensional. Consequently
\[
 p\bigl(\operatorname{div}(\omega_m)+mD\bigr)=mE.
\]
The least positive integer trivializing $K_J+D-Z$ is also $p$.
Writing $\omega=b\eta^{\otimes p}$ in a rational canonical frame
gives $p(K_J+D-Z)=-\operatorname{div}(b)$. Conversely, a
trivialization of $r(K_J+D-Z)$ gives a rational $r$-canonical form
whose logarithmic zero divisor is $rZ\ge0$, hence a nonzero
logarithmic pluriform of degree $r$. Minimality gives $r\ge p$.

Here is a second way to check the generic discrepancy-sheaf rank
condition in Corollary~\ref{cor:reduced-integral-moduli}. Take a log
resolution $J'\to J$, let $D_{J'}$ be the strict boundary plus the
reduced exceptional divisor, and put
\[
 E_{J'}=\operatorname{div}(\omega|_{J'})+pD_{J'}\ge0.
\]
The crepant transform of the auxiliary generic subboundary is
$D_{J'}-E_{J'}/p$. Its coefficients are at most one. If $A^*$ is
the discrepancy b-divisor with coefficient $-1$ replaced by zero,
the rounded trace is effective, and a positive coefficient can
occur only on $\operatorname{Supp}E_{J'}$. Coefficientwise,
\[
 0\le \lceil A^*\rceil_{J'}
   \le\lceil E_{J'}/p\rceil\le E_{J'}.
\]
The first inequality uses the omission of discrepancy $-1$.
Logarithmic birational invariance gives $\kappa(E_{J'})=0$.
The associated space of rational functions therefore has dimension
at most one and contains the constants. This proves rank one on
every sufficiently high resolution, including when the auxiliary
subboundary has negative coefficients.

\subsection{A toroidal proof of logarithmic extension}

The extension part of Lemma~\ref{lem:reduced-cyclic-top} also has a
direct coordinate proof on a toroidal resolution of the SNC
cyclic-cover construction. First resolve the supports on $J$,
retaining strict plus reduced exceptional logarithmic boundary.
The minimal index and the fixed-zero assertions are unchanged.
In a local SNC chart write
\[
 \omega=\varepsilon\prod_i z_i^{a_i-pd_i}
                (dz_1\wedge\cdots\wedge dz_r)^{\otimes p},
 \qquad a_i\ge0,\quad d_i\in\{0,1\},
\]
where $\varepsilon$ is a unit. Include all coordinates; outside
$D+\operatorname{Supp}E$ put $d_i=a_i=0$. After a local unramified
change the unit has a root and contributes no divisorial order.
Let $v$ be a toroidal prime on the resolved root cover, and put
$n_i=v(z_i)\ge0$. The nonzero ray defining this prime has at least
one positive $n_i$. Rewriting with logarithmic differentials gives
\begin{equation}\label{p1logcbfalt:toroidal-root-order}
 \operatorname{ord}_v(\Omega)
   =\sum_i n_i\left(1-d_i+\frac{a_i}{p}\right)-1.
\end{equation}
The logarithmic volume form has order $-1$ along a toroidal prime,
and the change of torus lattice multiplies it by a nonzero constant.
Every coefficient in the sum is nonnegative. It vanishes exactly
when $d_i=1$ and $a_i=0$, namely along the coefficient-one pole
boundary of the auxiliary subpair.

If $v$ is outside the inverse image of this pole boundary, then
$n_i=0$ for each such zero coefficient. Some remaining $n_i$ is
positive, so the sum in \eqref{p1logcbfalt:toroidal-root-order} is
strictly positive. The order of the rational top form $\Omega$ is
integral, so that positive sum is at least one, giving regularity.
On the inverse pole boundary the nonnegative sum gives order at
least $-1$. Primes in the smooth unramified locus have the
regularity already checked on the finite normalization. Thus
$\Omega$ extends as a top form with logarithmic poles allowed only
along the inverse pole boundary. This provides the coordinate
alternative to the discrepancy-scaling proof.

\subsection{The normalized discriminant on higher models}

Let $x:X'\to W$ and $D'$ be the projective reduced-boundary
data of Lemma~\ref{lem:rankone-normalization}. Write $s$ for its
relative rational $p$-canonical generator, and set
\[
 t_P=\min_{E\mapsto P}
 \frac{p^{-1}\operatorname{ord}^{pK_{X'/W}}_E(s)+d_E}
 {\operatorname{ord}_E(x^*P)},\qquad
 T_W=\sum_Pt_PP,
\]
\[
 \Delta_0=-p^{-1}\operatorname{div}_{K_{X'/W}}(s),
 \qquad \Delta=\Delta_0+x^*T_W.
\]
In this subsection $T_W$ denotes the entire minimum-rule base
term, before its decomposition into discriminant and moduli parts;
the boundary denoted $T$ in Section~\ref{rankone:section} is the
discriminant part. We have
$K_{X'}+\Delta\sim_{\Q}x^*(K_W+T_W)$.
The auxiliary divisor is required to be sub-lc over the generic
point of $W$. Its coefficients at source divisors with image of
codimension at least two are not asserted to be at most one.

For a prime $P\subset W$, let
\[
 c_P=\operatorname{lct}_{\eta_P}(X',\Delta;x^*P),
 \qquad b_P=1-c_P,
 \qquad B_W=\sum_Pb_PP.
\]
Over the generic point of $P$ the minimum rule gives
$\Delta\le D'$, including the horizontal divisors. Log canonicity
therefore gives $c_P\ge0$. A component $E$ attaining the minimum
has coefficient exactly $d_E\in\{0,1\}$ in $\Delta$, and hence
\[
 c_P\le\frac{1-d_E}{\operatorname{ord}_E(x^*P)}\le1.
\]
It follows that $0\le B_W\le1$. If $P$ lies in a required inverse
boundary from a base $Y$, all components over $P$ have $d_E=1$,
so $c_P=0$ and $b_P=1$.

For the exact model comparison, let $\beta:W_1\to W_0$ be a
higher smooth base model and let $\pi:X_1\to X_0$ resolve the
induced diagram, with $x_i:X_i\to W_i$. Choose compatible
rational top forms for the canonical divisors. First form the
crepant transform
\[
 K_{X_1}+\Delta_{\mathrm{crepant}}
    =\pi^*(K_{X_0}+\Delta_{\mathrm{old}}).
\]
Using the new logarithmic boundary and the same relative
trivialization $s$, recompute the minimum-rule divisor $T_{W_1}$
and the normalized auxiliary divisor $\Delta_{\mathrm{new}}$.
Comparison of their definitions gives
\begin{equation}\label{p1logcbfalt:model-twist}
 \begin{split}
 Q&=K_{W_1}+T_{W_1}-\beta^*(K_{W_0}+T_{W_0}),\\
 \Delta_{\mathrm{new}}
   &=\Delta_{\mathrm{crepant}}+x_1^*Q.
 \end{split}
\end{equation}
Thus the twist includes the canonical discrepancy of the base.
Adding $x^*Q$ lowers the threshold at $P$ by
$\operatorname{coeff}_P Q$, adding $Q$ both to the discriminant
and to the total base term. It leaves the moduli part unchanged.
Birational pullback of the lc-trivial fibration also preserves its
moduli b-divisor; these compatibilities agree with the Hodge
extension description in \cite[Remarks~6.22 and~6.25]{BFMT}.

Resolve the finite supports on $W$ and resolve the induced source
diagram. At an unchanged codimension-one point, logarithmic forms
satisfying the old tests remain logarithmic on the source
resolution, while the strict old minimizing components persist.
Thus the minimum and discriminant trace there are unchanged.
New support is contained in the base-exceptional divisor. Its union
with the old strict support is SNC on the chosen resolution;
recomputing and applying the preceding coefficient argument yields
an effective SNC discriminant with the same coefficient-one inverse
boundary and the same moduli b-divisor.

Finally, if $m=ap$ and $s_m=c s^a$ on the generic fiber, with
$c\in\C(W)^*$, the minimum construction gives
\[
 T_W(s_m)=T_W(s)+m^{-1}\operatorname{div}(c).
\]
The two occurrences of $\operatorname{div}(c)$ cancel in the
definition of the normalized auxiliary subpair. It is unchanged.
The minimal-index cover is therefore compatible with all common
divisible degrees used for section descent.

\raggedbottom
\begin{thebibliography}{99}

\bibitem{AdiprasitoLiuTemkinSemistable}
Karim Adiprasito, Gaku Liu and Michael Temkin.
\emph{\href{https://arxiv.org/abs/1810.03131v2}{Semistable reduction in characteristic {0}}}.
Preprint, 2019.
\href{https://arxiv.org/abs/1810.03131v2}{arXiv:1810.03131v2}, 17 June 2019; 22-page full preprint.

\bibitem{AmbroBoundary}
Florin Ambro.
\emph{\href{https://arxiv.org/abs/math/0210271v1}{{Shokurov}'s boundary property}}.
Journal of Differential Geometry \textbf{67} (2004), no.~2, 229--255.
Theorem numbering is cited from arXiv:math/0210271v1.

\bibitem{AmbroModuli}
Florin Ambro.
\emph{\href{https://doi.org/10.1112/S0010437X04001071}{The moduli b-divisor of an lc-trivial fibration}}.
Compositio Mathematica \textbf{141} (2005), no.~2, 385--403.
\href{https://doi.org/10.1112/S0010437X04001071}{doi:\nolinkurl{10.1112/S0010437X04001071}}.

\bibitem{AndreMonodromy}
Yves Andr{\'e}.
\emph{\href{https://www.numdam.org/item/CM_1992__82_1_1_0/}{{Mumford--Tate} groups of mixed {Hodge} structures and the theorem of the fixed part}}.
Compositio Mathematica \textbf{82} (1992), no.~1, 1--24.

\bibitem{BakkerMatsushita}
Benjamin Bakker.
\emph{\href{https://benjamin-bakker.github.io/matsushita.pdf}{A short proof of a conjecture of {Matsushita}}}.
Advances in Mathematics \textbf{482} (2025), 110554.
\href{https://doi.org/10.1016/j.aim.2025.110554}{doi:\nolinkurl{10.1016/j.aim.2025.110554}}.
Part A.

\bibitem{BakkerBrunebarbeTsimerman}
Benjamin Bakker, Yohan Brunebarbe and Jacob Tsimerman.
\emph{\href{https://arxiv.org/abs/1811.12230v3}{o-minimal {GAGA} and a conjecture of {Griffiths}}}.
Inventiones Mathematicae \textbf{232} (2023), no.~1, 163--228.
\href{https://doi.org/10.1007/s00222-022-01166-1}{doi:\nolinkurl{10.1007/s00222-022-01166-1}}.
Theorem locators refer to arXiv:1811.12230v3, 23 October 2022.

\bibitem{BFMT}
Benjamin Bakker, Stefano Filipazzi, Mirko Mauri and Jacob Tsimerman.
\emph{\href{https://arxiv.org/abs/2508.19215v2}{{Baily--Borel} compactifications of period images and the b-semiampleness conjecture}}.
Preprint, 2025.
\href{https://arxiv.org/abs/2508.19215v2}{arXiv:2508.19215v2}, 18 December 2025.

\bibitem{BakkerTsimermanAxSchanuel}
Benjamin Bakker and Jacob Tsimerman.
\emph{\href{https://arxiv.org/abs/1712.05088v1}{The {Ax--Schanuel} conjecture for variations of {Hodge} structures}}.
Inventiones Mathematicae \textbf{217} (2019), no.~1, 77--94.
\href{https://doi.org/10.1007/s00222-019-00863-8}{doi:\nolinkurl{10.1007/s00222-019-00863-8}}.
The cited locators refer to \href{https://arxiv.org/abs/1712.05088v1}{arXiv:1712.05088v1}.

\bibitem{BierstoneMilmanFunctorial}
Edward Bierstone and Pierre D. Milman.
\emph{\href{https://ems.press/journals/prims/articles/2154}{Functoriality in resolution of singularities}}.
Publications of the Research Institute for Mathematical Sciences \textbf{44} (2008), no.~2, 609--639.
\href{https://doi.org/10.2977/PRIMS/1210167338}{doi:\nolinkurl{10.2977/PRIMS/1210167338}}.

\bibitem{BirkarIitakaSix}
Caucher Birkar.
\emph{\href{https://arxiv.org/abs/0806.4413}{The {Iitaka} conjecture {$C_{n,m}$} in dimension six}}.
Compositio Mathematica \textbf{145} (2009), no.~6, 1442--1446.
\href{https://doi.org/10.1112/S0010437X09004187}{doi:\nolinkurl{10.1112/S0010437X09004187}}.
arXiv:0806.4413, first posted 27 June 2008.

\bibitem{BCHM}
Caucher Birkar, Paolo Cascini, Christopher D. Hacon and James McKernan.
\emph{\href{https://arxiv.org/abs/math/0610203v2}{Existence of minimal models for varieties of log general type}}.
Journal of the American Mathematical Society \textbf{23} (2010), no.~2, 405--468.
\href{https://doi.org/10.1090/S0894-0347-09-00649-3}{doi:\nolinkurl{10.1090/S0894-0347-09-00649-3}}.
The cited locators refer to the 81-page \href{https://arxiv.org/abs/math/0610203v2}{arXiv:math/0610203v2}.

\bibitem{BoucksomVolume}
S{\'e}bastien Boucksom.
\emph{\href{https://arxiv.org/abs/math/0201031v1}{On the volume of a line bundle}}.
International Journal of Mathematics \textbf{13} (2002), no.~10, 1043--1063.
\href{https://doi.org/10.1142/S0129167X02001575}{doi:\nolinkurl{10.1142/S0129167X02001575}}.
Theorem locators refer to arXiv:math/0201031v1, 5 January 2002.

\bibitem{BEGZ}
S{\'e}bastien Boucksom, Philippe Eyssidieux, Vincent Guedj and Ahmed Zeriahi.
\emph{\href{https://www.math.univ-toulouse.fr/~guedj/fichierspdf/BEGZ10.pdf}{{Monge--Amp{\`e}re} equations in big cohomology classes}}.
Acta Mathematica \textbf{205} (2010), 199--262.
\href{https://doi.org/10.1007/s11511-010-0054-7}{doi:\nolinkurl{10.1007/s11511-010-0054-7}}.

\bibitem{BrunebarbeCadorel}
Yohan Brunebarbe and Beno{\^i}t Cadorel.
\emph{\href{https://arxiv.org/abs/1707.01327v1}{Hyperbolicity of varieties supporting a variation of {Hodge} structure}}.
International Mathematics Research Notices \textbf{2020} (2020), no.~6, 1601--1609.
\href{https://doi.org/10.1093/imrn/rny054}{doi:\nolinkurl{10.1093/imrn/rny054}}.
The cited locators refer to \href{https://arxiv.org/abs/1707.01327v1}{arXiv:1707.01327v1}.

\bibitem{Campana2004}
Fr{\'e}d{\'e}ric Campana.
\emph{\href{https://www.numdam.org/item/10.5802/aif.2027/}{Orbifolds, special varieties and classification theory}}.
Annales de l'Institut Fourier \textbf{54} (2004), no.~3, 499--630.
\href{https://doi.org/10.5802/aif.2027}{doi:\nolinkurl{10.5802/aif.2027}}.

\bibitem{CampanaOrbifolds}
Fr{\'e}d{\'e}ric Campana.
\emph{\href{https://arxiv.org/abs/0705.0737v8}{Orbifoldes g{\'e}om{\'e}triques sp{\'e}ciales et classification bim{\'e}romorphe des vari{\'e}t{\'e}s k{\"a}hl{\'e}riennes compactes}}.
Journal of the Institute of Mathematics of Jussieu \textbf{10} (2011), no.~4, 809--934.
\href{https://doi.org/10.1017/S1474748010000101}{doi:\nolinkurl{10.1017/S1474748010000101}}.
The cited locators refer to \href{https://arxiv.org/abs/0705.0737v8}{arXiv:0705.0737v8}.

\bibitem{CaoSurfaces}
Junyan Cao.
\emph{\href{https://content.algebraicgeometry.nl/2018-6/2018-6-020.pdf}{{Kodaira} dimension of algebraic fiber spaces over surfaces}}.
Algebraic Geometry \textbf{5} (2018), no.~6, 728--741.
\href{https://doi.org/10.14231/AG-2018-020}{doi:\nolinkurl{10.14231/AG-2018-020}}.

\bibitem{CaoPaunAbelian}
Junyan Cao and Mihai P{\u a}un.
\emph{\href{https://arxiv.org/abs/1504.01095v3}{{Kodaira} dimension of algebraic fiber spaces over abelian varieties}}.
Inventiones Mathematicae \textbf{207} (2017), no.~1, 345--387.
\href{https://doi.org/10.1007/s00222-016-0672-6}{doi:\nolinkurl{10.1007/s00222-016-0672-6}}.
Theorem locators refer to arXiv:1504.01095v3, submitted 3 August 2016.

\bibitem{CKS}
Eduardo Cattani, Aroldo Kaplan and Wilfried Schmid.
\emph{\href{https://annals.math.princeton.edu/1986/123-3/p02}{Degeneration of {Hodge} structures}}.
Annals of Mathematics (2) \textbf{123} (1986), no.~3, 457--535.
\href{https://doi.org/10.2307/1971333}{doi:\nolinkurl{10.2307/1971333}}.

\bibitem{ChenLimits}
Qianyu Chen.
\emph{\href{https://qymath.github.io/Steenbrink.pdf}{Limits of {Hodge} structures via holonomic {$D$}-modules}}.
Preprint, 2021.
arXiv:2103.03983; the cited Lemma 8.1 is in the \href{https://qymath.github.io/Steenbrink.pdf}{61-page author manuscript}.

\bibitem{DeligneHodgeII}
Pierre Deligne.
\emph{Th\'eorie de Hodge. II}.
Publications Math\'ematiques de l'IH\'ES \textbf{40} (1971), 5--57.
\href{https://doi.org/10.1007/BF02684692}{doi:\nolinkurl{10.1007/BF02684692}}.

\bibitem{DemaillyRegularization}
Jean-Pierre Demailly.
\emph{\href{https://www-fourier.univ-grenoble-alpes.fr/~demailly/manuscripts/regularization.pdf}{Regularization of closed positive currents and intersection theory}}.
Journal of Algebraic Geometry \textbf{1} (1992), no.~3, 361--409.
The cited locator refers to the \href{https://www-fourier.univ-grenoble-alpes.fr/~demailly/manuscripts/regularization.pdf}{author manuscript}.

\bibitem{DemaillyPaun2004}
Jean-Pierre Demailly and Mihai P{\u{a}}un.
\emph{\href{https://annals.math.princeton.edu/2004/159-3/p05}{Numerical characterization of the {K{\"a}hler} cone of a compact {K{\"a}hler} manifold}}.
Annals of Mathematics (2) \textbf{159} (2004), no.~3, 1247--1274.
\href{https://doi.org/10.4007/annals.2004.159.1247}{doi:\nolinkurl{10.4007/annals.2004.159.1247}}.

\bibitem{Douady}
Adrien Douady.
\emph{\href{https://www.numdam.org/item/AIF_1966__16_1_1_0/}{Le probl{\`e}me des modules pour les sous-espaces analytiques compacts d'un espace analytique donn{\'e}}}.
Annales de l'Institut Fourier \textbf{16} (1966), no.~1, 1--95.
\href{https://doi.org/10.5802/aif.226}{doi:\nolinkurl{10.5802/aif.226}}.

\bibitem{EnokizonoHashizumeSemistable}
Makoto Enokizono and Kenta Hashizume.
\emph{\href{https://arxiv.org/abs/2404.05124v3}{Semistable reduction for complex analytic spaces}}.
Transactions of the American Mathematical Society \textbf{378} (2025), 7667--7688.
\href{https://doi.org/10.1090/tran/9365}{doi:\nolinkurl{10.1090/tran/9365}}.
The cited locators refer to \href{https://arxiv.org/abs/2404.05124v3}{arXiv:2404.05124v3}, 12 October 2024.

\bibitem{FujikiDouady}
Akira Fujiki.
\emph{\href{https://ems.press/content/serial-article-files/41937}{Closedness of the {Douady} spaces of compact {K{\"a}hler} spaces}}.
Publications of the Research Institute for Mathematical Sciences \textbf{14} (1978), no.~1, 1--52.
\href{https://doi.org/10.2977/PRIMS/1195189279}{doi:\nolinkurl{10.2977/PRIMS/1195189279}}.

\bibitem{FujinoWeakPositivity}
Osamu Fujino.
\emph{\href{https://www.math.kyoto-u.ac.jp/~fujino/weak-posi11.pdf}{Notes on the weak positivity theorems}}.
In \emph{Algebraic Varieties and Automorphism Groups}.
Advanced Studies in Pure Mathematics 75.
Mathematical Society of Japan, 2017, pp.~73--118.
\href{https://doi.org/10.2969/aspm/07510073}{doi:\nolinkurl{10.2969/aspm/07510073}}.
The cited locators refer to \href{https://www.math.kyoto-u.ac.jp/~fujino/weak-posi11.pdf}{author version 0.54, 30 June 2015}.

\bibitem{FujinoAnalyticAdjunction}
Osamu Fujino.
\emph{\href{https://www.math.kyoto-u.ac.jp/~fujino/on-inversion-of-adjunction-pja-final.pdf}{Log canonical inversion of adjunction}}.
Proceedings of the Japan Academy, Series A, Mathematical Sciences \textbf{100} (2024), no.~2, 7--11.
\href{https://doi.org/10.3792/PJAA.100.002}{doi:\nolinkurl{10.3792/PJAA.100.002}}.

\bibitem{FujinoFujisawa}
Osamu Fujino and Taro Fujisawa.
\emph{\href{https://ems.press/content/serial-article-files/41240}{Variations of mixed {Hodge} structure and semipositivity theorems}}.
Publications of the Research Institute for Mathematical Sciences \textbf{50} (2014), no.~4, 589--661.
\href{https://doi.org/10.4171/PRIMS/145}{doi:\nolinkurl{10.4171/PRIMS/145}}.

\bibitem{FujinoFujisawaVMHS}
Osamu Fujino and Taro Fujisawa.
\emph{\href{https://www.math.kyoto-u.ac.jp/~fujino/vmhs-applications8.pdf}{Variation of mixed {Hodge} structure and its applications}}.
International Journal of Mathematics \textbf{36} (2025), no.~9, 2550020.
Published article: 43 pp.; theorem locators refer to the 36-page \href{https://www.math.kyoto-u.ac.jp/~fujino/vmhs-applications8.pdf}{author version 0.55, 11 March 2025}.

\bibitem{FujinoGongyo}
Osamu Fujino and Yoshinori Gongyo.
\emph{\href{https://www.numdam.org/item/10.5802/aif.2894.pdf}{On the moduli b-divisors of lc-trivial fibrations}}.
Annales de l'Institut Fourier \textbf{64} (2014), no.~4, 1721--1735.
\href{https://doi.org/10.5802/aif.2894}{doi:\nolinkurl{10.5802/aif.2894}}.

\bibitem{FujinoMori}
Osamu Fujino and Shigefumi Mori.
\emph{\href{https://www.math.kyoto-u.ac.jp/~fujino/Canonical.pdf}{A canonical bundle formula}}.
Journal of Differential Geometry \textbf{56} (2000), no.~1, 167--188.

\bibitem{GrauertDirectImages}
Hans Grauert.
\emph{\href{https://www.numdam.org/item/PMIHES_1960__5__5_0/}{Ein Theorem der analytischen Garbentheorie und die Modulr{\"a}ume komplexer Strukturen}}.
Publications Math{\'e}matiques de l'IH{\'E}S \textbf{5} (1960), 5--64.
\href{https://doi.org/10.1007/BF02684746}{doi:\nolinkurl{10.1007/BF02684746}}.
With corrigendum, Publ. Math. IH{\'E}S 16 (1963), 35--36.

\bibitem{GrauertDirectImagesCorr}
Hans Grauert.
\emph{\href{https://www.numdam.org/articles/10.1007/BF02684291/}{Berichtigung: Ein Theorem der analytischen Garbentheorie und die Modulr{\"a}ume komplexer Strukturen}}.
Publications Math{\'e}matiques de l'IH{\'E}S \textbf{16} (1963), 35--36.
\href{https://doi.org/10.1007/BF02684291}{doi:\nolinkurl{10.1007/BF02684291}}.

\bibitem{GriffithsPeriodsIII}
Phillip A. Griffiths.
\emph{\href{https://www.numdam.org/item/PMIHES_1970__38__125_0/}{Periods of integrals on algebraic manifolds, {III}: Some global differential-geometric properties of the period mapping}}.
Publications Math{\'e}matiques de l'IH{\'E}S \textbf{38} (1970), 125--180.
\href{https://doi.org/10.1007/BF02684654}{doi:\nolinkurl{10.1007/BF02684654}}.

\bibitem{GuedjZeriahiEnergy}
Vincent Guedj and Ahmed Zeriahi.
\emph{\href{https://arxiv.org/abs/math/0612630v2}{The weighted {Monge--Amp{\`e}re} energy of quasiplurisubharmonic functions}}.
Journal of Functional Analysis \textbf{250} (2007), no.~2, 442--482.
\href{https://doi.org/10.1016/j.jfa.2007.04.018}{doi:\nolinkurl{10.1016/j.jfa.2007.04.018}}.
The cited locators refer to \href{https://arxiv.org/abs/math/0612630v2}{arXiv:math/0612630v2}.

\bibitem{HaconPopaSchnellAbelian}
Christopher D. Hacon, Mihnea Popa and Christian Schnell.
\emph{\href{https://people.math.harvard.edu/~mpopa/papers/HPS.pdf}{Algebraic fiber spaces over abelian varieties: around a recent theorem by {Cao} and {P{\u a}un}}}.
In \emph{Local and Global Methods in Algebraic Geometry}.
Contemporary Mathematics 712.
American Mathematical Society, Providence, RI, 2018, pp.~143--195.
\href{https://doi.org/10.1090/conm/712/14346}{doi:\nolinkurl{10.1090/conm/712/14346}}.
Theorem locators refer to the author manuscript dated 9 May 2017.

\bibitem{Hashizume}
Kenta Hashizume.
\emph{\href{https://arxiv.org/abs/1902.10923v2}{Log {Iitaka} conjecture for abundant log canonical fibrations}}.
Proceedings of the Japan Academy, Series A, Mathematical Sciences \textbf{96} (2020), no.~10, 87--92.
\href{https://doi.org/10.3792/pjaa.96.017}{doi:\nolinkurl{10.3792/pjaa.96.017}}.
Theorem numbering is cited from arXiv:1902.10923v2.

\bibitem{HironakaFlattening}
Heisuke Hironaka.
\emph{\href{https://doi.org/10.2307/2373721}{Flattening theorem in complex-analytic geometry}}.
American Journal of Mathematics \textbf{97} (1975), no.~2, 503--547.
\href{https://doi.org/10.2307/2373721}{doi:\nolinkurl{10.2307/2373721}}.

\bibitem{IitakaOriginal}
Shigeru Iitaka.
\emph{\href{https://www.jstage.jst.go.jp/article/jmath1948/23/2/23_2_356/_article}{On {$D$}-dimensions of algebraic varieties}}.
Journal of the Mathematical Society of Japan \textbf{23} (1971), no.~2, 356--373.
\href{https://doi.org/10.2969/jmsj/02320356}{doi:\nolinkurl{10.2969/jmsj/02320356}}.

\bibitem{IitakaLog}
Shigeru Iitaka.
\emph{\href{https://doi.org/10.1017/CBO9780511569197.014}{On logarithmic {Kodaira} dimension of algebraic varieties}}.
In \emph{Complex Analysis and Algebraic Geometry: A Collection of Papers Dedicated to {K. Kodaira}}.
Walter L. Baily, Jr. and Tetsuji Shioda (eds.).
Iwanami Shoten and Cambridge University Press, 1977.
\href{https://doi.org/10.1017/CBO9780511569197.014}{doi:\nolinkurl{10.1017/CBO9780511569197.014}}.

\bibitem{IitakaClassification}
Shigeru Iitaka.
\emph{\href{https://www.jstage.jst.go.jp/article/pjab1945/53/3/53_3_103/_pdf}{Classification of Algebraic Varieties}}.
Proceedings of the Japan Academy \textbf{53} (1977), no.~3, 103--105.
\href{https://doi.org/10.3792/pja/1195517881}{doi:\nolinkurl{10.3792/pja/1195517881}}.

\bibitem{KawamataCurves}
Yujiro Kawamata.
\emph{\href{https://link.springer.com/article/10.1007/BF01404756}{{Kodaira} dimension of algebraic fiber spaces over curves}}.
Inventiones Mathematicae \textbf{66} (1982), no.~1, 57--71.
\href{https://doi.org/10.1007/BF01404756}{doi:\nolinkurl{10.1007/BF01404756}}.

\bibitem{KawamataKodaira}
Yujiro Kawamata.
\emph{\href{https://doi.org/10.1515/crll.1985.363.1}{Minimal models and the {Kodaira} dimension of algebraic fiber spaces}}.
Journal f{\"u}r die reine und angewandte Mathematik \textbf{363} (1985), 1--46.
\href{https://doi.org/10.1515/crll.1985.363.1}{doi:\nolinkurl{10.1515/crll.1985.363.1}}.

\bibitem{Kollar87}
J{\'a}nos Koll{\'a}r.
\emph{\href{https://doi.org/10.2969/aspm/01010361}{Subadditivity of the {Kodaira} dimension: fibers of general type}}.
In \emph{Algebraic Geometry, Sendai, 1985}.
Advanced Studies in Pure Mathematics 10.
Kinokuniya and North-Holland, Tokyo and Amsterdam, 1987, pp.~361--398.
\href{https://doi.org/10.2969/aspm/01010361}{doi:\nolinkurl{10.2969/aspm/01010361}}.

\bibitem{KovacsPatakfalvi}
S{\'a}ndor J. Kov{\'a}cs and Zsolt Patakfalvi.
\emph{\href{https://sites.math.washington.edu/~kovacs/current/papers/Kovacs_Patakfalvi__Projectivity.pdf}{Projectivity of the moduli space of stable log-varieties and subadditivity of log-{Kodaira} dimension}}.
Journal of the American Mathematical Society \textbf{30} (2017), no.~4, 959--1021.
\href{https://doi.org/10.1090/jams/871}{doi:\nolinkurl{10.1090/jams/871}}.
The cited locators refer to the \href{https://sites.math.washington.edu/~kovacs/current/papers/Kovacs_Patakfalvi__Projectivity.pdf}{62-page author manuscript}.

\bibitem{MaeharaIitaka}
Kazuhisa Maehara.
\emph{\href{https://arxiv.org/abs/0910.0604v3}{{Iitaka--Viehweg} conjectures {$C$} and {$C^{++}$}}}.
arXiv preprint, 2010.
Version 3, 16 May 2010; first posted 4 October 2009.

\bibitem{MilneAlgebraicGroups}
James S. Milne.
\emph{\href{https://www.jmilne.org/math/Books/iAG2022.pdf}{Algebraic Groups: The Theory of Group Schemes of Finite Type over a Field}}.
Cambridge Studies in Advanced Mathematics 170.
Cambridge University Press, 2017.
\href{https://doi.org/10.1017/9781316711736}{doi:\nolinkurl{10.1017/9781316711736}}.
Cited from the corrected 2022 reprint.

\bibitem{NamikawaProjectivity}
Yoshinori Namikawa.
\emph{\href{https://arxiv.org/abs/math/0101019v6}{Projectivity criterion of {Moishezon} spaces and density of projective symplectic varieties}}.
International Journal of Mathematics \textbf{13} (2002), no.~2, 125--135.
Theorem locators refer to arXiv:math/0101019v6, 12 November 2001.

\bibitem{IitakaVariationCompanion}
{\raggedright
{OpenAI}.
\emph{Logarithmic Kodaira dimension and whole-fiber variation}.
OpenAI Math Release preprint
\href{https://github.com/openai/math/blob/main/preprints/Logarithmic-Kodaira-dimension-and-whole-fiber-variation-September-26-2026/paper.pdf}{OAI:Logarithmic-Kodaira-dimension-and-whole-fiber-variation-September-26-2026}, 2026.\par}

\bibitem{IitakaProjectiveCompanion}
{\raggedright
{OpenAI}.
\emph{Projective Hodge lines and ordinary Iitaka subadditivity}.
OpenAI Math Release preprint
\href{https://github.com/openai/math/blob/main/preprints/Projective-Hodge-lines-and-ordinary-Iitaka-subadditivity-September-27-2026/paper.pdf}{OAI:Projective-Hodge-lines-and-ordinary-Iitaka-subadditivity-September-27-2026}, 2026.\par}

\bibitem{IitakaAdditivityCompanion}
{\raggedright
{OpenAI}.
\emph{The reverse logarithmic Kodaira inequality and additivity}.
OpenAI Math Release preprint
\href{https://github.com/openai/math/blob/main/preprints/The-reverse-logarithmic-Kodaira-inequality-and-additivity-September-26-2026/paper.pdf}{OAI:The-reverse-logarithmic-Kodaira-inequality-and-additivity-September-26-2026}, 2026.\par}

\bibitem{OpenAILogAbundance2026}
{\raggedright
{OpenAI}.
\emph{Log abundance in characteristic zero}.
OpenAI Math Release preprint
\href{https://github.com/openai/math/blob/main/preprints/Log-abundance-in-characteristic-zero-September-24-2026/paper.pdf}{OAI:Log-abundance-in-characteristic-zero-September-24-2026}, 2026.\par}

\bibitem{PopaSchnellKahler}
Giuseppe Pareschi, Mihnea Popa and Christian Schnell.
\emph{\href{https://msp.org/gt/2017/21-4/gt-v21-n4-p12-s.pdf}{{Hodge} modules on complex tori and generic vanishing for compact {K{\"a}hler} manifolds}}.
Geometry \& Topology \textbf{21} (2017), 2419--2460.
\href{https://doi.org/10.2140/gt.2017.21.2419}{doi:\nolinkurl{10.2140/gt.2017.21.2419}}.

\bibitem{PatakfalviFibered}
Zsolt Patakfalvi.
\emph{\href{https://arxiv.org/abs/1208.1787v4}{Fibered stable varieties}}.
Transactions of the American Mathematical Society \textbf{368} (2016), no.~3, 1837--1869.
\href{https://doi.org/10.1090/tran/6386}{doi:\nolinkurl{10.1090/tran/6386}}.

\bibitem{Schmid}
Wilfried Schmid.
\emph{\href{https://doi.org/10.1007/BF01389674}{Variation of {Hodge} structure: The singularities of the period mapping}}.
Inventiones Mathematicae \textbf{22} (1973), no.~3--4, 211--319.
\href{https://doi.org/10.1007/BF01389674}{doi:\nolinkurl{10.1007/BF01389674}}.

\bibitem{SommesePeriodMapping}
Andrew John Sommese.
\emph{\href{https://www.numdam.org/item/ASNSP_1978_4_5_4_683_0/}{On the rationality of the period mapping}}.
Annali della Scuola Normale Superiore di Pisa, Classe di Scienze (4)
\textbf{5} (1978), no.~4, 683--717.

\bibitem{TsujiDirectImages}
Hajime Tsuji.
\emph{\href{https://arxiv.org/abs/1012.0884v1}{Global generation of the direct images of relative pluricanonical systems}}.
arXiv preprint, 2010.
Version 1, 4 December 2010.

\bibitem{UenoClassification}
Kenji Ueno.
\emph{\href{https://doi.org/10.1007/BFb0070570}{Classification Theory of Algebraic Varieties and Compact Complex Spaces}}.
Lecture Notes in Mathematics 439.
Springer, Berlin, 1975.
\href{https://doi.org/10.1007/BFb0070570}{doi:\nolinkurl{10.1007/BFb0070570}}.

\bibitem{ViehwegAdditivity}
Eckart Viehweg.
\emph{\href{https://esaga.uni-due.de/eckart.viehweg/ecpubl/}{Weak positivity and the additivity of the {Kodaira} dimension for certain fibre spaces}}.
In \emph{Algebraic Varieties and Analytic Varieties (Tokyo, 1981)}.
Advanced Studies in Pure Mathematics 1.
Kinokuniya and North-Holland, Tokyo and Amsterdam, 1983, pp.~329--353.
\href{https://doi.org/10.2969/aspm/00110329}{doi:\nolinkurl{10.2969/aspm/00110329}}.

\bibitem{VilladsenHodgeKahler}
Mads Bach Villadsen.
\emph{\href{https://arxiv.org/abs/2401.09544v1}{{Hodge} modules and {K{\"a}hler} morphisms}}.
Preprint, 2024.
\href{https://arxiv.org/abs/2401.09544v1}{arXiv:2401.09544v1}, 17 January 2024.

\bibitem{WangKahlerIitaka}
Juanyong Wang.
\emph{\href{https://afst.centre-mersenne.org/item/10.5802/afst.1690.pdf}{On the {Iitaka} Conjecture {$C_{n,m}$} for {K{\"a}hler} Fibre Spaces}}.
Annales de la Facult{\'e} des Sciences de Toulouse. Math{\'e}matiques (6) \textbf{30} (2021), no.~4, 813--897.
\href{https://doi.org/10.5802/afst.1690}{doi:\nolinkurl{10.5802/afst.1690}}.

\end{thebibliography}

\end{document}
